% !TEX program = xelatex
% !LW recipe = XeLaTeX
\documentclass[11pt, a4paper]{extarticle}
\usepackage{amsmath, amsthm, xspace, enumitem, comment, etoolbox, thmtools, needspace}
\usepackage[table]{xcolor}
\usepackage{tikz}
\usetikzlibrary{positioning, arrows.meta, shapes.geometric}
\usepackage{booktabs}
\usepackage{amssymb}

% Цвета из вашего шаблона
\definecolor{darkred}{RGB}{150,0,0}
\definecolor{darkblue}{RGB}{0,0,130}
\definecolor{darkgreen}{RGB}{0,100,0}
\definecolor{darkviolet}{RGB}{80,0,120}
\definecolor{darkteal}{RGB}{0,118,118}
 
% --- ПРЕАМБУЛА (XeLaTeX) ---
\usepackage{fontspec} % Главный пакет для работы со шрифтами в XeLaTeX
\usepackage{unicode-math} % Современный пакет для математики
\setmainfont{XCharter}
\setmathfont{XCharter Math}
%\setsansfont{XCharter}
\AtBeginDocument{%
   \renewcommand{\eqdef}{\ifmmode{\,\color{darkred}\mathop{\overset{\text{\tiny def}}{\rightleftharpoons}}\,}\fi}%
}
\usepackage[russian]{babel}

% Геометрия страницы: одна непрерывная страница с сохранением ширины и отступов
\usepackage[
paperwidth=30cm,   % Ширина
paperheight=14400pt,  % Очень большая высота для непрерывной страницы
left=12cm, 
right=5mm, 
top=10mm, 
bottom=10mm, 
%showframe
]{geometry}

\tolerance=9999 % Разрешить умеренно плохие пробелы

% Настройка теорем и определений
\declaretheoremstyle[
    headfont=\bfseries\color{darkblue},
    bodyfont=\normalfont,
    mdframed={
        linewidth=1pt,
        linecolor=darkblue,
        backgroundcolor=blue!5,
    },
    preheadhook={\needspace{8\baselineskip}},
    postheadhook={\leavevmode},
]{thmblue}

\declaretheoremstyle[
    headfont=\bfseries\color{darkteal},
    bodyfont=\normalfont,
    mdframed={
        linewidth=1pt,
        linecolor=darkteal,
        backgroundcolor=darkteal!8,
    },
    preheadhook={\needspace{8\baselineskip}},
    postheadhook={\leavevmode},
]{definitionstyle}

\declaretheoremstyle[
    headfont=\bfseries\color{darkgreen},
    bodyfont=\small\itshape,
    spaceabove=6pt, 
    spacebelow=6pt,
]{remarkstyle}

\declaretheoremstyle[
    headfont=\bfseries\color{darkred},
    bodyfont=\normalfont,
    mdframed={
        linewidth=1pt,
        linecolor=darkred,
        backgroundcolor=red!5,
    },
    preheadhook={\needspace{8\baselineskip}},
    postheadhook={\leavevmode},
]{exerstyle}

\declaretheorem[style=definitionstyle, name=Определение, within=section]{definition}
\declaretheorem[style=thmblue, name=Теорема, within=section]{theorem}
\declaretheorem[style=thmblue, name=Следствие, within=section]{corollary}
\declaretheorem[style=remarkstyle, name=Комментарий, numbered=no]{remark}
\declaretheorem[style=exerstyle, name=Блок задач]{exercise}

% Стиль для примеров
\declaretheoremstyle[
    headfont=\bfseries\color{darkgreen},
    bodyfont=\normalfont,
    mdframed={
        linewidth=1pt,
        linecolor=darkgreen,
        backgroundcolor=green!5,
    },
    preheadhook={\needspace{8\baselineskip}},
    postheadhook={\leavevmode},
]{examplestyle}

\declaretheorem[style=examplestyle, name=Пример]{example}

% Стиль для аксиом (отдельный цвет: фиолетовый)
\declaretheoremstyle[
    headfont=\bfseries\color{darkviolet},
    bodyfont=\normalfont,
    mdframed={
        linewidth=1pt,
        linecolor=darkviolet,
        backgroundcolor=violet!6,
    },
    preheadhook={\needspace{8\baselineskip}},
    postheadhook={\leavevmode},
]{axiomstyle}
\declaretheorem[style=axiomstyle, name=Аксиома, within=section]{axiom}

\newcommand{\eqdef}{\,\mathop{\color{darkred}\overset{\text{\tiny def}}{\rightleftharpoons}}\,}

% Команда для аннотации
\newcommand{\annotation}[1]{%
    \begin{abstract}
        #1
    \end{abstract}
}

% СОКРАЩЕНИЯ И НОВЫЕ ОПЕРАТОРЫ И СИМВОЛЫ
\newcommand{\al}{\alpha}
\newcommand{\be}{\beta}
\newcommand{\ga}{\gamma}
\newcommand{\de}{\delta}
\newcommand{\ep}{\varepsilon}
\newcommand{\la}{\lambda}
\newcommand{\La}{\Lambda}
\newcommand{\ka}{\varkappa}
\newcommand{\si}{\sigma}
\newcommand{\Si}{\Sigma}
\newcommand{\ph}{\varphi}
\newcommand{\om}{\omega}
\newcommand{\Om}{\Omega}
\newcommand{\De}{\Delta}
\newcommand{\A}{{\mathbb A}}
\newcommand{\AZ}{{\mathbb AZ}}
\newcommand{\e}{{\mathsf e}}
\newcommand{\F}{{\mathbb F}}
\newcommand{\T}{{\mathbb T}}
\newcommand{\TP}{{\mathbb TP}}
\newcommand{\W}{{\mathfrak W}}
\newcommand{\B}{{\mathbb B}}
\newcommand{\Rec}{{\mathfrak R}}
\newcommand{\Pcal}{{\mathcal P}}
\newcommand{\Dcal}{{\mathcal D}}
\newcommand{\Gcal}{{\mathcal G}}
\newcommand{\Fcal}{{\mathcal F}}
\newcommand{\Scal}{{\mathcal S}}
\newcommand{\Rcal}{{\mathcal R}}
\newcommand{\Qcal}{{\mathcal Q}}
\newcommand{\Tcal}{{\mathcal T}}
\newcommand{\Lcal}{{\mathcal L}}
\newcommand{\Ncal}{{\mathcal N}}
\newcommand{\Ucal}{{\mathcal U}}
\newcommand{\Mcal}{{\mathcal M}}
\newcommand{\Acal}{{\mathcal A}}
\newcommand{\Bcal}{{\mathcal B}}
\newcommand{\Ccal}{{\mathcal C}}
\newcommand{\Kcal}{{\mathcal K}}
\newcommand{\E}{{\mathsf E}}
\newcommand{\U}{{\mathfrak U}}
\newcommand{\M}{{\mathcal M}}
\newcommand{\Z}{{\mathbb Z}}
\newcommand{\R}{{\mathbb R}}
\newcommand{\RP}{\mathbb{RP}}
\newcommand{\QP}{\mathbb{QP}}
\newcommand{\CP}{\mathbb{CP}}
\newcommand{\HR}{{\mathbb H\hspace{-0.38em}\mathbb R}}
\newcommand{\Los}{{{\fontfamily{lmr}\selectfont\L}o\'s}}
\newcommand{\Lang}{{\mathcal L}}
\newcommand{\Sur}{{\mathsf{No}}}
\newcommand{\C}{{\mathbb C}}
\let\oldH\H
\renewcommand{\H}{\ifmmode\mathbb{H}\else\oldH\fi}
\newcommand{\Harm}{{\mathcal H}}
\newcommand{\Hcal}{{\mathcal H}}
\renewcommand{\O}{{\mathbb O}}
\newcommand{\cgot}{{\mathfrak c}}
\newcommand{\mgot}{{\mathfrak m}}
\newcommand{\rgot}{{\mathfrak r}}
\newcommand{\qgot}{{\mathfrak q}}
\newcommand{\fgot}{{\mathfrak f}}
\newcommand{\agot}{{\mathfrak a}}
\newcommand{\bgot}{{\mathfrak b}}
\newcommand{\pgot}{{\mathfrak p}}
\newcommand{\Agot}{{\mathfrak A}}
\newcommand{\Bgot}{{\mathfrak B}}
\newcommand{\Cgot}{{\mathfrak C}}
\newcommand{\Ugot}{{\mathfrak U}}
\newcommand{\Vgot}{{\mathfrak V}}
\newcommand{\Fgot}{{\mathfrak F}}
\newcommand{\Sgot}{{\mathfrak S}}
\newcommand{\Mgot}{{\mathfrak M}}
\newcommand{\Rgot}{{\mathfrak R}}
\newcommand{\Pgot}{{\mathfrak P}}
\newcommand{\Tgot}{{\mathfrak T}}
\newcommand{\hartogs}[1]{{\mathfrak h(#1)}}
\newcommand{\Q}{{\mathbb Q}}
\newcommand{\N}{{\mathbb N}}
\newcommand{\f}{{\mathfrak f}}
\newcommand{\g}{{\mathfrak g}}
\renewcommand{\le}{\leqslant}
\renewcommand{\ge}{\geqslant}
\newcommand{\lep}{\preccurlyeq}
\newcommand{\gep}{\succcurlyeq}
\renewcommand{\bar}[1]{\overline{#1}}
\newcommand{\const}{\mathrm{const}}
\newcommand{\diam}{\mathrm{diam}}
\newcommand{\step}{\mathrm{st}}
\newcommand{\Spec}{\mathrm{Spec}}
\newcommand{\id}{\mathrm{id}}
\newcommand{\supp}{\mathop{\mathrm{supp}}}
\newcommand{\dddoteq}{\raisebox{-.1em}{\textbf{=}}\nolinebreak\hspace{-.85em}\raisebox{.4em}{...}}
%\newcommand{\eqdef}{\,\mathop{\color{red}\overset{\mbox{\tiny def}}{\rightleftharpoons}}\,}
%\newcommand{\eqrule}[1]{\,\mathop{\overset{\mbox{\tiny $#1$}}{=}}\,}
\newcommand{\lradef}{\,\mathop{\color{red}\overset{\mbox{\tiny def}}{\leftrightarrow}}\,}
\newcommand{\Code}{\mathop{\mathrm{Code}}\nolimits}
\newcommand{\cov}{\mathop{\mathrm{cov}}}
\newcommand{\upuparrow}{\mathop{\uparrow\uparrow}}
\newcommand{\Trans}{\mathop{\mathrm{Trans}}}
\newcommand{\Ord}{\mathop{\mathrm{Ord}}}
\renewcommand{\atop}[2]{\genfrac{}{}{0pt}{1}{#1}{#2}}
\newcommand{\ceil}[1]{\left\lceil#1\right\rceil}
\newcommand{\AC}{\ifmmode \mathsf{\color{darkred}AC}\else ${\mathsf{\color{darkred}AC}}$\fi\xspace}
\newcommand{\DC}{\ifmmode \mathsf{DC}\else ${\mathsf{DC}}$\fi\xspace}
\newcommand{\AD}{\ifmmode \mathsf{AD}\else ${\mathsf{AD}}$\fi\xspace}
\newcommand{\SI}{\ifmmode \mathsf{SI}\else ${\mathsf{SI}}$\fi\xspace}
\newcommand{\ZL}{\ifmmode \mathsf{\color{darkred}ZL}\else ${\mathsf{\color{darkred}ZL}}$\fi\xspace}
\newcommand{\HMP}{\ifmmode \mathsf{\color{darkred}HMP}\else ${\mathsf{\color{darkred}HMP}}$\fi\xspace}
\newcommand{\ZT}{\ifmmode \mathsf{\color{darkred}ZT}\else ${\mathsf{\color{darkred}ZT}}$\fi\xspace}
\newcommand{\CH}{\ifmmode \mathsf{\color{darkred}CH}\else ${\mathsf{\color{darkred}CH}}$\fi\xspace}
\newcommand{\GB}{\ifmmode \mathsf{\color{darkred}GB}\else ${\mathsf{\color{darkred}GB}}$\fi\xspace}
\newcommand{\GC}{\ifmmode \mathsf{GC}\else ${\mathsf{GC}}$\fi\xspace}
\newcommand{\ZF}{\ifmmode \mathsf{\color{darkred}ZF}\else ${\mathsf{\color{darkred}ZF}}$\fi\xspace}
\newcommand{\Zermelo}{\ifmmode \mathsf{\color{darkred}Z}\else ${\mathsf{\color{darkred}Z}}$\fi\xspace}
\newcommand{\HF}{\ifmmode \mathsf{\color{darkred}HF}\else ${\mathsf{\color{darkred}HF}}$\fi\xspace}
\newcommand{\Qar}{\ifmmode \mathsf{\color{darkred}Q}\else ${\mathsf{\color{darkred}Q}}$\fi\xspace}
\newcommand{\ZFC}{\ifmmode \mathsf{\color{darkred}ZFC}\else ${\mathsf{\color{darkred}ZFC}}$\fi\xspace}
\newcommand{\ZFA}{\ifmmode \mathsf{\color{darkred}ZFA}\else ${\mathsf{\color{darkred}ZFA}}$\fi\xspace}
\newcommand{\ZFD}{\ifmmode \mathsf{ZFD}\else ${\mathsf{ZFD}}$\fi\xspace}
\newcommand{\NGB}{\ifmmode \mathsf{NGB}\else ${\mathsf{NGB}}$\fi\xspace}
\newcommand{\Math}{\ifmmode {\color{darkred}\mathfrak M\mathfrak a \mathfrak t \mathfrak h}\else ${\color{darkred}\mathfrak M\mathfrak a \mathfrak t \mathfrak h}$\fi\xspace}
\newcommand{\GCH}{\ifmmode \mathsf{GCH}\else ${\mathsf{\color{darkred}GCH}}$\fi\xspace}
\newcommand{\PA}{\ifmmode \mathsf{\color{darkred}PA}\else ${\mathsf{\color{darkred}PA}}$\fi\xspace}
\newcommand{\RCF}{\ifmmode \mathsf{\color{darkred}RCF}\else ${\mathsf{\color{darkred}RCF}}$\fi\xspace}
\newcommand{\ACF}{\ifmmode \mathsf{\color{darkred}ACF}\else ${\mathsf{\color{darkred}ACF}}$\fi\xspace}
\newcommand{\DLO}{\ifmmode \mathsf{\color{darkred}DLO}\else ${\mathsf{\color{darkred}DLO}}$\fi\xspace}
\newcommand{\Hart}{\ifmmode \mathsf{Hart}\else $\mathsf{Hart}$\fi\xspace}
\newcommand{\DAG}{\ifmmode \mathsf{\color{darkred}DAG}\else ${\mathsf{\color{darkred}DAG}}$\fi\xspace}
\newcommand{\DTFA}{\ifmmode \mathsf{\color{darkred}DTFA}\else ${\mathsf{\color{darkred}DTFA}}$\fi\xspace}
\newcommand{\EA}{\ifmmode \mathsf{EA}\else ${\mathsf{EA}}$\fi\xspace}
\newcommand{\versus}{\quad\mathsf{vs}\quad}
\newcommand{\Rel}{{\mathsf{R}}}
\newcommand{\Srel}{{\mathsf{S}}}
\newcommand{\Fsf}{{\mathsf{F}}}
\newcommand{\Ep}{{\mathsf{\varepsilon}}}
\newcommand{\Sb}{{\mathsf{S}}}
\newcommand{\Aut}{{\mathsf{Aut}}}
\newcommand{\On}{\mathop{\mathrm{On}}}
% unicode-math резервирует \Succ под символ U+2ABC («double succeeds»), без глифа в XCharter Math;
% переопределяем после инициализации символов:
\AtBeginDocument{\renewcommand{\Succ}{\mathop{\mathrm{Succ}}}}
% Краткая запись для окружения proof (amsthm); иначе \pf/\epf не определены
\newcommand{\pf}{\begin{proof}}
\newcommand{\epf}{\end{proof}}
\newcommand{\Pred}{\mathop{\mathrm{Pred}}}
\newcommand{\Lim}{\mathop{\mathrm{Lim}}}
\newcommand{\cf}{\mathop{\mathrm{cf}}}
\newcommand{\dom}{\mathop{\mathrm{dom}}}
\newcommand{\Sch}[1]{\mathop{\mathrm{Sch}#1}}
\newcommand{\content}{\mathop{\mathrm{cont}}}
\newcommand{\Ker}{\mathop{\mathrm{Ker}}}
\newcommand{\sgn}{\mathop{\mathrm{sgn}}}
\newcommand{\ran}{\mathop{\mathrm{ran}}}
\newcommand{\rem}{\mathop{\mathrm{rem}}}
\newcommand{\Con}{\mathop{\mathrm{Con}}}
\newcommand{\Th}{\mathop{\mathrm{Th}}}
\newcommand{\Exp}{\mathop{\mathrm{Exp}}}
\newcommand{\ssup}{\mathop{\mathrm{ssup}}}
\newcommand{\Card}{\mathop{\mathrm{Card}}}
\newcommand{\SO}{\mathop{\mathrm{SO}}}
\newcommand{\GL}{\mathop{\mathrm{GL}}}
\newcommand{\Galois}{{\mathsf{G}}}
\newcommand{\Shift}{{\mathsf{Shift}}}
\newcommand{\PGL}{\mathop{\mathrm{PGL}}}
\newcommand{\PSL}{\mathop{\mathrm{PSL}}}
\newcommand{\SL}{\mathop{\mathrm{SL}}}
\newcommand{\SU}{\mathop{\mathrm{SU}}}
\newcommand{\Python}{\hbox{\textsf{Python}}}
\newcommand{\pr}{\mathrm{pr}}
\renewcommand{\Pr}{\mathrm{Pr}}
\newcommand{\In}{\mathop{\mathrm{In}}}
\renewcommand{\exp}{\mathop{\mathrm{exp}}}
\renewcommand{\Exp}{\mathop{\mathrm{Exp}}}
\newcommand{\Res}{\mathop{\mathrm{Res}}}
\newcommand{\Numcode}[1]{\ulcorner #1\urcorner}
\newcommand{\numeral}[1]{\underline{#1}}
\newcommand{\isom}{\cong}
\newcommand{\ravnom}{\simeq}
\newcommand{\esssup}{\mathrm{ess } \sup}
%\renewcommand{\bar}[1]{\overline{#1}}
\newcommand{\D}{\Variance}
\renewcommand{\P}{\Prob}
\newcommand{\I}{\mathrm{I}}
\newcommand{\red}{\mathrm{red}}
\newcommand{\rad}{\mathrm{rad}}
\renewcommand{\S}{\mathbb S}
\newcommand{\Mod}{\mathsf M}
\newcommand{\specialcell}[2][c]{\begin{tabular}[#1]{@{}l@{}}#2\end{tabular}}
\newcommand{\rank}{\mathrm{rank}}
\newcommand{\Prog}{\mathrm{Prog}}
%\renewcommand{\gcd}{\mbox{НОД}}
%\renewcommand{\nok}{\mbox{НОК}}
\newcommand{\factorial}{!}
\newcommand{\uspace}{\sqcup}
\newcommand{\chk}{\textcolor{darkgreen}{\checkmark}}
\newcommand{\none}{\textcolor{darkred}{\textsf{x}}}
\newcommand{\minus}{\textcolor{gray}{--}} % в таблицах сравнения
\newcommand{\PAn}[1]{{\color{darkred}\textsf{PA}#1}}
%\newcommand{\iif}{\quad\;\Longleftrightarrow\quad\;}

% Настройка содержания без номеров страниц
\usepackage{tocloft}
\cftpagenumbersoff{section}
\cftpagenumbersoff{subsection}
\cftpagenumbersoff{subsubsection}

% Hyperref для ссылок и навигации (должен быть загружен последним)
\usepackage[colorlinks=true, linkcolor=darkblue, citecolor=darkgreen, urlcolor=darkred, pdfencoding=auto, linktoc=all]{hyperref}


\usepackage{xargs}

\newcommandx{\marginfigure}[3][1=0.5,2=0.5,usedefault]{%
\par\noindent\hspace*{-13cm}\hspace*{#2\oddsidemargin}%
    \raisebox{-#1\height}[0pt][0pt]{%
        \makebox[0pt][c]{%
            #3%
        }%
    }%
    \vspace{-\baselineskip}%
    \par\noindent\ignorespaces%
}

\begin{document}

\title{Архитектура математических моделей}
\maketitle

\phantomsection
%\addcontentsline{toc}{section}{\contentsname}
\pdfbookmark[0]{\contentsname}{contents}
\tableofcontents

\pagestyle{empty}

\section{Теория множеств}

\subsection{Аксиомы \ZF: исторический обзор и формулировки}

Теория множеств Цермело-Френкеля (\ZF) была разработана в ответ на кризис оснований математики начала XX века. Парадоксы наивной теории множеств (такие как парадокс Рассела) показали необходимость аксиоматического подхода.

\begin{definition}[Аксиомы \ZF]
Теория \ZF состоит из следующих аксиом:
\begin{enumerate}
    \item \textbf{Аксиома экстенсиональности:} Два множества равны, если они содержат одни и те же элементы:
    $$\forall x \forall y (\forall z (z \in x \leftrightarrow z \in y) \rightarrow x = y)$$
    
    \textbf{Исторический контекст:} Эта аксиома была явно сформулирована Фреге (1884) и отражает основную интуицию о том, что множество полностью определяется своими элементами. До этого множества часто определялись через свойства, что приводило к парадоксам.
    
    \item \textbf{Аксиома пары:} Для любых двух множеств существует множество, содержащее ровно эти два множества:
    $$\forall x \forall y \exists z \forall w (w \in z \leftrightarrow (w = x \lor w = y))$$
    
    \textbf{Исторический контекст:} Введена Цермело (1908) для построения упорядоченных пар $\langle x,y \rangle = \{\{x\}, \{x,y\}\}$, необходимых для определения функций и отношений.
    
    \item \textbf{Аксиома объединения:} Для любого множества существует множество, содержащее все элементы элементов этого множества:
    $$\forall x \exists y \forall z (z \in y \leftrightarrow \exists w (w \in x \land z \in w))$$
    
    \textbf{Исторический контекст:} Позволяет строить объединения семейств множеств, критически важно для работы с бесконечными структурами.
    
    \item \textbf{Аксиома степени множества:} Для любого множества существует множество всех его подмножеств:
    $$\forall x \exists y \forall z (z \in y \leftrightarrow z \subseteq x)$$
    
    \textbf{Исторический контекст:} Введена Цермело (1908). Ключевая аксиома для теоремы Кантора о том, что $|X| < |\Pcal(X)|$, что показывает иерархию мощностей.
    
    \item \textbf{Аксиома бесконечности:} Существует индуктивное множество (содержащее $\emptyset$ и замкнутое относительно операции $x \mapsto x \cup \{x\}$):
    $$\exists x (\emptyset \in x \land \forall y (y \in x \rightarrow y \cup \{y\} \in x))$$
    
    \textbf{Исторический контекст:} Введена Цермело (1908) для гарантии существования бесконечных множеств. До этого бесконечность часто постулировалась неявно, что создавало проблемы.
    
    \item \textbf{Схема аксиом выделения (ограниченного понимания):} Для любого свойства $\ph$ и множества $x$ существует множество элементов $x$, удовлетворяющих $\ph$:
    $$\forall x \exists y \forall z (z \in y \leftrightarrow (z \in x \land \ph(z)))$$
    
    \textbf{Исторический контекст:} Это схема (бесконечное семейство аксиом), а не одна аксиома. Введена для устранения парадокса Рассела: вместо «множество всех $x$ таких, что $\ph(x)$» мы берём «множество всех $x \in y$ таких, что $\ph(x)$», где $y$ уже существует. Это ограничение предотвращает парадоксы.
    
    \item \textbf{Схема аксиом подстановки:} Если $F$ --- функциональное отношение, то образ любого множества при $F$ является множеством:
    $$\forall x \exists! y \ph(x,y) \rightarrow \forall u \exists v \forall r (r \in v \leftrightarrow \exists s (s \in u \land \ph(s,r)))$$
    
    \textbf{Исторический контекст:} Добавлена Френкелем (1922) и независимо Сколемом (1922). Позволяет строить множества через «подстановку» элементов одного множества в другой. Критически важна для работы с ординалами и кардиналами.
    
    \item \textbf{Аксиома регулярности (фундированности):} Каждое непустое множество содержит элемент, не пересекающийся с этим множеством:
    $$\forall x (x \neq \emptyset \rightarrow \exists y (y \in x \land y \cap x = \emptyset))$$
    
    \textbf{Исторический контекст:} Введена фон Нейманом (1925) для предотвращения «плохих» множеств типа $x \in x$ или бесконечных нисходящих цепей $\ldots \in x_2 \in x_1 \in x_0$. Гарантирует, что отношение $\in$ является фундированным (well-founded), что позволяет использовать трансфинитную индукцию.
\end{enumerate}
\end{definition}

\begin{remark}
\textbf{Эволюция аксиоматики:}
\begin{itemize}
    \item \textbf{Цермело (1908):} Первая аксиоматизация теории множеств, включала аксиомы 1--6. Не содержала схемы подстановки, что ограничивало построение больших ординалов.
    \item \textbf{Френкель (1922):} Добавил схему подстановки, что позволило строить множества типа $V_{\omega+\omega}$ и выше. Это стало основой современного \ZF.
    \item \textbf{Сколем (1922):} Независимо предложил схему подстановки, подчеркнув её необходимость для теории ординалов.
    \item \textbf{фон Нейман (1925):} Добавил аксиому регулярности, завершив формирование \ZF.
\end{itemize}
\end{remark}

\subsection{Логический анализ аксиом \ZF: зависимости и избыточность}

\quad 

\begin{theorem}[Зависимости между аксиомами]\;
\begin{enumerate}
    \item \textbf{Аксиома пары выводима:} Из схемы аксиом подстановки и аксиомы степени множества можно вывести аксиому пары.
    
    \textbf{Доказательство (набросок):} Для множеств $x$ и $y$ рассмотрим функцию $F$ такую, что $F(\emptyset) = x$ и $F(\{\emptyset\}) = y$. Применяя схему подстановки к множеству $\{\emptyset, \{\emptyset\}\}$ (которое существует по аксиоме степени), получаем пару $\{x, y\}$.
    
    \item \textbf{Пустое множество выводимо:} Из аксиомы бесконечности можно вывести существование пустого множества.
    
    \textbf{Доказательство (набросок):} Аксиома бесконечности прямо утверждает существование пустого множества.
    
\end{enumerate}
\end{theorem}

\begin{theorem}[Роль аксиомы регулярности]
Аксиома регулярности имеет особый статус в \ZF:
\begin{enumerate}
    \item \textbf{Практическая неиспользуемость:} В большинстве теорем математики аксиома регулярности не используется. Обычные математические объекты (числа, функции, пространства) не требуют её применения.
    
    \item \textbf{Теоретическая необходимость:} Аксиома регулярности
    наводит порядок в универсуме всех множеств, а именно: с ее помощью доказывается, что всякое множество содержится в каком-то универсуме фон Неймана, а значит, имеет ординальный ранг.

    \item \textbf{Независимость:} Можно построить модели \ZF без регулярности, где существуют множества $x$ такие, что $x \in x$ (модель Куайна). Однако такие модели не имеют фон Неймановской структуры.
\end{enumerate}
\end{theorem}

\begin{definition}[Аксиома коллекции (схема)]
Для любого отношения $\ph(x, y)$ (не обязательно функционального) и множества $u$ существует множество $v$ такое, что для каждого $x \in u$, если существует $y$ с $\ph(x, y)$, то существует $y \in v$ с $\ph(x, y)$:
$$\forall u \exists v \forall x \in u (\exists y \ph(x,y) \rightarrow \exists y \in v \ph(x,y))$$

\textbf{Отличие от подстановки:} 
\begin{itemize}
    \item \textbf{Подстановка} требует функциональности: $\forall x \exists! y \ph(x,y)$ (единственность $y$ для каждого $x$). Тогда образ множества $u$ --- это множество всех $y$ таких, что $\ph(x,y)$ для некоторого $x \in u$.
    \item \textbf{Коллекция} работает с произвольным отношением $\ph(x,y)$ (может быть много $y$ для одного $x$). Она только гарантирует, что если для $x \in u$ существует хотя бы один $y$ с $\ph(x,y)$, то в $v$ найдётся хотя бы один такой $y$ (но не обязательно все такие $y$, и не обязательно единственный).
\end{itemize}
\end{definition}


\begin{definition}[Теория \Zermelo (теория Цермело)]
Теория \Zermelo (теория Цермело) получается из \ZF удалением схемы аксиом подстановки. Она включает:
\begin{itemize}
    \item Аксиомы экстенсиональности, пары, объединения, степени множества, бесконечности
    \item Схему аксиом выделения
\end{itemize}
\textbf{Исторический контекст:} Это оригинальная аксиоматика Цермело (1908), которая была достаточна для большинства приложений, но ограничивала построение больших ординалов.
\end{definition}

\begin{definition}[\ZFC]
Теория \ZFC получается добавлением к \ZF аксиомы выбора (\AC), которая будет рассмотрена ниже.
\end{definition}

\subsection{Аксиома Выбора (\AC) и Лемма Цорна (\ZL)}

\begin{definition}[Частично упорядоченное множество]
Частично упорядоченное множество (poset) --- это пара $\langle P, \le \rangle$, где $P$ --- множество, а $\le$ --- бинарное отношение на $P$, удовлетворяющее:
\begin{enumerate}
    \item \textbf{Рефлексивность:} $\forall x \in P: x \le x$
    \item \textbf{Транзитивность:} $\forall x, y, z \in P: (x \le y \land y \le z) \Rightarrow x \le z$
    \item \textbf{Антисимметричность:} $\forall x, y \in P: (x \le y \land y \le x) \Rightarrow x = y$
\end{enumerate}
\end{definition}

\begin{example}[Примеры частично упорядоченных множеств]\;
\begin{enumerate}
    \item \textbf{Линейный порядок:} Множество натуральных чисел $\N$ с обычным порядком $\le$. Любые два элемента сравнимы.
    
    \begin{center}
    \begin{tikzpicture}[node distance=1.5cm]
        \node (1) {$1$};
        \node[right of=1] (2) {$2$};
        \node[right of=2] (3) {$3$};
        \node[right of=3] (4) {$4$};
        \node[right of=4] (dots) {$\ldots$};
        \draw[->] (1) -- (2);
        \draw[->] (2) -- (3);
        \draw[->] (3) -- (4);
        \draw[->] (4) -- (dots);
    \end{tikzpicture}
    \end{center}
    
    \item \textbf{Частичный порядок (не все элементы сравнимы):} Множество подмножеств $\{a, b, c\}$ с отношением включения $\subseteq$.
     
    \marginfigure[][]{
        \begin{tikzpicture}[node distance=1.2cm]
            \node (abc) {$\{a,b,c\}$};
            \node[below left of=abc, xshift=-1cm] (ab) {$\{a,b\}$};
            \node[below of=abc] (ac) {$\{a,c\}$};
            \node[below right of=abc, xshift=1cm] (bc) {$\{b,c\}$};
            \node[below of=ab] (a) {$\{a\}$};
            \node[below of=ac] (b) {$\{b\}$};
            \node[below of=bc] (c) {$\{c\}$};
            \node[below of=b] (empty) {$\emptyset$};
            
            \draw[->] (empty) -- (a);
            \draw[->] (empty) -- (b);
            \draw[->] (empty) -- (c);
            \draw[->] (a) -- (ab);
            \draw[->] (a) -- (ac);
            \draw[->] (b) -- (ab);
            \draw[->] (b) -- (bc);
            \draw[->] (c) -- (ac);
            \draw[->] (c) -- (bc);
            \draw[->] (ab) -- (abc);
            \draw[->] (ac) -- (abc);
            \draw[->] (bc) -- (abc);
        \end{tikzpicture}%
    }
    
    Заметим, что $\{a\}$ и $\{b\}$ не сравнимы (ни одно не является подмножеством другого).
    
    \item \textbf{Дерево:} Множество вершин дерева с корнем, где $x \le y$ означает, что $x$ --- предок $y$.

    \marginfigure[][0.9]{
    \begin{tikzpicture}[node distance=1.2cm]
        \node (root) {$r$};
        \node[below left of=root] (v1) {$v_1$};
        \node[below of=root] (v2) {$v_2$};
        \node[below right of=root] (v3) {$v_3$};
        \node[below left of=v1] (v4) {$v_4$};
        \node[below of=v1] (v5) {$v_5$};
        
        \draw[->] (root) -- (v1);
        \draw[->] (root) -- (v2);
        \draw[->] (root) -- (v3);
        \draw[->] (v1) -- (v4);
        \draw[->] (v1) -- (v5);
    \end{tikzpicture}
    }
    Здесь $r \le v_1$, $r \le v_4$, но $v_2$ и $v_3$ не сравнимы с $v_1$.
\end{enumerate}
\end{example}

\begin{definition}[Цепь и верхняя грань]
Пусть $\langle P, \le \rangle$ --- частично упорядоченное множество.
\begin{enumerate}
    \item \textbf{Цепь} --- это подмножество $C \subseteq P$, в котором любые два элемента сравнимы: $\forall x, y \in C: (x \le y \lor y \le x)$
    \item \textbf{Верхняя грань} цепи $C$ --- это элемент $u \in P$ такой, что $\forall x \in C: x \le u$
    \item \textbf{Максимальный элемент} --- это элемент $m \in P$ такой, что $\nexists x \in P: m < x$ (где $m < x$ означает $m \le x \land m \neq x$)
\end{enumerate}
\end{definition}

\begin{example}[Цепи и максимальные элементы]
Рассмотрим частично упорядоченное множество подмножеств $\{a, b, c\}$ с включением:
\begin{enumerate}
    \item \textbf{Пример цепи:} $C = \{\emptyset, \{a\}, \{a,b\}, \{a,b,c\}\}$ --- это цепь, так как все элементы сравнимы по включению.
    
    \marginfigure{
    \begin{tikzpicture}[node distance=1cm]
        \node (empty) {$\emptyset$};
        \node[above of=empty] (a) {$\{a\}$};
        \node[above of=a] (ab) {$\{a,b\}$};
        \node[above of=ab] (abc) {$\{a,b,c\}$};
        \draw[->] (empty) -- (a);
        \draw[->] (a) -- (ab);
        \draw[->] (ab) -- (abc);
    \end{tikzpicture}
    }
    
    \item \textbf{Пример множества без цепи:} В множестве $\{\{a\}, \{b\}, \{c\}\}$ нет цепи длины больше 1, так как элементы попарно не сравнимы.
    
    \item \textbf{Максимальный элемент:} В множестве всех подмножеств $\{a, b, c\}$ элемент $\{a,b,c\}$ является максимальным (и единственным максимальным).
    
    \item \textbf{Множество с несколькими максимальными элементами:} $P = \{\{a\}, \{b\}, \{c\}, \{a,b\}\}$ с включением. Здесь $\{a,b\}$ и $\{c\}$ --- оба максимальные элементы, но они не сравнимы.
\end{enumerate}
\end{example}

\begin{theorem}[Лемма Цорна (\ZL)]
Если в частично упорядоченном множестве $\langle P, \le \rangle$ каждая цепь имеет верхнюю грань, то $P$ содержит максимальный элемент.
\end{theorem}

\begin{theorem}[Аксиома Выбора (\AC)]
Для любого семейства непустых множеств $\{X_i\}_{i \in I}$ существует функция выбора $f: I \to \bigcup_{i \in I} X_i$ такая, что $f(i) \in X_i$ для всех $i \in I$.

Эквивалентная формулировка: произведение непустых множеств непусто: если $\forall i \in I: X_i \neq \emptyset$, то $\prod_{i \in I} X_i \neq \emptyset$.
\end{theorem}

\begin{theorem}[Эквивалентность \AC и \ZL]
Аксиома Выбора (\AC) эквивалентна Лемме Цорна (\ZL) (в рамках \ZF).
\end{theorem}

\begin{theorem}[Парадокс Банаха-Тарского]
Шар в $\R^3$ можно разбить на конечное число частей, которые можно переставить (используя только изометрии) так, чтобы получить два шара того же радиуса, что и исходный.

\textbf{Примечание:} Доказательство использует Аксиому Выбора (\AC) и показывает её неконструктивный характер.
\end{theorem}

\begin{theorem}[Существование базиса]
В любом векторном пространстве существует базис (максимальное линейно независимое множество).

\textbf{Доказательство (набросок):} Рассмотрим частично упорядоченное множество всех линейно независимых подмножеств пространства, упорядоченное по включению. Каждая цепь имеет верхнюю грань (объединение элементов цепи). По Лемме Цорна (\ZL) существует максимальный элемент --- базис.
\end{theorem}

\begin{remark}
В моделировании \AC часто отвечает за возможность перехода от «локальной разрешимости» к «глобальному состоянию». Если мы моделируем агента, который должен сделать бесконечное количество выборов, \AC гарантирует (неконструктивно), что такая стратегия существует.
\end{remark}

\begin{theorem}[Сравнение подстановки и коллекции]\;
\begin{enumerate}
    \item \textbf{Коллекция влечёт подстановку над теорией \Zermelo:} Над теорией \Zermelo (т.е. \ZF без подстановки и регулярности) из схемы аксиом коллекции следует схема аксиом подстановки.
    
    \textbf{Доказательство:} Если отношение $\ph(x,y)$ функционально (т.е. $\forall x \exists! y \ph(x,y)$), то для каждого $x \in u$ существует единственный $y$ с $\ph(x,y)$. Аксиома коллекции гарантирует существование множества $v$ такого, что для каждого $x \in u$ существует $y \in v$ с $\ph(x,y)$. Поскольку $y$ единственен для каждого $x$, множество $v$ содержит в точности все такие $y$, то есть является образом $u$ при $\ph$. Этот образ можно выделить из $v$ при помощи аксиомы выделения, что и есть утверждение аксиомы подстановки.
    
    \item \textbf{Подстановка влечёт коллекцию над \Zermelo+\AC:} Над теорией \Zermelo с добавлением аксиомы выбора из схемы аксиом подстановки следует схема аксиом коллекции.
    
    \textbf{Доказательство:} Пусть дано произвольное отношение $\ph(x,y)$ и множество $u$. По аксиоме выбора можно выбрать функцию $f$ такую, что для каждого $x \in u$, если существует $y$ с $\ph(x,y)$, то $\ph(x, f(x))$ выполнено. Отношение $\langle x, f(x) \rangle$ функционально, поэтому к нему применима аксиома подстановки. Применяя подстановку, получаем образ множества $u$, который содержит хотя бы один подходящий $y$ для каждого $x \in u$, что и требуется аксиомой коллекции.
    
    \item \textbf{Независимость над \Zermelo:} Над теорией \Zermelo (без аксиомы выбора) подстановка не влечёт коллекцию. Существуют модели теории \Zermelo с полной схемой подстановки, но без схемы коллекции.
    
    \textbf{Объяснение:} Подстановка работает только с функциональными отношениями. Без аксиомы выбора нельзя выбрать функцию из произвольного отношения, поэтому подстановка не может дать коллекцию для нефункциональных отношений.
    
    \item \textbf{Практическое значение:} 
    \begin{itemize}
        \item В теории \Zermelo коллекция является более сильной аксиомой, чем подстановка.
        \item В теории \Zermelo+\AC (и тем более в \ZFC) коллекция и подстановка равносильны, что объясняет, почему в классической теории множеств часто используется только подстановка.
        \item В конструктивной математике и теориях без аксиомы выбора коллекция часто предпочтительнее подстановки.
    \end{itemize}
\end{enumerate}
\end{theorem}

\begin{exercise}[Максимальные элементы и цепи]\; 
\begin{enumerate}
    \item \textbf{Максимальные элементы:} Рассмотрим частично упорядоченное множество всех двух- и трехэлементных подмножеств $\{a, b, c, d\}$, упорядоченное по включению $\subseteq$.
    \begin{enumerate}
        \item Нарисуйте диаграмму Хассе для этого множества.
        \item Сколько элементов в этом множестве?
        \item Перечислите все максимальные элементы этого множества.
        \item Перечислите все минимальные элементы.
    \end{enumerate}
    
    \item \textbf{Цепи:} В частично упорядоченном множестве подмножеств $\{a, b, c\}$:
    \begin{enumerate}
        \item Найдите все максимальные цепи (сквозные цепи от минимального до максимального элемента).
        \item Сколько всего максимальных цепей в этом множестве?
        \item Сколько элементов в самой длинной цепи?
    \end{enumerate}
    
    \item \textbf{Максимальные элементы в простом poset:} Рассмотрим частично упорядоченное множество $\{1, 2, 3, 4, 6, 12\}$ с отношением делимости (т.е. $x \le y$ если $x$ делит $y$).
    \begin{enumerate}
        \item Постройте диаграмму Хассе для этого множества.
        \item Является ли это множество изоморфным какому-либо булеану (с отношением вложения)? Почему?
        \item Постройте множество корневых цепей (цепей, начинающихся с минимального элемента) этого множества и укажите все его максимальные элементы (по отношению продолжения цепей).
        \item Какого вида графом является это множество корневых цепей?
    \end{enumerate}
    \item Приведите примеры максимальных антицепей из предыдущих трех задач.
    \item \textbf{Теорема Шпернера:}
        Пусть $\Pcal(n)$ — множество всех подмножеств множества $\{1, \dots, n\}$, упорядоченное по включению $\subseteq$.
        Докажите, что для любого $n$ самая длинная антицепь (максимальное число взаимоисключающих вариантов) достигается на семействе всех подмножеств мощности $\lfloor n/2 \rfloor$.
   
    \item \textbf{Бесконечные антицепи в деревьях:}
        Рассмотрим бесконечное двоичное дерево $\Tcal = \{0, 1\}^*$ (множество всех конечных последовательностей нулей и единиц). Порядок $x \le y$, если $x$ — начало последовательности $y$.
        \begin{enumerate}
            \item Приведите пример бесконечной цепи в этом дереве.
            \item Приведите пример бесконечной антицепи (бесконечного множества элементов, где ни один не является префиксом другого).
        \end{enumerate}
    
        \item \textbf{Линеаризация Иерархии (Код Прюфера):}
    Любое конечное дерево с пронумерованными вершинами можно взаимно-однозначно закодировать последовательностью чисел. Это показывает, как сложная топология сводится к линейному тексту.
    
    \textbf{Алгоритм кодирования:} Пока в дереве больше 2 вершин:
    \begin{itemize}
        \item Находим лист (вершину со степенью 1) с \textit{наименьшим} номером.
        \item Записываем номер его единственного соседа в код.
        \item Удаляем этот лист из дерева.
    \end{itemize}
    
    \begin{enumerate}
        \item \textbf{Прямая задача (Сжатие):} 
        Дано дерево с вершинами $\{1, 2, 3, 4, 5, 6\}$:
        \begin{center}
        \begin{tikzpicture}[node distance=1.5cm, main/.style = {draw, circle}] 
        \node[main] (1) {1}; 
        \node[main] (2) [below left of=1] {2};
        \node[main] (3) [below right of=1] {3}; 
        \node[main] (4) [below left of=2] {4};
        \node[main] (5) [below right of=2] {5};
        \node[main] (6) [right of=3] {6};
        \draw (1) -- (2);
        \draw (1) -- (3);
        \draw (2) -- (4);
        \draw (2) -- (5);
        \draw (3) -- (6);
        \end{tikzpicture} 
        \end{center}
        Вычислите его код Прюфера. (Подсказка: длина кода должна быть $n-2 = 4$ числа).
        
        \item \textbf{Обратная задача (Восстановление):}
        Дан код Прюфера: $\langle 3, 3, 3, 3 \rangle$. 
        Восстановите структуру дерева из 6 вершин. Какую топологию оно имеет (цепь, звезда, вилка)?
        
        \item \textbf{Интерпретация:}
        Код Прюфера имеет длину $n-2$, а алфавит состоит из $n$ символов.
        Общее количество таких кодов — $n^{n-2}$.
        Что это говорит нам о количестве возможных иерархий, которые можно построить на $n$ фиксированных точках? Быстро ли растет это число с ростом $n$?
    \end{enumerate}
\end{enumerate}
\end{exercise}

\begin{exercise}[Функции в микро-вселенной]\;

В этом упражнении элементы обозначаются строчными буквами, множества --- заглавными.

\textbf{Вселенная элементов:} $U = \{a, b, c, d, e, f, g, h\}$.

\textbf{Именованные множества:}
$$A = \{a, b, c\}, \quad B = \{b, c, d\}, \quad C = \{c, d, e, f\}, \quad D = \{e, f, g\}, \quad E = \{a, f, g, h\}, \quad F = \{b, d, f, h\}.$$

Определим функцию $\ph: U \to U$ набором пар:
$$\ph = \{\langle a,c \rangle,\; \langle b,d \rangle,\; \langle c,b \rangle,\; \langle d,e \rangle,\; \langle e,a \rangle,\; \langle f,d \rangle,\; \langle g,h \rangle,\; \langle h,a \rangle\}$$
(Напомним: $\langle x,y \rangle \in \ph$ означает $\ph(x) = y$. Например, $\ph(a) = c$.)

\begin{enumerate}
\item \textbf{Разминка: значения функции.}
\begin{enumerate}[label=(\roman*)]
    \item Вычислите $\ph(b)$, $\ph(e)$, $\ph(g)$.
    \item Верно ли, что $\ph(a) \in A$?
    \item Верно ли, что $\ph(a) \in B$?
    \item Вычислите $\ph(\ph(a))$ и $\ph(\ph(\ph(a)))$.
    \item Верно ли, что $\forall x \in A\; (\ph(x) \in B)$?
    \item Существует ли $x \in U$ такой, что $\ph(x) = f$? А такой, что $\ph(x) = g$?
\end{enumerate}

\item \textbf{Образы и прообразы.}

Напомним: \textbf{образ множества} --- $\ph[S] = \{\ph(x) \mid x \in S\}$.
\textbf{Прообраз множества} --- $\ph^{-1}[S] = \{x \in U \mid \ph(x) \in S\}$.

Вычислите:
\begin{enumerate}[label=(\roman*), start=7]
    \item $\ph[A]$. Совпадает ли результат с одним из множеств $A$--$F$?
    \item $\ph[B]$
    \item $\ph[D]$
    \item $\ph[\emptyset]$
    \item $\ph^{-1}[\{a\}]$ \quad (все элементы, которые $\ph$ отправляет в $a$)
    \item $\ph^{-1}[\{d\}]$
    \item $\ph^{-1}[A]$
    \item $\ph^{-1}[D]$
\end{enumerate}

\item \textbf{Образы и теоретико-множественные операции.}

Вычислите обе стороны и сравните --- совпадают ли?
\begin{enumerate}[label=(\roman*), start=15]
    \item $\ph[A \cap C]$ и $\ph[A] \cap \ph[C]$
    \item $\ph[A \cup B]$ и $\ph[A] \cup \ph[B]$
    \item $\ph[A \setminus B]$ и $\ph[A] \setminus \ph[B]$
    \item $\ph[B \setminus C]$ и $\ph[B] \setminus \ph[C]$
\end{enumerate}

\item \textbf{Прообразы и теоретико-множественные операции.}

Аналогичное задание для прообразов:
\begin{enumerate}[label=(\roman*), start=19]
    \item $\ph^{-1}[A \cap C]$ и $\ph^{-1}[A] \cap \ph^{-1}[C]$
    \item $\ph^{-1}[A \cup C]$ и $\ph^{-1}[A] \cup \ph^{-1}[C]$
    \item $\ph^{-1}[A \setminus C]$ и $\ph^{-1}[A] \setminus \ph^{-1}[C]$
\end{enumerate}
Заметили ли вы разницу в поведении образов и прообразов? 

Для каких операций ($\cap$, $\cup$, $\setminus$) равенство сохраняется, а для каких нет? \textbf{Укажите все случаи.}

\item \textbf{Образ прообраза и прообраз образа.}
\begin{enumerate}[label=(\roman*), start=22]
    \item $\ph^{-1}[\ph[A]]$. Содержит ли результат $A$? Равен ли $A$?
    \item $\ph[\ph^{-1}[B]]$. Содержится ли результат в $B$? Равен ли $B$?
    \item $\ph[\ph[A]]$
\end{enumerate}

\item \textbf{Свойства функции $\ph$.}
\begin{enumerate}[label=(\roman*), start=25]
    \item Является ли $\ph$ инъекцией (разные элементы $\to$ разные образы)? Если нет, приведите контрпример.
    \item Является ли $\ph$ сюръекцией (каждый элемент $U$ является образом)? Если нет, какие элементы не попадают в область значений?
    \item Является ли $\ph$ биекцией?
\end{enumerate}

\item \textbf{Композиция функций.}

Определим вторую функцию $\psi: U \to U$:
$$\psi = \{\langle a,f \rangle,\; \langle b,a \rangle,\; \langle c,g \rangle,\; \langle d,b \rangle,\; \langle e,c \rangle,\; \langle f,h \rangle,\; \langle g,e \rangle,\; \langle h,d \rangle\}$$
\begin{enumerate}[label=(\roman*), start=28]
    \item Вычислите $(\psi \circ \ph)(a)$, $(\psi \circ \ph)(b)$, $(\psi \circ \ph)(e)$.
    \item Вычислите $(\psi \circ \ph)[A]$ и $\psi[\ph[A]]$. Совпадают ли?
    \item Вычислите $(\psi \circ \ph)^{-1}[D]$ и $\ph^{-1}[\psi^{-1}[D]]$. Совпадают ли?
    \item Является ли $\psi$ инъекцией? Сюръекцией? Биекцией?
    \item Является ли композиция $\psi \circ \ph$ инъекцией? Сюръекцией?
\end{enumerate}

\end{enumerate}
\end{exercise}

\subsection{Ординалы и кардиналы}

Если кардинальные числа описывают «размер» бесконечности, то ординальные числа (ординалы) описывают её «структуру» и «этапы». Для моделирования иерархий и процессов ординалы являются более тонким инструментом.

\subsubsection{Вполне упорядоченные множества}

\quad

\begin{definition}[Вполне упорядоченное множество]
Линейно упорядоченное множество $\langle A, \le \rangle$ называется \textit{вполне упорядоченным} (well-ordered), если любое его непустое подмножество имеет наименьший элемент.
\end{definition}

\quad 


\begin{example}
\begin{itemize}
    \item $\N = \{0, 1, 2, \dots\}$ — вполне упорядочено.
    \item $\Z = \{\dots, -1, 0, 1, \dots\}$ — не вполне упорядочено (множество всех целых чисел не имеет минимума).
    \item $[0, 1]$ (интервал вещественных чисел) — не вполне упорядочено (интервал $(0, 1)$ не имеет минимума).
\end{itemize}
\end{example}

\begin{definition}[Ординал]
Ординал — это транзитивное множество, вполне упорядоченное отношением принадлежности $\in$. То есть, множество $\al$ является ординалом, если:
\begin{enumerate}
    \item Отношение $\in$ транзитивно на $\al$: если $x \in y \in \al$, то $x \in \al$.
    \item Отношение $\in$ строго линейно упорядочивает $\al$: для любых $x, y \in \al$ либо $x \in y$, либо $x = y$, либо $y \in x$.
    \item Любое непустое подмножество $\al$ имеет $\in$-наименьший элемент.
\end{enumerate}
\end{definition}

\begin{example}[Примеры ординалов]\;
\begin{enumerate}
    \item \textbf{Конечные ординалы:} Благодаря аксиоме бесконечности, можно построить последовательность конечных ординалов:
    \begin{itemize}
        \item $0 = \emptyset$ — пустое множество является ординалом.
        \item $1 = \{0\} = \{\emptyset\}$ — множество, содержащее только $\emptyset$, является ординалом.
        \item $2 = \{0, 1\} = \{\emptyset, \{\emptyset\}\}$ — множество $\{0, 1\}$ с порядком $\in$ изоморфно $\{0 < 1\}$.
        \item $n = \{0, 1, 2, \dots, n-1\}$ для любого натурального $n$.
    \end{itemize}
    
    \item \textbf{Первый бесконечный ординал:} $\om = \{0, 1, 2, \dots\}$ — множество всех конечных ординалов (натуральных чисел). Существование $\om$ гарантируется аксиомой бесконечности: мы берем пересечение всех индуктивных множеств.
    
    \item \textbf{Трансфинитные ординалы:}
    \begin{itemize}
        \item $\om + 1 = \om \cup \{\om\} = \{0, 1, 2, \dots, \om\}$ — следующий ординал после $\om$.
        \item $\om + 1 + 1 = (\om + 1) \cup \{\om + 1\} = \{0, 1, 2, \dots, \om, \om + 1\}$ — следующий ординал после $\om + 1$.
    \end{itemize}
\end{enumerate}
\end{example}

\begin{theorem}[Свойства ординалов]
\begin{enumerate}
    \item Любое вполне упорядоченное множество изоморфно единственному ординалу.
    \item Ординалы сравнимы: для любых $\al, \be$ либо $\al \in \be$, либо $\al = \be$, либо $\be \in \al$.
    \item Класс всех ординалов $\On$ не является множеством (парадокс Бурали-Форти).
    \item Ординал совпадает с множеством всех меньших ординалов: $\al = \{\be \mid \be \in \al\}$.
\end{enumerate}
\end{theorem}

\begin{theorem}[Трансфинитная индукция]
Пусть $P(\al)$ — некоторое свойство ординалов. Если для любого ординала $\al$ из того, что $P(\be)$ выполнено для всех $\be < \al$, следует $P(\al)$, то $P(\al)$ выполнено для всех ординалов $\al$.

\textbf{Формулировка через классы:} Если класс $C$ ординалов таков, что для любого ординала $\al$ из $\{\be \mid \be < \al\} \subseteq C$ следует $\al \in C$, то $C = \On$ (класс всех ординалов).

\textbf{Интуиция:} Если свойство «наследуется вверх» по ординалам (т.е. если оно выполнено для всех меньших ординалов, то оно выполнено и для данного), то оно выполнено для всех ординалов.
\end{theorem}

\begin{remark}
На практике трансфинитная индукция часто используется не на всей шкале ординалов, а начиная с некоторого ординала (например, с $\om$) или только для ординалов, меньших некоторого фиксированного ординала $\al$. В таких случаях база индукции формулируется для начального ординала, а шаг индукции — для всех последующих ординалов в рассматриваемом диапазоне.
\end{remark}

\begin{theorem}[Теорема о рекурсивных определениях (трансфинитная рекурсия)]
Пусть $\Phi$ — класс-функция (генератор рекурсии), которая каждой функции $g$ с областью определения, являющейся ординалом, ставит в соответствие некоторое значение $\Phi(g)$.

Тогда существует единственная класс-функция $F$ на ординалах такая, что для каждого ординала $\al$:
$$F(\al) = \Phi(F \upharpoonright \al),$$
где $F \upharpoonright \al$ — сужение функции $F$ на множество $\al$ (т.е. на все ординалы, меньшие $\al$).

\textbf{Интуиция:} Значение $F(\al)$ определяется на основе всех предыдущих значений $F(\be)$ для $\be < \al$.
\end{theorem}

Важно отметить, что ординалы бывают трех видов:
\begin{itemize}
    \item \textbf{Ноль:} $\al = 0$ (пустое множество)
    \item \textbf{Последователь (преемник):} $\al = \be + 1$ для некоторого ординала $\be$
    \item \textbf{Предельный:} $\al$ — предельный ординал (т.е. $\al > 0$ и для любого $\be < \al$ выполнено $\be + 1 < \al$)
\end{itemize}

Чаще всего рекурсивное определение задается разными формулами для этих трех классов ординалов. Это позволяет явно указать базовый случай (для нуля), шаг рекурсии (для последователей) и предельный случай (для предельных ординалов).

\subsubsection{Иерархия фон Неймана и ранг множества}

Иерархия фон Неймана ($V_\al$) — это способ построения всех множеств из пустого множества с помощью трансфинитной рекурсии.

\begin{definition}[Иерархия фон Неймана]
\marginfigure[1][0.7]{
\begin{tikzpicture}[scale=0.6, every node/.style={font=\tiny}]
    % Метки рангов в левом столбце (сдвинуты левее)
    \node at (-7.5,0) {\scriptsize $V_0$};
    \node at (-7.5,-1) {\scriptsize $V_1$};
    \node at (-7.5,-2) {\scriptsize $V_2$};
    \node at (-7.5,-3) {\scriptsize $V_3$};
    \node at (-7.5,-4) {\scriptsize $V_4$};
    
    % Ранг 0
    \node (v0) at (0,0) {$\emptyset$};
    
    % Ранг 1
    \node (v1) at (0,-1) {$\{\emptyset\}$};
    
    % Ранг 2
    \node (v2b) at (0,-2) {$\{\emptyset,\{\emptyset\}\}$};
    
    % Ранг 3 (увеличено расстояние)
    \node (v3b) at (-4,-3) {$\{\{\emptyset\}\}$};
    \node (v3d) at (4,-3) {$\{\emptyset,\{\emptyset\},\{\{\emptyset\}\}\}$};
    
    % Ранг 4 (новые элементы, не входящие в V₃, увеличено расстояние)
    \node (v4a) at (-5.5,-4) {$\{\emptyset,\{\{\emptyset\}\}\}$};
    \node (v4b) at (-2,-4) {$\{\{\emptyset\},\{\emptyset,\{\emptyset\}\}\}$};
    \node (v4c) at (3,-4) {$\{\{\{\emptyset\}\},\{\emptyset,\{\emptyset\}\}\}$};
    \node (v4d) at (7.5,-4) {$\{\emptyset,\{\emptyset\},\{\{\emptyset\}\},\{\emptyset,\{\emptyset\}\}\}$};
\end{tikzpicture}
}
Иерархия фон Неймана определяется трансфинитной рекурсией по ординалам:
\begin{itemize}
    \item $V_0 = \emptyset$
    \item $V_{\al + 1} = \Pcal(V_\al)$ (множество всех подмножеств $V_\al$)
    \item $V_\la = \bigcup_{\be < \la} V_\be$ для предельного ординала $\la$
\end{itemize}

Для натуральных чисел $n$ мощность $V_n$ вычисляется рекурсивно:
$$|V_0| = 0, \quad |V_{n+1}| = 2^{|V_n|},
\quad |V_{n+1}|=2\uparrow\uparrow n\mbox{ (нотация Кнута)}.$$
\end{definition}

\begin{theorem}[Фундаментальность иерархии фон Неймана]
Благодаря аксиоме регулярности, всякое множество $x$ попадает в некоторый универсум $V_\al$ иерархии фон Неймана. То есть, для любого множества $x$ существует ординал $\al$ такой, что $x \in V_{\al + 1}$.
\end{theorem}

\begin{definition}[Ранг множества]
Ранг множества $x$ (обозначается $\rank(x)$) определяется как наименьший ординал $\al$ такой, что $x \in V_{\al + 1}$. Эквивалентно:
\begin{itemize}
    \item $\rank(\emptyset) = 0$
    \item $\rank(x) = \sup\{\rank(y) + 1 \mid y \in x\}$ для непустого множества $x$
\end{itemize}

Для натуральных чисел $n$ количество множеств ранга $n$ равно:
$$|\{x \mid \rank(x) = n\}| = |V_{n+1}| - |V_n| = 2^{|V_n|} - |V_n|=2\uparrow\uparrow n-2\uparrow\uparrow (n-1).$$
\end{definition}

\begin{remark}
Иерархия фон Неймана отражает «уровень сложности» множества: чем выше ранг, тем «глубже» структура множества. Например:
\begin{itemize}
    \item Ранг конечного ординала $n$ равен $n$.
    \item Ранг ординала $\om$ равен $\om$.
    \item Ранг любого множества является ординалом.
\end{itemize}
\end{remark}

\subsubsection{Ординальная арифметика}

В отличие от обычных натуральных чисел, арифметика ординалов некоммутативна. Это отражает структуру процесса: «сделать бесконечно много шагов, а потом еще один» отличается от «сделать один шаг, а потом бесконечно много».

\begin{definition}[Сумма ординалов (рекурсивное определение)]
Сумма ординалов $\al + \be$ определяется трансфинитной рекурсией по второму аргументу $\be$:
\begin{itemize}
    \item $\al + 0 = \al$
    \item $\al + (\be + 1) = (\al + \be) + 1 = (\al + \be) \cup \{(\al + \be)\}$ — следующий ординал после $\al + \be$
    \item $\al + \la = \sup\{\al + \be \mid \be < \la\}$ для предельного ординала $\la$
\end{itemize}

\textbf{Интуиция:} Сумма $\al + \be$ — это результат приписывания $\be$ «после» $\al$. Для каждого шага рекурсии мы берем уже построенную сумму и добавляем следующий элемент.
\end{definition}


\begin{definition}[Произведение ординалов (рекурсивное определение)]
Произведение ординалов $\al \cdot \be$ определяется трансфинитной рекурсией по второму аргументу $\be$:
\begin{itemize}
    \item $\al \cdot 0 = 0$
    \item $\al \cdot (\be + 1) = \al \cdot \be + \al$ — приписывание $\al$ «после» уже построенного произведения
    \item $\al \cdot \la = \sup\{\al \cdot \be \mid \be < \la\}$ для предельного ординала $\la$
\end{itemize}

\textbf{Интуиция:} Произведение $\al \cdot \be$ можно представить как лексикографический порядок на $\be \times \al$: сначала все пары с первой координатой $0$, затем с первой координатой $1$, и т.д. Это соответствует сумме $\be$ копий ординала $\al$.
\end{definition}

\begin{definition}[Степень ординалов (рекурсивное определение)]
Степень ординалов $\al^{\be}$ определяется трансфинитной рекурсией по показателю $\be$:
\begin{itemize}
    \item $\al^0 = 1$
    \item $\al^{\be + 1} = \al^{\be} \cdot \al$ — умножение уже построенной степени на $\al$
    \item $\al^{\la} = \sup\{\al^{\be} \mid 0<\be < \la\}$ для предельного ординала $\la$
\end{itemize}

Все три операции (сумма, произведение, степень) определяются с помощью теоремы о рекурсивных определениях: для суммы генератор рекурсии $\Phi$ по функции $g$ определяет $\Phi(g) = \al + \sup(\ran(g)) + 1$ (если $\be$ — преемник) или $\Phi(g) = \sup(\ran(g))$ (если $\be$ — предел), аналогично для произведения и степени.
\end{definition}


\begin{theorem}[Непрерывность операций]
Арифметические операции над ординалами непрерывны по второму аргументу.
Пусть $\lambda$ — предельный ординал, и $\{\beta_\xi\}_{\xi < \lambda}$ — возрастающая последовательность ординалов с пределом $\beta = \sup_{\xi < \lambda} \beta_\xi$. Тогда:
\begin{enumerate}
    \item \textbf{Для сложения:} $\alpha + \sup_{\xi < \lambda} \beta_\xi = \sup_{\xi < \lambda} (\alpha + \beta_\xi)$ (для любых $\alpha$).
    \item \textbf{Для умножения:} $\alpha \cdot \sup_{\xi < \lambda} \beta_\xi = \sup_{\xi < \lambda} (\alpha \cdot \beta_\xi)$ (для любых $\alpha$).
    \item \textbf{Для степени:} $\alpha^{\sup_{\xi < \lambda} \beta_\xi} = \sup_{\xi < \lambda} (\alpha^{\beta_\xi})$ (для любых $\be_\xi>0$ и любых $\alpha$).
\end{enumerate}
\textbf{Геометрический смысл:} «Пристыковка» фиксированного хвоста $\alpha$ (или возведение в степень) к растущему множеству не меняет предельного поведения. Операция коммутирует с супремумом.
\end{theorem}



\begin{exercise}[Арифметика ординалов]
\begin{enumerate}
    \item \textbf{Доказать свойства суммы ординалов:}
    \begin{itemize}
        \item Ассоциативность: $(\al + \be) + \ga = \al + (\be + \ga)$
        \item Нейтральный элемент: $\al + 0 = 0 + \al = \al$
        \item Левая сократимость: если $\al + \be = \al + \ga$, то $\be = \ga$
        \item Правая сократимость неверна: показать, что $1 + \omega = 0 + \omega$, но $1 \neq 0$
    \end{itemize}
    
    \item \textbf{Проверить некоммутативность суммы:} Показать, что $1 + \omega = \omega$, но $\omega + 1 \neq \omega$. Объяснить геометрически (через порядковые типы).
    
    \item \textbf{Доказать свойства произведения ординалов:}
    \begin{itemize}
        \item Ассоциативность: $(\al \cdot \be) \cdot \ga = \al \cdot (\be \cdot \ga)$
        \item Дистрибутивность слева: $\al \cdot (\be + \ga) = \al \cdot \be + \al \cdot \ga$
        \item Нейтральный элемент: $1 \cdot \al = \al \cdot 1 = \al$
        \item Отсутствие делителей нуля: если $\al \cdot \be = 0$, то $\al = 0$ или $\be = 0$
        \item Левая сократимость: если $\al > 0$ и $\al \cdot \be = \al \cdot \ga$, то $\be = \ga$
        \item Правая сократимость неверна: показать, что $2 \cdot \omega = 1 \cdot \omega$, но $2 \neq 1$
    \end{itemize}
    
    \item \textbf{Проверить некоммутативность произведения:} Показать, что $2 \cdot \omega = \omega$, но $\omega \cdot 2 = \omega + \omega \neq \omega$.
    
    \item \textbf{Проверить недистрибутивность справа:} Показать, что $(1 + 1) \cdot \omega = 2 \cdot \omega = \omega$, но $1 \cdot \omega + 1 \cdot \omega = \omega + \omega \neq \omega$.
    
    \item \textbf{Вычислить простейшие примеры степени:}
    \begin{itemize}
        \item $0^0$, $0^\al$ (для $\al > 0$), $\al^0$
        \item Объяснить, почему $0^0 = 1$ по определению
    \end{itemize}
    
    \item \textbf{Доказать свойства степени ординалов:}
    \begin{itemize}
        \item $\al^{\be + \ga} = \al^{\be} \cdot \al^{\ga}$
        \item $(\al^{\be})^{\ga} = \al^{\be \cdot \ga}$
    \end{itemize}
    
    \item \textbf{Связь операций с порядком:}
    \begin{itemize}
        \item Доказать, что если $\al < \be$, то $\al + \ga \le \be + \ga$ и $\ga + \al \le \ga + \be$ (но строгое неравенство может не выполняться)
        \item Доказать, что если $\al < \be$ и $\ga > 0$, то $\al \cdot \ga \le \be \cdot \ga$ и $\ga \cdot \al \le \ga \cdot \be$ (но строгое неравенство может не выполняться)
        \item Доказать, что если $\al < \be$ и $\ga > 1$, то $\al^{\ga} \le \be^{\ga}$ и $\ga^{\al} \le \ga^{\be}$ (но строгое неравенство может не выполняться)
        \item Показать, что строгая монотонность справа не всегда верна: $2 < 3$, но $2 + \omega = 3 + \omega = \omega$, $2 \cdot \omega = 3 \cdot \omega = \omega$, $2^{\omega} = 3^{\omega}=\omega$.
    \end{itemize}
\end{enumerate}

\textit{Подсказка:} Для доказательства свойств используйте трансфинитную индукцию по соответствующему аргументу, опираясь на рекурсивные определения операций. Для проверки некоммутативности и недистрибутивности достаточно вычислить конкретные примеры по рекурсивным определениям.
\end{exercise}

\begin{theorem}[Канторовская нормальная форма ординалов]
    Всякий ординал $\al > 0$ может быть единственным образом представлен в виде:
    $$\al = \om^{\be_1} \cdot n_1 + \om^{\be_2} \cdot n_2 + \dots + \om^{\be_k} \cdot n_k,$$
    где $k \in \N$, $\be_1 > \be_2 > \dots > \be_k$ --- ординалы, и $n_1, n_2, \dots, n_k \in \N \setminus \{0\}$ --- натуральные числа.
    
    \textbf{Интуиция:} Это аналог позиционной системы счисления, но с основанием $\om$ вместо конечного числа. Каждый ординал записывается как «полином» по степеням $\om$.
    \end{theorem}
    
    \begin{remark}
    На основе канторовской нормальной формы можно определить \textit{коммутативные} операции сложения и умножения (Гессенберга) на ординалах (натуральная сумма и натуральное произведение), которые отличаются от обычных некоммутативных операций, рассмотренных выше. Эти коммутативные операции получаются путем «покомпонентного» сложения и умножения коэффициентов в нормальной форме, что делает их коммутативными и дистрибутивными с обеих сторон.
    \end{remark}
    
    \begin{remark}[Ординал $\ep_0$ и эффективность КНФ]
    Ординал $\ep_0$ определяется как наименьшая неподвижная точка функции $\al \mapsto \om^\al$, то есть $\ep_0 = \sup\{\om, \om^\om, \om^{\om^\om}, \dots\}$.
    
    Канторовская нормальная форма эффективно работает для всех ординалов, меньших $\ep_0$. Более того, любой ординал $\al < \ep_0$ может быть представлен как \textit{конечная запись} в алфавите, состоящем из символов $+$, $0$, $\om$ и показателя степени. Это означает, что такие ординалы имеют эффективное (вычислимое) представление, что делает их особенно важными для теории доказательств и вычислимости.
    \end{remark}
    
\begin{exercise}[Парадокс «Сделка с Чёртом»]
Бизнесмен заключил с чёртом сделку:
\begin{enumerate}
    \item У бизнесмена есть кошелек с \textbf{произвольным} конечным набором купюр любых натуральных достоинств.
    \item В любой момент он может обменять у чёрта \textbf{одну} купюру достоинства $k$ ($k>1$) на \textbf{любое конечное количество} купюр любых меньших достоинств ($<k$).
    \item Чёрт не ограничивает количество выдаваемых мелких купюр (их может быть сколь угодно много).
    \item Каждый день бизнесмен обязан тратить ровно 1 рубль (купюру достоинства 1). Если рубля нет, он вынужден разменять более крупную купюру.
\end{enumerate}
\textbf{Вопрос:} Докажите, что бизнесмен гарантированно обанкротится за конечное число дней, каким бы огромным ни был его стартовый капитал и какой бы хитрой стратегии размена он ни придерживался.
\end{exercise}

\begin{proof}[Решение: Метод спуска по ординалам]
Сопоставим состоянию кошелька ординал. Пусть у бизнесмена есть $a_k$ купюр достоинства $k$. Тогда состоянию кошелька соответствует ординал в Канторовской нормальной форме:
$$ \alpha = \omega^{k} \cdot a_k + \omega^{k-1} \cdot a_{k-1} + \dots + \omega \cdot a_1 + a_0 $$
где $k$ — максимальный номинал в текущий момент.

Рассмотрим динамику процесса:
\begin{itemize}
    \item \textbf{Трата рубля:} Уменьшает свободный член $a_0$. Ординал строго убывает.
    \item \textbf{Размен:} Мы заменяем одну купюру достоинства $m$ (слагаемое $\omega^m$) на конечный набор купюр меньших достоинств. В терминах ординалов это замена слагаемого $\omega^m$ на конечную сумму слагаемых вида $\omega^j$, где $j < m$.
\end{itemize}
По свойствам ординальной арифметики, для любого $m > 0$ и любых натуральных коэффициентов $N_j$:
$$ \omega^m > \sum_{j=0}^{m-1} \omega^j \cdot N_j $$
Следовательно, на каждом шаге ассоциированный ординал \textbf{строго убывает}: $\alpha_0 > \alpha_1 > \alpha_2 > \dots$
Поскольку множество ординалов до $\omega^\omega$ вполне упорядочено, этот процесс не может продолжаться бесконечно.
\end{proof}

\begin{remark}[Границы доказуемости: $\omega^\omega$ vs $\epsilon_0$]
Эта задача иллюстрирует разницу между «большим» и «недостижимо большим».
\begin{itemize}
    \item Множество всех возможных состояний кошелька имеет порядковый тип $\omega^\omega$. Фундированность этого множества (и завершение игры для любого старта) \textbf{доказуема} в Арифметике Пеано (\PA).
    \item Если бы мы разрешили купюрам иметь номиналы, являющиеся «купюрами от купюр» (древовидную структуру, как в задаче о Гидре Кирби-Пэриса), сложность системы выросла бы до ординала $\epsilon_0$. Теорема о том, что такой процесс всегда завершается, является истинной, но \textbf{недоказуемой} в \PA (теорема Гудстейна).
\end{itemize}
\end{remark}


\subsubsection{Мощности множеств и сравнение мощностей}

Для работы с «размером» множеств независимо от их структуры вводится понятие мощности.

\begin{definition}[Равномощные множества]
Два множества $A$ и $B$ называются \textit{равномощными} (обозначается $|A| = |B|$), если существует биекция $f: A \to B$ между ними.
\end{definition}

\begin{definition}[Сравнение мощностей]
Мощность множества $A$ не превосходит мощности множества $B$ (обозначается $|A| \le |B|$), если существует инъекция $f: A \to B$.

Мощность множества $A$ строго меньше мощности множества $B$ (обозначается $|A| < |B|$), если $|A| \le |B|$ и $|A| \neq |B|$.
\end{definition}

Для доказательства равенства мощностей двух множеств достаточно показать, что существуют инъекции в обе стороны. Это доказывают следующие теоремы.

\begin{theorem}[Кнастера-Тарского о неподвижной точке]
Пусть $L$ — полная решетка (т.е. любое подмножество имеет точную верхнюю и точную нижнюю грань) и $f: L \to L$ — монотонная функция (т.е. если $x \le y$, то $f(x) \le f(y)$). Тогда $f$ имеет неподвижную точку, т.е. существует $x \in L$ такой, что $f(x) = x$.

Множество всех неподвижных точек функции $f$ является полной подрешеткой $L$.
\end{theorem}

\begin{theorem}[Кантора-Бернштейна-Шредера]
Если существуют инъекции $f: A \to B$ и $g: B \to A$, то существует биекция между $A$ и $B$, т.е. $|A| = |B|$.

\textbf{Доказательство (набросок):} Применяя теорему Кнастера-Тарского к полной решетке подмножеств множества $A$ и монотонной функции, построенной из $f$ и $g$, получаем неподвижную точку, которая позволяет построить биекцию между $A$ и $B$.
\end{theorem}

\begin{example}[Применение теоремы Кантора-Бернштейна-Шредера]
\begin{enumerate}
    \item \textbf{Счетность множества целых чисел:} Множество $\Z$ равномощно $\N$. Инъекция $\N \to \Z$ очевидна (тождественное вложение), а инъекция $\Z \to \N$ задается, например, отображением $n \mapsto 2|n| + [n < 0]$, где $[n < 0]$ равно 1, если $n < 0$, и 0 иначе. По теореме Кантора-Бернштейна-Шредера $|\Z| = |\N|$.
    
    \item \textbf{Счетность множества рациональных чисел:} Множество $\Q$ также равномощно $\N$. Инъекция $\N \to \Q$ очевидна. Для построения инъекции $\Q \to \N$ заметим, что каждое рациональное число можно представить как дробь $\frac{p}{q}$, где $p \in \Z$, $q \in \Z \setminus \{0\}$, причем $\gcd(p,q) = 1$. Это задает инъекцию $\Q \to \Z \times \Z$ (отображение $\frac{p}{q} \mapsto \langle p,q \rangle$). Далее, $\Z \times \Z$ можно вложить в $\Z$ следующим образом: используя биекцию $\phi: \Z \to \N$ из предыдущего пункта, пару $\langle m,n \rangle \in \Z \times \Z$ кодируем как элемент $\N \times \N$ с помощью $\langle \phi(m), \phi(n) \rangle$, а затем применяем диагональную нумерацию (канторовскую) для $\N \times \N \to \N$: пару $\langle a,b \rangle \in \N \times \N$ кодируем как $\frac{(a+b)(a+b+1)}{2} + b$.
    
    \marginfigure[][0.7]{
        \begin{tikzpicture}[scale=1.5, >=stealth]
            % Настройки стилей
            \tikzset{
                grid node/.style={circle, draw=gray!40, fill=white, inner sep=2pt, minimum size=6mm, font=\small},
                path edge/.style={->, thick, blue!80!black, rounded corners=3pt},
                axis label/.style={font=\footnotesize, color=gray}
            }
        
            % 1. Рисуем координатную сетку и оси
            \draw[->, gray!50] (-0.5,0) -- (3.5,0) node[right, black] {$x$}; % Ось X
            \draw[->, gray!50] (0,-0.5) -- (0,3.5) node[above, black] {$y$}; % Ось Y
        
            % Подписи координат
            \foreach \i in {0,1,2,3} {
                \node[axis label] at (\i, -0.3) {$\i$};
                \node[axis label] at (-0.3, \i) {$\i$};
            }
        
            % 2. Размещаем узлы с номерами (n)
            % Формула для змейки (Zig-Zag):
            % Если x+y четное: направление вверх-вправо, если нечетное: вниз-влево.
            % Здесь просто расставим узлы в соответствии с визуальным путем для наглядности.
            
            % Координаты (x,y) и значения n
            \node[grid node] (n00) at (0,0) {$0$};
            
            \node[grid node] (n01) at (0,1) {$1$};
            \node[grid node] (n10) at (1,0) {$2$};
            
            \node[grid node] (n20) at (2,0) {$3$};
            \node[grid node] (n11) at (1,1) {$4$};
            \node[grid node] (n02) at (0,2) {$5$};
            
            \node[grid node] (n03) at (0,3) {$6$};
            \node[grid node] (n12) at (1,2) {$7$};
            \node[grid node] (n21) at (2,1) {$8$};
            \node[grid node] (n30) at (3,0) {$9$};
        
            % Остальные пустые узлы для полноты картины 4x4
            \foreach \x in {0,...,3} {
                \foreach \y in {0,...,3} {
                    % Если узел еще не определен (сумма индексов > 3), рисуем пустую точку
                    \ifnum\numexpr\x+\y\relax>3
                        \node[circle, fill=gray!20, inner sep=1pt, minimum size=2mm] at (\x,\y) {};
                    \fi
                }
            }
        
            % 3. Рисуем непрерывный путь (Змейка)
            % Порядок: (0,0) -> (0,1) -> (1,0) -> (2,0) -> (1,1) -> (0,2) -> (0,3) -> ...
            \draw[path edge] (n00) -- (n01);
            \draw[path edge] (n01) -- (n10);
            \draw[path edge] (n10) -- (n20); % Переход на новую диагональ
            \draw[path edge] (n20) -- (n11);
            \draw[path edge] (n11) -- (n02);
            \draw[path edge] (n02) -- (n03); % Переход на новую диагональ
            \draw[path edge] (n03) -- (n12);
            \draw[path edge] (n12) -- (n21);
            \draw[path edge] (n21) -- (n30);
            
            % Стрелка в бесконечность
            \draw[dashed, blue!50, ->] (n30) -- (3.5, -0.5);
        
            % Название
            \node[align=center] at (1.5, -1.0) {Диагональная нумерация Кантора\\(зигзаг)};
        \end{tikzpicture}
    }
    
    Поскольку $|\Z| = |\N|$ (из предыдущего пункта), получаем $|\Q| = |\N|$.
    
    \item \textbf{Равномощность интервала и прямой:} Множество $[0,1]$ равномощно $\R$. Инъекция $[0,1] \to \R$ очевидна (тождественное вложение), а инъекция $\R \to [0,1]$ можно построить, например, с помощью функции $\arctan$ с подходящим масштабированием. По теореме Кантора-Бернштейна-Шредера $|[0,1]| = |\R|$.
\end{enumerate}
\end{example}

\begin{theorem}[Кантора]
    Для любого множества $X$ справедливы следующие равносильные утверждения (эквивалентность не требует \AC):
    \begin{enumerate}
        \item Не существует сюръекции $f: X \to \Pcal(X)$, где $\Pcal(X)$ --- множество всех подмножеств $X$.
        \item Не существует инъекции $g: \Pcal(X) \to X$.
    \end{enumerate}
    
    \textbf{Следствие:} $|X| < |\Pcal(X)|$ для любого множества $X$. В частности, $|\N| < |\Pcal(\N)| = 2^{\aleph_0}$.
    \end{theorem}

\subsubsection{Кардиналы как начальные ординалы}

\begin{definition}[Кардинал]
Кардинал — это \textit{начальный ординал}, то есть наименьший ординал данной мощности.
\begin{itemize}
    \item $|\omega| = |\omega + 1| = |\omega^2| = \aleph_0$ (все они счетны).
    \item $\omega_1$ (или $\Omega$) — первый несчетный ординал. Это множество всех счетных ординалов. Его мощность обозначается $\aleph_1$.
\end{itemize}
\end{definition}

\begin{theorem}[Теорема Цермело (эквивалент \AC)]
Всякое множество может быть вполне упорядочено. То есть, для любого множества $X$ существует отношение порядка $\le$ на $X$, которое делает $\langle X, \le \rangle$ вполне упорядоченным множеством.

Эта теорема эквивалентна аксиоме выбора (\AC) в рамках \ZF.
\end{theorem}

\begin{theorem}[Следствие из \AC]
Всякое множество равномощно некоторому кардиналу. Такой кардинал называется \textit{мощностью} данного множества и обозначается $|X|$ для множества $X$.

\textbf{Доказательство:} По теореме Цермело, множество $X$ может быть вполне упорядочено. По теореме о свойствах ординалов, всякое вполне упорядоченное множество изоморфно единственному ординалу $\al$. Тогда $X$ равномощно $\al$. Пусть $\ka$ — наименьший ординал, равномощный $\al$ (т.е. начальный ординал мощности $|\al|$). Тогда $\ka$ является кардиналом, и $|X| = \ka$.
\end{theorem}

\begin{remark}[Кардинал Хартогса]
В отсутствие аксиомы выбора (\AC) нельзя гарантировать, что всякое множество равномощно некоторому кардиналу. Однако можно показать (без \AC), что для любого множества $X$ существует \textit{кардинал Хартогса} $\Hart(X)$ — наименьший ординал, который не может быть вложен в $X$ (т.е. не существует инъекции $\Hart(X) \to X$). 

Кардинал Хартогса $\Hart(X)$ служит верхней кардинальной оценкой для $X$: если $X$ равномощно некоторому кардиналу $\ka$, то $\ka < \Hart(X)$. Более того, $\Hart(X)$ всегда является кардиналом (начальным ординалом), и его существование доказывается без использования \AC.

Если $X=\tau$ --- бесконечный кардинал, то $\Hart(\tau)>\tau$.
\end{remark}

\begin{theorem}[Базовые свойства кардиналов]
\begin{enumerate}
    \item \textbf{Собственный класс:} Класс всех кардиналов не является множеством (аналогично классу всех ординалов).
    
    \item \textbf{Вполне упорядоченность:} Кардиналы вполне упорядочены отношением принадлежности $\in$ (или, что эквивалентно, отношением $\le$ на мощностях). То есть, для любых кардиналов $\ka$ и $\la$ выполнено либо $\ka \in \la$, либо $\ka = \la$, либо $\la \in \ka$.
    
    \item \textbf{Объединение и пересечение:} Если $A$ — множество кардиналов, то $\bigcup A$ и $\bigcap A$ (если $A \neq \emptyset$) также являются кардиналами.
    
    \item \textbf{Предельность бесконечных кардиналов:} Всякий бесконечный кардинал является предельным ординалом. То есть, если $\ka$ — бесконечный кардинал, то не существует ординала $\al$ такого, что $\ka = \al + 1$.
\end{enumerate}
\end{theorem}

\begin{definition}[Кардинал, следующий за данным]
Для любого кардинала $\tau$ существует наименьший кардинал, строго больший $\tau$. Он обозначается $\tau^+$ и называется \textit{кардиналом, следующим за $\tau$}.

\textbf{Свойство:} Если $\tau$ — кардинал, то $\tau^+$ — наименьший кардинал такой, что $\tau < \tau^+$. Кроме того, если $\tau$ --- бесконечный кардинал, то $\tau^+=Hart(\tau)$.
\end{definition}

\begin{theorem}
Множество всех порядковых типов (вполне упорядочений) мощности $\tau$ имеет мощность $\tau^+$, если $\tau$ --- бесконечный кардинал.

\textbf{Интуиция:} Для множества данной мощности существует много различных способов его вполне упорядочить, и количество таких способов (с точностью до изоморфизма) равно следующему кардиналу $\tau^+$.
\end{theorem}

\begin{definition}[Алефы]
Алефы определяются трансфинитной рекурсией по ординалам:
\begin{itemize}
    \item $\aleph_0 = \omega$ (наименьший бесконечный кардинал)
    \item $\aleph_{\al + 1} = (\aleph_\al)^+$ (кардинал, следующий за $\aleph_\al$)
    \item $\aleph_\la = \sup\{\aleph_\be \mid \be < \la\}$ для предельного ординала $\la$
\end{itemize}

\textbf{Интуиция:} Алефы — это все бесконечные кардиналы, упорядоченные по возрастанию. Каждый алеф $\aleph_\al$ является начальным ординалом $\omega_\al$.
\end{definition}



\subsection{Кардинальная арифметика}

\begin{definition}[Кардинальная арифметика]\;
Для кардиналов $\ka$ и $\la$:
\begin{enumerate}
    \item \textbf{Сумма:} $\ka + \la = |A \sqcup B|$, где $A$ и $B$ --- дизъюнктные множества мощностей $\ka$ и $\la$ соответственно
    \item \textbf{Произведение:} $\ka \cdot \la = |A \times B|$, где $A$ и $B$ --- множества мощностей $\ka$ и $\la$ соответственно
    \item \textbf{Степень:} $\ka^{\la} = |A^B|$, где $A^B$ --- множество всех функций из $B$ в $A$
\end{enumerate}
\end{definition}

\begin{theorem}[Свойства кардинальной арифметики]\;
\begin{enumerate}
    \item \textbf{Коммутативность:} $\ka + \la = \la + \ka$, $\ka \cdot \la = \la \cdot \ka$
    \item \textbf{Ассоциативность:} $(\ka + \la) + \mu = \ka + (\la + \mu)$, $(\ka \cdot \la) \cdot \mu = \ka \cdot (\la \cdot \mu)$
    \item \textbf{Дистрибутивность:} $\ka \cdot (\la + \mu) = \ka \cdot \la + \ka \cdot \mu$
    \item \textbf{Степень произведения:} $(\ka \cdot \la)^{\mu} = \ka^{\mu} \cdot \la^{\mu}$
    \item \textbf{Степень степени:} $(\ka^{\la})^{\mu} = \ka^{\la \cdot \mu}$
    \item \textbf{Сумма степеней:} $\ka^{\la + \mu} = \ka^{\la} \cdot \ka^{\mu}$
    \item Для бесконечных кардиналов: $\ka + \la = \ka \cdot \la = \max(\ka, \la)$ (если хотя бы один из них бесконечен)
    \item \textbf{Теорема Кантора:} $\ka < 2^{\ka}$ для любого кардинала $\ka$
    \item \textbf{Теорема Гессенберга:} Для любого бесконечного кардинала $\ka$ выполнено $\ka^2 = \ka$.
\end{enumerate}
\end{theorem}

\begin{remark}[Эквивалентность теоремы о квадрате и \AC]
Теорема о том, что для произвольного бесконечного множества $X$ выполнено $|X \times X| = |X|$, эквивалентна аксиоме выбора (\AC). 

В частности, без \AC могут существовать бесконечные множества $X$, для которых $|X \times X| > |X|$. Теорема Гессенберга гарантирует, что для кардиналов (которые определены через вполне упорядочения) квадрат всегда равен исходному кардиналу, но для произвольных множеств это требует \AC.
\end{remark}

\begin{theorem}[Связь степени и булеана]
Для любого множества $X$: $2^{|X|} = |\Pcal(X)|$, где $2 = \{0, 1\}$.

\textbf{Доказательство:} Биекция между $\Pcal(X)$ и $2^X$ задается характеристическими функциями: $A \mapsto \chi_A$, где $\chi_A(x) = 1$ если $x \in A$, и $0$ иначе.
\end{theorem}


\begin{definition}[Beth-числа]
    Beth-числа (обозначаются $\beth_\al$) определяются трансфинитной рекурсией по ординалам через булеан (степень множества):
    \begin{itemize}
        \item $\beth_0 = \aleph_0 = |\N|$ (счетная мощность)
        \item $\beth_{\al + 1} = 2^{\beth_\al} = |\Pcal(X)|$, где $|X| = \beth_\al$ (мощность булеана множества мощности $\beth_\al$)
        \item $\beth_\la = \sup\{\beth_\be \mid \be < \la\}$ для предельного ординала $\la$
    \end{itemize}
    
    \textbf{Интуиция:} Beth-числа описывают иерархию мощностей, получаемую последовательным применением операции взятия булеана. В частности, $\beth_1 = 2^{\aleph_0}$ — это мощность континуума.
    \end{definition}
    
    \begin{remark}
    В отличие от иерархии алефов, которая гарантированно является классом (все бесконечные кардиналы упорядочены и образуют собственный класс), иерархия beth не гарантированно является классом. Без аксиомы выбора (\AC) мы не можем доказать даже существование алефа континуума, то есть не можем гарантировать, что $\beth_1 = 2^{\aleph_0}$ является кардиналом (начальным ординалом). 
    
    Более того, без \AC может оказаться, что некоторые beth-числа не существуют как кардиналы. Так, в модели Соловея континуум невозможно вполне упрядочить, поэтому $\beth_1$ не существует как алеф.
    \end{remark}
    
    
    
    

\begin{theorem}[Гипотеза Континуума (\CH)]
$2^{\aleph_0} = \aleph_1$

\textbf{Обобщенная гипотеза континуума (\GCH):} $2^{\aleph_{\al}} = \aleph_{\al+1}$ для всех ординалов $\al$.

\textbf{Статус:} \CH независима от \ZFC (Гёдель, 1940; Коэн, 1963). То есть \ZFC не может ни доказать, ни опровергнуть \CH.
\end{theorem}

\begin{example}[Примеры кардинальной арифметики]\;
\begin{enumerate}
    \item $\aleph_0 + \aleph_0 = \aleph_0$ (объединение двух счетных множеств счетно)
    \item $\aleph_0 \cdot \aleph_0 = \aleph_0$ (декартово произведение двух счетных множеств счетно)
    \item $2^{\aleph_0} = |\R|$ (мощность континуума)
    \item $\aleph_0^{\aleph_0} = 2^{\aleph_0}$ (множество всех последовательностей натуральных чисел имеет мощность континуума)
    \item Если \CH верна, то $2^{\aleph_0} = \aleph_1$
\end{enumerate}
\end{example}

\begin{exercise}[Арифметика кардиналов]
\begin{enumerate}
    \item \textbf{Доказать свойства суммы кардиналов:}
    \begin{itemize}
        \item Коммутативность: $\ka + \la = \la + \ka$
        \item Ассоциативность: $(\ka + \la) + \mu = \ka + (\la + \mu)$
        \item Нейтральный элемент: $\ka + 0 = 0 + \ka = \ka$
        \item Для бесконечных кардиналов: если $\ka$ или $\la$ бесконечен, то $\ka + \la = \max(\ka, \la)$
    \end{itemize}
    
    \item \textbf{Доказать свойства произведения кардиналов:}
    \begin{itemize}
        \item Коммутативность: $\ka \cdot \la = \la \cdot \ka$
        \item Ассоциативность: $(\ka \cdot \la) \cdot \mu = \ka \cdot (\la \cdot \mu)$
        \item Дистрибутивность: $\ka \cdot (\la + \mu) = \ka \cdot \la + \ka \cdot \mu$ и $(\ka + \la) \cdot \mu = \ka \cdot \mu + \la \cdot \mu$
        \item Нейтральный элемент: $1 \cdot \ka = \ka \cdot 1 = \ka$
        \item Нулевой элемент: $0 \cdot \ka = \ka \cdot 0 = 0$
        \item Для бесконечных кардиналов: если $\ka$ или $\la$ бесконечен, то $\ka \cdot \la = \max(\ka, \la)$
    \end{itemize}
    
    \item \textbf{Доказать свойства степени кардиналов:}
    \begin{itemize}
        \item $\ka^{\la + \mu} = \ka^{\la} \cdot \ka^{\mu}$
        \item $(\ka \cdot \la)^{\mu} = \ka^{\mu} \cdot \la^{\mu}$
        \item $(\ka^{\la})^{\mu} = \ka^{\la \cdot \mu}$
        \item $\ka^0 = 1$, $\ka^1 = \ka$, $1^{\ka} = 1$, $0^{\ka} = 0$ (для $\ka > 0$)
    \end{itemize}
    
    \item \textbf{Связь операций с порядком:}
    \begin{itemize}
        \item Доказать, что если $\ka \le \la$ и $\mu \le \nu$, то $\ka + \mu \le \la + \nu$ и $\ka \cdot \mu \le \la \cdot \nu$
        \item Доказать, что если $\ka < \la$ и $\mu > 0$, то $\ka^{\mu} \le \la^{\mu}$ и $\mu^{\ka} \le \mu^{\la}$
        \item Показать, что строгое неравенство может не сохраняться: $\aleph_0 < 2^{\aleph_0}$, но $\aleph_0^{\aleph_0} = 2^{\aleph_0}$
    \end{itemize}
    
    \item \textbf{Вычислить конкретные примеры:}
    \begin{itemize}
        \item $\aleph_0 + \aleph_0$, $\aleph_0 \cdot \aleph_0$, $\aleph_0^{\aleph_0}$
        \item $\aleph_1 + \aleph_0$, $\aleph_1 \cdot \aleph_0$, $\aleph_0^{\aleph_1}$
        \item $2^{\aleph_0} + 2^{\aleph_0}$, $2^{\aleph_0} \cdot 2^{\aleph_0}$, $(2^{\aleph_0})^{\aleph_0}$
    \end{itemize}
    
    \item \textbf{Сравнение алефов и beth:}
    \begin{itemize}
        \item Показать, что $\beth_\al \ge \aleph_\al$ для всех ординалов $\al$
        \item Показать, что при \GCH выполнено $\beth_\al = \aleph_\al$ для всех $\al$
        \item Объяснить, почему без \AC может быть $\beth_1 \neq \aleph_1$ (и даже $\beth_1$ может не существовать как алеф)
    \end{itemize}
\end{enumerate}

\textit{Подсказка:} Для доказательства свойств используйте определения операций через мощности множеств и биекции. Для бесконечных кардиналов используйте теорему Гессенберга и свойства максимума. Для теоремы Кантора используйте диагональный аргумент.
\end{exercise}

\subsubsection{Типы кардиналов и конфинальность}

\begin{definition}[Конфинальность кардинала]
\textit{Конфинальностью} кардинала $\ka$ (обозначается $\cf(\ka)$) называется наименьший кардинал $\la$ такой, что существует функция $f: \la \to \ka$ с $\sup(\ran(f)) = \ka$.

\textbf{Интуиция:} Конфинальность показывает, насколько «мелко» можно покрыть кардинал меньшими кардиналами. Кардинал $\ka$ называется \textit{регулярным}, если $\cf(\ka) = \ka$, и \textit{сингулярным}, если $\cf(\ka) < \ka$.
\end{definition}

\begin{definition}[Изолированный и предельный кардинал]
\begin{itemize}
    \item Кардинал $\ka$ называется \textit{изолированным} (или \textit{преемником}), если существует кардинал $\tau$ такой, что $\ka = \tau^+$.
    \item Кардинал $\ka$ называется \textit{предельным}, если он не является изолированным, то есть не существует кардинала $\tau$ такого, что $\ka = \tau^+$.
\end{itemize}

\textbf{Примеры:} $\aleph_1, \aleph_2, \aleph_{\omega+1}$ — изолированные кардиналы. $\aleph_0, \aleph_\omega, \aleph_{\omega_\omega}$ — предельные кардиналы.
\end{definition}

\begin{definition}[Слабо недостижимый кардинал]
Кардинал $\ka$ называется \textit{слабо недостижимым}, если:
\begin{enumerate}
    \item $\ka$ — регулярный (т.е. $\cf(\ka) = \ka$)
    \item $\ka$ — предельный (т.е. не является изолированным)
    \item $\ka > \aleph_0$ (несчетный)
\end{enumerate}

\textbf{Интуиция:} Слабо недостижимый кардинал нельзя «достичь» ни как предел меньших кардиналов (из-за регулярности), ни как преемник другого кардинала (из-за предельности).
\end{definition}

\begin{definition}[Сильно недостижимый кардинал]
Кардинал $\ka$ называется \textit{сильно недостижимым}, если:
\begin{enumerate}
    \item $\ka$ — регулярный
    \item $\ka$ — предельный
    \item $\ka > \aleph_0$
    \item Для любого кардинала $\la < \ka$ выполнено $2^{\la} < \ka$ (замкнутость относительно операции булеана)
\end{enumerate}

\textbf{Связь:} Всякий сильно недостижимый кардинал является слабо недостижимым. Обратное верно при предположении \GCH.
\end{definition}

\begin{theorem}[Модель \ZF в универсуме фон Неймана]
Если $\ka$ — сильно недостижимый кардинал, то универсум фон Неймана $V_\ka$ является моделью теории \ZF.

\textbf{Следствие:} Если существует сильно недостижимый кардинал (обозначим это утверждение как \SI), то $\Con(\ZF + \SI) \to \Con(\ZF)$. То есть, если теория \ZF с аксиомой о существовании сильно недостижимого кардинала непротиворечива, то и теория \ZF непротиворечива.

\textbf{Интуиция:} Сильно недостижимый кардинал $\ka$ настолько «большой», что универсум $V_\ka$ содержит все множества, которые можно построить в \ZF, и при этом сам $V_\ka$ ведет себя как модель \ZF.
\end{theorem}

\begin{theorem}[О невозможности доказать непротиворечивость \SI в \ZF]
Невозможно доказать в \ZF, что аксиому о существовании сильно недостижимого кардинала (\SI) нельзя опровергнуть. То есть, \ZF не доказывает $\Con(\ZF) \to \Con(\ZF + \SI)$.

\textbf{Доказательство:} Предположим противное: \ZF доказывает $\Con(\ZF) \to \Con(\ZF + \SI)$.

Рассмотрим теорию $T = \ZF + \Con(\ZF)$. Покажем, что тогда в $T$ можно вывести $\Con(T)$.
\begin{align*}
T &\vdash \ZF+\Con(\ZF+\SI) \\
& \vdash\ZF + \Con(\ZF+\Con(\ZF)) \\
& \vdash\Con(T)
\end{align*}

Это противоречит теореме Гёделя о неполноте в отношении теории $T$.
\end{theorem}


\section{Как конструируются математические структуры}

\subsection{Базовые конструкции}

Перечислим справочно основные элементарные теоретико-множественные конструкции, необходимые для построения математических структур:

\begin{definition}[Базовые конструкции]
\begin{enumerate}
\item \textbf{Пара и упорядоченная пара:} $\{x,y\}$; $\langle x, y \rangle \eqdef \{\{x\}, \{x, y\}\}$. 

\textit{Главное свойство:} $\langle a, b \rangle = \langle c, d \rangle \iff a=c \land b=d$.
        
\item \textbf{Декартово произведение:} $A \times B \eqdef \{\langle x, y \rangle \mid x \in A \land y \in B\}$.
        
\item \textbf{Бинарные отношения:} $R \subseteq A \times B$.
\begin{itemize}
    \item \textit{Рефлексивное:} $\forall x \in A: xRx$; иррефлексивное: $\forall x \in A: \neg(xRx)$.
    \item \textit{Симметричное:} $\forall x,y \in A: xRy \to yRx$.
    \item \textit{Транзитивное:} $\forall x,y,z \in A: xRy \land yRz \to xRz$.
    \item \textit{Антисимметричное:} $\forall x,y \in A: xRy \land yRx \to x=y$.
    \item \textit{Связное:} $\forall x,y \in A: xRy \lor yRx \lor x=y$.
    \item \textit{Всюду определенное:} $\forall x\in A:\exists y\in B:xRy$.
    \item \textit{Всюду значное:} $\forall y\in B:\exists x\in A:xRy$.
    \item \textit{Однозначное (функциональное):} $\forall x\in A:\forall y,z\in B:xRy\land xRz\to y=z$.
    \item \textit{Обратно однозначное:} $\forall x,z\in A:\forall y\in B:xRy\land zRy\to x=z$.    
\end{itemize}

В том числе порядки:
\begin{itemize}
    \item \textit{Предпорядок:} рефлексивное и транзитивное.
    \item \textit{Частичный порядок:} антисимметричный предпорядок.
    \item \textit{Линейный порядок:} связный частичный порядок.
    \item \textit{Вполне упорядочение:} линейный порядок, для которого $\forall X \subseteq A: X \neq \emptyset \implies \exists x \in X: \forall y \in X: xRy$.
\end{itemize}

В том числе функции:
\begin{itemize}
    \item \textit{Функция $f$ из $A$ в $B$:} всюду определенное однозначное отношение. Запись: $f:A\to B$, вместо $xfy$ пишем $f(x)=y$.
    \item \textit{Сюръекция:} $\forall y\in B:\exists x\in A:\;f(x)=y$.
    \item \textit{Инъекция:} $\forall x,z\in A:f(x)=f(z)\to x=z$.
    \item \textit{Биекция:} инъекция и сюръекция одновременно.
    \item \textit{Обратная функция:} $f^{-1}:B\to A$ --- обратная функция к $f:A\to B$, если $f(f^{-1}(y))=y$ для всех $y\in B$ и $f^{-1}(f(x))=x$ для всех $x\in A$.
    \item \textit{Композиция функций:} $g \circ f:A\to C$ --- композиция функций $f:A\to B$ и $g:B\to C$: $(g \circ f)(x) = z \iff \exists y\in B: f(x)=y \land g(y)=z$.
    \item \textit{Образ множества:} $f[X] \eqdef \{f(x) \mid x \in X\}$.
    \item \textit{Прообраз множества:} $f^{-1}[Y] \eqdef \{x \mid f(x) \in Y\}$.
    \item \textit{Область значений:} $ran(f) \eqdef \{y \mid \exists x \in A: f(x) = y\}$.
    \item \textit{Область определения:} $dom(f) \eqdef \{x \mid \exists y \in B: f(x) = y\}$.
    \item \textit{Последовательность:} $f:\om\to A$ --- функция из натуральных чисел в $A$, часто обозначается как $\langle a_n \rangle_{n\in\om}$.
\end{itemize}
    
\textbf{Отношение эквивалентности:} бинарное отношение $\sim$ на множестве $X$ (подмножество $X \times X$), которое рефлексивно, симметрично и транзитивно. Класс эквивалентности: $[x]_\sim \eqdef \{y \in X \mid x \sim y\}$. Основное свойство: разбиение множества на непересекающиеся классы эквивалентности.
        
\textbf{Фактор-множество:} $X/{\sim} \eqdef \{[x]_\sim \mid x \in X\}$ --- множество классов эквивалентности.

\item \textbf{Кортеж:} $\langle x_1, x_2, \dots, x_n \rangle$ --- функция, определенная на множестве $\{1,2,\dots,n\}$. При $n=0$ это пустой кортеж $\langle\rangle$ (пустое множество).

\item \textit{$n$-арное декартово произведение:} $A_1 \times A_2 \times \dots \times A_n \eqdef \{\langle a_1, a_2, \dots, a_n \rangle \mid a_i \in A_i\}$. Если все $A_i = A$, то это $A^n$ (при $n=2$ это $A \times A$, при $n=1$ это $A$, при $n=0$ это $A^0 = \{\emptyset\}$).

\item \textit{$n$-арное отношение:} $R \subseteq A_1 \times A_2 \times \dots \times A_n$ --- подмножество $n$-арного декартового произведения.

\item \textit{$n$-арная функция:} $f:A_1 \times A_2 \times \dots \times A_n \to B$ --- функция из $n$-арного декартового произведения в $B$.

\item \textbf{Звезда Клини:} $X^* = \bigcup_{n=0}^{\infty} X^n = X^0 \cup X \cup (X \times X) \cup (X \times X \times X) \cup \dots$ --- множество всех конечных списков (слов, кортежей), если мы не знаем заранее их длину.

\begin{example}
\begin{itemize}
    \item \textbf{Текст:} Множество всех книг — это $\Sigma^*$, где $\Sigma$ — алфавит. Это объединение множества книг длиной 1 символ, 2 символа и т.д.
    \item \textbf{Лес:} Если мы определили множество деревьев $\mathcal{T}$, то упорядоченный лес (набор деревьев) — это элемент $\mathcal{T}^*$.
    \item \textbf{Пространство многочленов:} $K[x]$ задается на носителе, который можно определить как $K^*$, т.е. множество всех конечных последовательностей из элементов кольца $K$ (которые являются коэффициентами многочлена).
\end{itemize}
\end{example}

Звезду Клини можно расширить до прямого произведения $X^\om$, которое представляет собой множество всех функций из $\om$ в $X$, т.е. всех бесконечных последовательностей элементов $X$. Это позволяет конструировать, например, пространство путей и определять на нем цилиндрическую топологию.
\end{enumerate}
\end{definition}

\subsection{Алгебраическая структура}
Само по себе множество (носитель) аморфно. Чтобы превратить его в пространство (математическую модель), мы должны навесить на него определенную структуру.

Базовые конструкции уже позволяют нам дать следующее определение:
\begin{definition}
Алгебраическая структура $\Mcal$ на множестве $M$ --- это пара $\Mcal = \langle M, \Sigma \rangle$, где $M$ — носитель (домен) структуры, а $\Sigma$ — сигнатура, включающая:
\begin{itemize}
    \item \textit{Операции} (функции $f: M^n \to M$) — законы преобразования элементов множества $M$.
    \item \textit{Отношения} (предикаты $P \subseteq M^n$) — законы соотношения элементов множества $M$. Отметим, что отношение равенства ($=$) является фундаментальным (заданным в логике предикатов) отношением и обычно не включается в сигнатуру алгебраической структуры в явном виде.
    \item \textit{Константы} (выделенные элементы $c \in M$).
\end{itemize}
\end{definition}

Однако даже в математике (не говоря уже о физике и других науках, где важны измерения, непрерывность и вероятность) нам часто требуется более сложная структура, чем просто множество с операциями и отношениями. Точнее, нам требуется добавить к базовой структуре вспомогательные (числовые) структуры и системы множеств, чтобы наделить ее более сложной эмерджентной семантикой.

В связи с этим дадим следующее определение:
\begin{definition}[Аналитическая структура]
Для задач анализа, геометрии и теории вероятностей алгебраическая структура расширяется до кортежа:
$$ \Mcal = \langle M, \Sigma, \mathcal{T}, \Phi \rangle, $$
который называется \textit{аналитической структурой}. 

\textbf{Алгебраическая сигнатура $\Sigma$} включает:
\begin{itemize}
    \item \textbf{Внутренние операции:} Отображения $f: M^n \to M$ (операции на домене).
    \item \textbf{Внешние операции:} Отображения $\cdot: R \times M \to M$, где $R$ — внешнее кольцо или поле (действие кольца на домене, превращающее абелеву группу в модуль над $R$).
    \item \textbf{Отношения:} Предикаты $P \subseteq M^n$.
    \item \textbf{Константы:} Элементы $c \in M$.
\end{itemize}

К сигнатуре $\Sigma$ добавляются два внешних слоя:
\begin{itemize}
    \item $\mathcal{T}$ (\textbf{Слой пространства}): Системы подмножеств $\mathcal{T} \subseteq \Pcal(M)$ для задания топологии, измеримости (алгебра множеств) и факторизации (фильтр).
    \item $\Phi$ (\textbf{Слой измерения}): Функционалы $F: M \to \mathbb{K}$ ($\R$ или $\C$), отображающие элементы домена в числовое поле $\mathbb{K}$.
\end{itemize}
\end{definition}


\marginfigure[0.1][0.75]{
\begin{tikzpicture}[
    % Стиль для блоков
    layer/.style={
        rectangle, 
        draw=darkblue, 
        very thick, 
        minimum width=10cm, 
        minimum height=1cm, 
        align=center, 
        rounded corners=3pt,
        font=\bfseries
    },
    % Стиль для стрелок
    arrow/.style={
        ->, 
        >={Stealth[length=3mm]}, 
        ultra thick, % Сделали стрелки потолще
        color=darkgray
    }
]

    % --- 4. Метрический слой (Вершина) ---
    \node[layer, fill=red!10] (metric) {
        4. Метрический слой ($\Phi$) \\ 
        \normalfont\footnotesize Функционалы: Метрика $d$, Норма $\|\cdot\|$, Мера $\mu$ и т.д.
    };

    % --- 3. Топологический слой ---
    % Увеличили отступ (below) до 1.5cm для видимости стрелок
    \node[layer, fill=orange!10, below=1cm of metric] (topo) {
        3. Топологический слой ($\mathcal{T}$) \\ 
        \normalfont\footnotesize Пространство: Топология $\tau$, Фильтры $\mathcal{F}$, $\sigma$-алгебры
    };

    % --- 2. Сигнатура (Алгебра + Отношения) ---
    \node[layer, fill=yellow!10, below=1cm of topo] (signature) {
        2. Алгебраическая сигнатура ($\Sigma$) \\ 
        \normalfont\footnotesize Внутр. опер. $f: M^n \to M$, Внеш. опер. $R \times M \to M$ \\ 
        \normalfont\footnotesize Отношения $P$, Константы $c$
    };

    % --- 1. Носитель (Фундамент) ---
    \node[layer, fill=blue!10, below=1cm of signature] (carrier) {
        1. Домен ($M$) \\ 
        \normalfont\footnotesize Исходное множество элементов
    };

    % --- Стрелки потока конструирования ---
    \draw[arrow] (carrier) -- (signature);
    \draw[arrow] (signature) -- (topo);
    \draw[arrow] (topo) -- (metric);

\end{tikzpicture}
}


\begin{example}[Сравнение подходов: Евклидово пространство $\R^3$]
\begin{itemize}
    \item Как \textit{алгебраическая структура}: Это векторное пространство $\langle \R^3, +, \cdot \rangle$. Есть только векторы и операции над ними.
    \item Как \textit{аналитическая структура}: Это Банахово пространство $\langle \R^3, +, \cdot, \|\cdot\| \rangle$. Мы добавляем норму (длину), что позволяет измерять расстояния и определять сходимость. Это позволяет нам определить понятия непрерывности, дифференцируемости и интегрируемости функций на $\R^3$.
\end{itemize}
\end{example}

Рассмотрим подробнее архитектуру обвеса аналитической структуры.

\subsubsection{Алгебраический слой (Операции)}
Задает законы \textit{взаимодействия} или \textit{преобразования} элементов. Элементы порождают новые элементы того же мира.

\textbf{Внутренние операции:} Отображения $f: M^n \to M$ (операции на самом домене).
\begin{itemize}
    \item \textbf{Смысл:} Динамика, синтез, движение, симметрии и т.д.
    \item \textbf{Примеры:}
    \begin{itemize}
        \item \textit{Сложение} ($+: \R \times \R \to \R$) — комбинирование величин.
        \item \textit{Обращение} ($^{-1}: G \to G$) — отмена действия (обратимость в группе).
        \item \textit{Композиция} ($\circ: \Fcal \times \Fcal \to \Fcal$) — операция на пространстве унарных операций $\Fcal$.
    \end{itemize}
\end{itemize}

\textbf{Внешние операции (Действие кольца/поля):} Отображения $\cdot: R \times M \to M$, где $R$ — внешнее кольцо или поле (не часть домена $M$).
\begin{itemize}
    \item \textbf{Смысл:} Шкалирование, масштабирование, умножение на скаляры. Превращает абелеву группу $(M, +)$ в \textit{модуль над кольцом $R$} (или векторное пространство, если $R$ — поле).
    \item \textbf{Примеры:}
    \begin{itemize}
        \item \textit{Умножение вектора на скаляр} ($\cdot: \R \times V \to V$) — превращает абелеву группу векторов в векторное пространство над $\R$.
        \item \textit{Действие группы на множестве} ($\cdot: G \times X \to X$) — группа $G$ действует на множестве $X$ (представление группы).
        \item \textit{Модуль над кольцом} ($\cdot: R \times M \to M$) — обобщение векторного пространства, где $R$ — произвольное кольцо (не обязательно поле).
    \end{itemize}
    \item \textbf{Философский смысл:} Внешние операции связывают структуру с «измерительной шкалой» (кольцом/полем), позволяя количественно оперировать элементами домена. Это мост между абстрактной структурой и числовыми системами.
\end{itemize}

\subsubsection{Реляционный слой (Предикаты и Законы)}
Если операции позволяют \textit{строить} новые объекты, то предикаты позволяют \textit{утверждать} истину о них. Без этого слоя невозможно сформулировать ни одного физического закона или аксиом стандартных математических структур.

\begin{itemize}
    \item \textbf{Тип объекта:} Отношение $P \subseteq M^n$. В логике это интерпретируется как функция $P: M^n \to \{True, False\}$.
    \item \textbf{Фундаментальные типы:}
    \begin{enumerate}
        \item \textbf{Равенство ($=$):} Базовый предикат любой модели.
        Определяет тождество объектов. Без него невозможно написать уравнение $F=ma$ (оно утверждает, что вектор силы и вектор $ma$ — это один и тот же объект).
        
        \item \textbf{Порядок ($\le$):} Предикат сравнения.
        Определяет направление процессов (энтропия $dS \ge 0$) и оптимизацию (энергия $E \to \min$).
        
        \item \textbf{Физические законы (Уравнения как предикаты):}
        Любое уравнение $f(x) = g(x)$ — это сложный предикат $P(x)$, который «вырезает» из всех возможных состояний только \textit{физически допустимые}.
        \textit{Пример:} Уравнение Шрёдингера $\hat{H}\psi = E\psi$ — это предикат на гильбертовом пространстве функций. Реальность — это подмножество функций, для которых этот предикат истинен.
    \end{enumerate}
\end{itemize}

\subsubsection{Топологический слой (Системы подмножеств)}
Задает \textit{текстуру} пространства, определяя понятия «близость», «локальность» и «шум» без использования линейки.
\begin{itemize}
    \item \textbf{Тип объекта:} Семейство $\mathcal{U} \subseteq \Pcal(M)$.
    \item \textbf{Смысл:} совместное (эмерджентное) свойство совокупностей индивидов из домена $M$.
    \item \textbf{Примеры:}
    \begin{itemize}
        \item \textit{Топология} ($\tau$): Набор открытых множеств. Определяет близость и непрерывность.
        \item \textit{Алгебра множеств}: Набор наблюдаемых событий. Определяет область применимости вероятности и измеримости.
        \item \textit{Фильтр}: Математическая модель понятия «большинство» или «окрестность». Фильтры (а точнее, ультрафильтры) используются, например, для факторизации домена.
    \end{itemize}
\end{itemize}

\subsubsection{Метрический слой (Функционалы)}
Проецирует абстрактную структуру на шкалу реальности (числа), позволяя производить  \textit{измерения, шкалирование и количественную оценку}.
\begin{itemize}
    \item \textbf{Тип объекта:} Отображение $F: X \to \text{Числовое поле } (\R, \C)$, где $X=M^n$ для некоторого $n$, или $X=\Tcal$ (топологический слой).
    \item \textbf{Смысл:} Оценка, стоимость, энергия, вероятность, взаимное расположение (угол) и т.д.
    \item \textbf{Примеры:} 
    \begin{itemize}
        \item \textit{Метрика} ($d: M \times M \to \R$): Расстояние между точками.
        \item \textit{Норма} ($\|x\|: M \to \R$): Размер (магнитуда) вектора.
        \item \textit{Мера} ($\mu: \Tcal \to \R$): Объем множества или вероятность события.
        \item \textit{Функция полезности} ($U: M \to \R$): Ценность исхода в экономике.
    \end{itemize}
\end{itemize}

\begin{remark}[Синтез: Многослойность структур]
Реальные математические модели обычно сочетают все четыре слоя.
Например, \textbf{Гильбертово пространство} $L^2$ (основа квантовой механики) устроено так:
\begin{enumerate}
    \item \textbf{Носитель:} Классы эквивалентности функций (факторизация).
    \item \textbf{Алгебра:} Сложение функций (векторное пространство).
    \item \textbf{Функционал:} Скалярное произведение $\langle f, g \rangle$.
    \item \textbf{Топология:} Порождается функционалом нормы $\|f\| = \sqrt{\langle f, f \rangle}$.
\end{enumerate}
\end{remark}

\begin{example}[Путь как связующий объект]
Рассмотрим \textit{путь} — непрерывное движение точки $\gamma: [0,1] \to M$.
\begin{itemize}
    \item Его существование обеспечивает \textbf{Топологический слой} $\mathcal{T}$ (требование непрерывности).
    \item Его геометрическая форма (траектория) возникает в \textbf{Реляционном слое} $\mathcal{R}$ (мы отождествляем движения с разной скоростью: $\gamma \sim \gamma \circ \varphi$).
    \item Его длину измеряет \textbf{Метрический слой} $\Phi$.
    \item А классы деформируемых путей (петли) образуют \textbf{Алгебраический слой} $\Sigma$ (фундаментальную группу $\pi_1(X)$), превращая пространство в группу.
\end{itemize}
Это показывает, что слои модели — это не изолированные полки, а взаимосвязанные механизмы. Более того, здесь мы даже наблюдаем рождение еще одной надструктуры --- фактор-множества путей.
\end{example}

\begin{remark}[Связка: Функционалы + Предикаты = Вариационные принципы]
В развитых моделях (теоретическая механика, теория поля) структура работает в связке:
\begin{enumerate}
    \item Слой \textbf{Функционалов} ($\Phi$) задает величины (например, Лагранжиан $\mathcal{L}$ или Действие $S$).
    \item Слой \textbf{Предикатов} ($\mathcal{R}$) формулирует закон природы в виде вариационного уравнения:
    $$ \delta S[\gamma] = 0 $$
\end{enumerate}
Таким образом, физический закон — это предикат, наложенный на функционал. Он утверждает, что природа выбирает траектории, экстремальные с точки зрения функционала.
\end{remark}


Дальше мы сформулируем и покажем на примерах ряд традиционных способов генерации структур, позволяющих строить и расширять домены структур, а также снабжать их архитектурно более сложным обвесом, включающим все перечисленные слои аналитической структуры.

\subsection{Метод Выделения}

    Метод выделения (спецификации) позволяет «высекать» нужную структуру из ранее определенного множества.
    Это реализация философского принципа \textit{Omnis determinatio est negatio} («Всякое определение есть отрицание», Спиноза): чтобы создать что-то определенное, нужно запретить всё остальное.
    
    \begin{definition}[Конструкция через ограничение]
    Пусть $U$ — «универсальное» множество (обычно булеан $\Pcal(X)$ или декартово произведение), содержащее все мыслимые сочетания элементов.
    Мы строим новое множество $S$ как подмножество $U$, элементы которого удовлетворяют предикату (свойству) $\phi$:
    $$ S = \{x \in U \mid \phi(x)\} $$
    \end{definition}
    
    \begin{example}[Графы с $k$ компонентами]
        Пусть мы хотим построить множество всех графов на $n$ вершинах, которые состоят ровно из $k$ несвязных «островов» (компонент связности).
        Мы действуем методом отсечения:
        
        \begin{enumerate}
            \item \textbf{Сырой материал (Универсум):} 
            Фиксируем множество вершин $V = \{1, \dots, n\}$.
            Универсальное множество $U$ — это множество \textbf{всех возможных графов} (т.е. бинарных отношений) на $V$.
            Поскольку граф определяется набором ребер, а всего возможных ребер $\binom{n}{2}$, то мощность универсума $|U| = 2^{\binom{n}{2}}$. Это «хаос» всех возможных связей.
        
            \item \textbf{Резец (Предикат):} 
            Нам нужно условие $\phi(G)$, выделяющее нужную топологию.
            Определим на вершинах отношение достижимости: $u \sim_G v$, если в графе $G$ существует путь между $u$ и $v$. Это отношение эквивалентности.
            Классы эквивалентности $V / \sim_G$ — это и есть компоненты связности.
            Наше условие:
            $$ \phi(G) \iff |V / \sim_G| = k $$
            (число классов эквивалентности равно $k$).
        
            \item \textbf{Результат:} 
            Искомое множество $S_k = \{G \in U \mid \text{число компонент } G \text{ равно } k\}$.
            \textit{Пример:} Если $k=1$, мы вырезаем множество всех \textbf{связных} графов. Если $k=n$, мы оставляем только один граф — пустой (без ребер).
        \end{enumerate}
        \end{example}
    
    \begin{example}[Топология как вырезание]
    Что такое топологическое пространство на множестве $X$?
    Это не произвольный набор подмножеств.
    \begin{enumerate}
        \item Берем булеан $\Pcal(X)$ — множество всех подмножеств.
        \item Вырезаем из него то семейство $\tau$, которое удовлетворяет аксиомам топологии (замкнутость относительно объединений и конечных пересечений, наличие $\emptyset$ и $X$).
    \end{enumerate}
    \end{example}
    
    
    \begin{remark}[Связь с рангом]
    Метод выделения **понижает энтропию**, не повышая онтологический ранг.
    $$ S \subseteq U \implies rank(S) \le rank(U) $$
    Мы не создаем новые сущности, мы лишь наводим порядок в существующих: оставляем существенное и отбрасываем избыточное.
    \end{remark}
    


\subsection{Факторизация}

Самый мощный инструмент абстракции --- это факторизация. Мы берем множество (домен уже построенной структуры) и объявляем разные элементы «одним и тем же» используя заданное на <<старом>> множестве отношение эквивалентности (или разбиение).

\begin{definition}[Фактор-множество]
Пусть $X$ --- множество, а $\sim$ --- отношение эквивалентности. Фактор-множество $X/{\sim}$ --- это множество классов эквивалентности $[x] = \{y \in X \mid x \sim y\}$.
\end{definition}


\begin{remark}[Условие конгруэнтности]
    Факторизация имеет смысл только тогда, когда отношение эквивалентности $\sim$ \textbf{уважает} сигнатуру (операции и отношения) исходной структуры. Это свойство называется \textit{конгруэнтностью}.
    
    Пусть $f$ — $n$-арная операция на $X$. Чтобы индуцировать соответствующую операцию на фактор-множестве $X/{\sim}$, необходимо, чтобы результат не зависел от выбора представителей классов:
    $$ x_1 \sim y_1, \dots, x_n \sim y_n \to f(x_1, \dots, x_n) \sim f(y_1, \dots, y_n) $$
    В противном случае операция на классах будет \textit{некорректно определена} (ill-defined), и факторизация не будет иметь смысла.

    Аналогичное требование предъявляется к отношениям:
    $$ x_1 \sim y_1, \dots, x_n \sim y_n \to P(x_1, \dots, x_n) \iff P(y_1, \dots, y_n).$$
    
    \textbf{Философский смысл:} Мы не можем произвольно объявлять разные объекты «одним и тем же». Абстракция легитимна только тогда, когда «законы» (операции и отношения исходной системы) не замечают различия между склеиваемыми объектами. Если закон $f$ чувствителен к различию между $x$ и $y$, то $x$ и $y$ онтологически различны в рамках этой системы законов. В этом случае мы либо не можем их склеивать, либо мы должны <<забыть>> про этот закон (исключить из факторизованной структуры).
    \end{remark}


\begin{remark}[Онтологическая цена абстракции: повышение ранга домена]
    Важно отметить, что факторизация, как правило, повышает \textit{ранг} множества в иерархии фон Неймана.
    
    Пусть $rank(X) = \alpha$. Элементы фактор-множества $X/{\sim}$ — это классы эквивалентности $[x]$, которые являются подмножествами $X$.
    Согласно свойствам ранга, $rank([x]) \le rank(X) = \alpha$.
    
    Тогда ранг самого фактор-множества (множества классов):
    $$ rank(X/{\sim}) = \sup \{rank([x]) + 1\} $$
    
    Здесь возможны два случая:
    \begin{enumerate}
        \item \textbf{Стандартный случай ($\alpha + 1$):} Если хотя бы один класс эквивалентности имеет тот же ранг $\alpha$, что и исходное множество (что верно для построения $\Z, \Q, \R$), то ранг возрастает: $rank(X/{\sim}) = \alpha + 1$.
        \item \textbf{Предельный случай ($\alpha$):} В редких случаях (для бесконечных предельных ординалов и «мелких» разбиений) ранг может остаться прежним.
    \end{enumerate}
    
    \textbf{Интерпретация:} Переход к абстракции (от конкретных объектов к их типам) обычно поднимает нас на один уровень выше в иерархии бытия. Мы платим усложнением структуры (ростом ранга) за семантическое упрощение.
    \end{remark}

\subsubsection{Где работает Факторизация: примеры}

Метод факторизации (отождествления) используется повсеместно для создания новых структур из имеющихся. Мы берем структуру и объявляем некоторые различия «несущественными».

\begin{enumerate}
    \item \textbf{Числовые системы:} Как показано ниже, $\Z$, $\Q$, $\R$ и $\C$ строятся как классы эквивалентности пар или последовательностей.
    
    \item \textbf{Модулярная арифметика («Часы»):} 
    Кольцо вычетов $\Z_n = \Z / n\Z$. Мы отождествляем числа, дающие одинаковый остаток от деления на $n$. 
    \textit{Пример:} На часах $13:00$ и $1:00$ — это одно и то же время ($13 \equiv 1 \pmod{12}$).
    
    \item \textbf{Геометрия: равенство фигур.}
    Мы отождествляем фигуры, если они конгруэнтны.
    \textit{Смысл:} Физические объекты (как предметы в реальном мире) не меняются при перемещении, вращении или отражении. Мы можем отождествлять фигуры, если они идентичны в этих преобразованиях.
    
    \item \textbf{Проективная геометрия.}
    Это пространство «лучей зрения». Мы отождествляем все точки на одной прямой, проходящей через начало координат: $\vec X \sim \lambda \vec X$ для $\lambda \neq 0$.
    \textit{Смысл:} Для глаза все точки на одном луче проецируются в один пиксель. Реальность наблюдателя — это фактор-пространство $\R^3$ по отношению коллинеарности.

    \item \textbf{Дифференциальная геометрия: Параметризация путей.}
    Путь в пространстве часто определяют как класс эквивалентности всех отображений $\gamma: [0,1] \to \R^n$, переходящих друг в друга при замене параметра (времени) $t \mapsto \varphi(t)$.
    \textit{Смысл:} Физическая траектория (мировая линия) объективна, а то, с какой скоростью мы движемся по ней в формулах — это лишь наш способ описания (выбор калибровки времени).

    \item \textbf{Функциональный анализ: Пространства $L^p$.}
    Строго говоря, элементами пространства $L^p$ (например, квадратично интегрируемых функций, важных для квантовой механики) являются не функции, а классы эквивалентности функций.
    Мы отождествляем $f \sim g$, если они отличаются только на множестве меры нуль ($\int |f-g| = 0$).
    \textit{Важность:} Без этой факторизации норма $\|f\|$ не была бы нормой (существовали бы ненулевые векторы с нулевой длиной), и метрическое пространство не работало бы. Мы «вырезаем» невидимые различия.
    
    \item \textbf{Топология: Склейка (Тор и Лента Мёбиуса).}
    Тор (поверхность бублика) можно построить из квадрата $[0,1] \times [0,1]$, если отождествить (склеить) противоположные стороны: $(x, 0) \sim (x, 1)$ и $(0, y) \sim (1, y)$.
    \textit{Физика:} Это модель «вселенной Pac-Man», где, улетая за правый край экрана, появляешься слева.
    
    \item \textbf{Алгебра: Фактор-группы $G/H$.}
    Если $H$ — нормальная подгруппа (симметрия части структуры), то $G/H$ — это группа, где эта симметрия «схлопнута» в единицу. Мы игнорируем внутреннюю структуру $H$, глядя на группу «снаружи».
    
    \item \textbf{Алгебра: Построение полей (Фактор-кольца).}
    Как придумать $i$, если его нет? Берем многочлены $\R[x]$ и отождествляем $x^2+1$ с $0$ (факторизуем по идеалу $(x^2+1)$). В полученном мире $x^2 = -1$.
    $$\C \cong \R[x] / (x^2 + 1)$$
    
    \item \textbf{Логика: Алгебра Линденбаума-Тарского.}
    Множество всех логических формул огромно. Но если мы отождествим формулы $\phi$ и $\psi$, когда $\phi \leftrightarrow \psi$ доказуемо (они логически эквивалентны), мы получим булеву алгебру, описывающую саму структуру истины в данной теории.
    
    \item \textbf{Физика: Калибровочная инвариантность (Электродинамика).}
    Векторный потенциал $\vec{A}$ не измерим напрямую. Потенциалы $\vec{A}$ и $\vec{A} + \nabla \chi$ описывают \textit{одно и то же} магнитное поле $\vec{B} = \nabla \times \vec{A}$.
    Физическое состояние — это не конкретный потенциал, а класс эквивалентности всех таких потенциалов.
    \textit{Вывод:} Реальность — это фактор-пространство математических переменных по группе калибровочных преобразований.
    
    \item \textbf{Физика: Пространство-время в ОТО.}
    В Общей Теории Относительности координаты не имеют физического смысла (они условны). Физическая реальность — это многообразие фактор-по-диффеоморфизмам (заменам координат). Точка пространства-времени определяется не координатами $x,y,z,t$, а событиями, происходящими в ней.
    
    \item \textbf{Лингвистика и Семантика.}
    «Смысл» слова можно определить как класс эквивалентности всех его употреблений или синонимов в разных контекстах. Абстракция — это факторизация конкретных примеров по отношению подобия.
\end{enumerate}


    \begin{remark}[Факторизация как выделение]
    Факторизация --- это выделение части булеана (поскольку классы эквивалентности суть подмножества базового множества). Но мы выделяем факторизацию как отдельный метод, поскольку это наиболее заметный и мощный из методов выделения, в особенности когда речь идет о связи с операциями.
    \end{remark}

    
    \subsection{Метод пополнения («Затыкание дыр»)}

    Если факторизация удаляет избыточную информацию, то \textbf{пополнение} добавляет недостающую. Мы обнаруживаем, что в нашем пространстве есть «дыры» — процессы, которые должны куда-то привести, но не приводят, или уравнения, которые должны иметь решение, но не имеют.
    
    Пополнение переводит «потенциальное» в «актуальное». Мы объявляем саму проблему (отсутствие предела или решения) новым объектом, который эту проблему решает.
    
    
    Рассмотрим основные виды пополнения:
    
    \subsubsection{Алгебраическое пополнение (Замыкание)}
    Цель: Сделать обратные операции возможными.
    \begin{itemize}
        \item \textbf{Группа Гротендика ($\N \to \Z$):} Мы хотим вычитать. Мы строим $\Z$ как пары $\langle a, b\rangle$ из $\N$, факторизованные по отношению эквивалентности.
        \item \textbf{Поле частных ($\Z \to \Q$):} Мы хотим делить. Строим дроби как пары (тоже факторизация).
        \item \textbf{Алгебраическое замыкание ($\R \to \C$):} Мы хотим, чтобы любой многочлен имел корень. Поле $\C$ «совершенно» в алгебраическом смысле.
    \end{itemize}
    \textbf{Механизм:} Обычно используется конструкция пар + факторизация.
    
    \subsubsection{Пополнение порядка (Дедекиндовы сечения)}
    Цель: Обеспечить существование точных граней (супремумов и инфимумов).
    Используется для перехода $\Q \to \R$. Определение сечения является фундаментальным понятием анализа и теории множеств.
    
        
    \begin{definition}[Сечение]
    Мы разрезаем линейно упорядоченное множество $L$ на два класса $L=A\sqcup B$ так, что все элементы $A$ лежат ниже всех элементов $B$. Сечением называется либо пара $\langle A, B\rangle$, либо (как правило) нижний класс сечения, т.е. множество $A$.
    \end{definition}
    Если в $A$ нет максимума, а в $B$ нет минимума, то между ними «дырка». Пополнение объявляет саму эту пару $\langle A, B\rangle$ новым числом, заполняющим дыру.

    Важный момент: так же, как и в случае факторизации, пополнение должно <<уважать>> имеющуюся структуру (или ту ее часть, которую мы хотим оставить в пополнении). Это значит, что для новых точек пространства должны быть корректным способом доопределены операции и отношения исходной структуры.
    
\begin{example}[Сложение сечений]
Пусть $A, B$ --- два нижних класса сечений рациональных чисел. Тогда их сумма:
\[ A + B \eqdef \{ q \in \Q \mid \exists a \in A,\, b \in B:\ q = a + b \} \]
Это подмножество всех рациональных чисел, которые можно получить как сумму одного числа из $A$ и одного из $B$.

Такие поэлементные сумма и произведение множеств называются, соответственно, суммой и произведением Минковского.

\textbf{Пример:} Пусть $A$ и $B$ --- сечения, соответствующие числам $\sqrt{2}$ и $\sqrt{3}$ (то есть, $A = \{q \in \Q: q < \sqrt{2}\}$, $B = \{q \in \Q: q < \sqrt{3}\}$). Тогда $A+B$ будет соответствовать сечению, определяющему число $\sqrt{2}+\sqrt{3}$.

Таким образом, мы \emph{продлеваем} операцию сложения с рациональных чисел на вещественные, определённые как сечения рациональных чисел.
\end{example}
    
\subsubsection{Метрическое пополнение (Фундаментальные последовательности)}

Если у нас задана метрика на множестве $X$, то мы можем построить пополнение этого пространства фундаментальными последовательностями.

\begin{definition}[Фундаментальная последовательность]
Фундаментальной последовательностью в метрическом пространстве $\langle X, d \rangle$ называется последовательность $\langle x_n \rangle$ такая, что для любого $\epsilon > 0$ существует $N\in\om$ такое, что для всех $m, n > N$ выполняется $d(x_m, x_n) < \epsilon$.
\end{definition}

На множестве фундаментальных последовательностей можно ввести отношение эквивалентности:
\begin{definition}[Эквивалентность фундаментальных последовательностей]
Две фундаментальные последовательности $\langle x_n \rangle$ и $\langle y_n \rangle$ эквивалентны, если для любого $\epsilon > 0$ существует $N\in\om$ такое, что для всех $m, n > N$ выполняется $d(x_m, y_m) < \epsilon$.
\end{definition}
Эквивалентные фундаментальные последовательности определяют один и тот же элемент пополнения. Это отношение эквивалентности называется отношением Коши.

\begin{definition}[Пополнение метрического пространства]
Пополнением метрического пространства $\langle X, d \rangle$ называется метрическое пространство $\langle X^*, d^* \rangle$, такое что:
\begin{itemize}
    \item $X^*$ --- фактор-множество множества всех фундаментальных последовательностей из $X$ по отношению эквивалентности.
    \item $d^*$ --- метрика на $X^*$, такая что для всех $x, y \in X^*$ выполняется $d^*(x, y) = \lim_{n\to\infty} d(x_n, y_n)$, где $\langle x_n \rangle$ и $\langle y_n \rangle$ --- фундаментальные последовательности из $X$, представляющие классы эквивалентности $x$ и $y$ соответственно.
\end{itemize}
\end{definition}

\begin{itemize}
    \item \textbf{От $\Q$ к $\R$ (Метод Кантора):} Вещественные числа — это классы эквивалентности фундаментальных последовательностей рациональных чисел.
    \item \textbf{От $C[a,b]$ к $L^2$:} 
    Пространство непрерывных функций с интегральной нормой $\|f\| = \sqrt{\int f^2}$ не является полным (предел последовательности непрерывных функций может быть разрывным). 
    Пополнив это пространство фундаментальными последовательностями, мы получаем пространство $L^2$ (функций, интегрируемых по Лебегу).
    \textit{Смысл:} Сложные, «рваные» функции Лебега возникают как пределы «гладких» непрерывных функций.
\end{itemize}
\textbf{Философия:} Динамический подход. Объект определяется не своей статической формой, а процессом приближения к нему.

\subsubsection{Топологическое пополнение (Компактификация)}
    Цель: Сделать пространство «компактным» (замкнутым и ограниченным), чтобы контролировать поведение на бесконечности.
    
    \begin{itemize}
        \item \textbf{Одноточечная компактификация (Александрова):}
        Мы берем плоскость $\R^2$ (которая бесконечна и незамкнута) и добавляем к ней одну точку $\infty$. Плоскость «сворачивается» в сферу (Сфера Римана $\C \cup \{\infty\}$).
        \item \textbf{Проективное пополнение:}
        Добавление «линии горизонта» к плоскости, где пересекаются параллельные прямые.
    \end{itemize}
    \textbf{Применение:} В комплексном анализе и алгебраической геометрии удобнее считать, что $1/0 = \infty$, чем говорить, что деление невозможно.
    
\subsubsection{Пополнение меры (Лебег и Каратеодори)}

\begin{itemize}
    \item \textbf{Алгебраическое пополнение (Лебег):}
    У нас уже есть готовое пространство с мерой (например, Борелевская мера на прямой). Мы обнаруживаем «баг»: существуют $A \subset B$, где $\mu(B)=0$, но $A$ не входит в $\sigma$-алгебру.
    \textit{Решение:} Мы вручную расширяем алгебру, добавляя туда все такие $A$. Это «ремонт» системы.

    Таким образом, мера Лебега — это пополнение меры Бореля.
    
    \item \textbf{Конструктивное пополнение (Каратеодори):}
    Более глубокий метод, используемый для построения мер «с нуля» (например, меры Хаусдорфа для фракталов).
    \begin{enumerate}
        \item Мы вводим \textit{внешнюю меру} $\mu^*$ для \textbf{всех} подмножеств (грубая оценка «сверху»).
        \item Мы объявляем «реальными» (измеримыми) только те множества $E$, которые проходят \textbf{тест Каратеодори}: множество $E$ измеримо, если оно «аккуратно разрезает» любое множество $X$:
        $$ \mu^*(X) = \mu^*(X \cap E) + \mu^*(X \setminus E) $$
    \end{enumerate}
    \textit{Результат:} Полученная мера автоматически является полной. «Ничтожные» множества проходят тест автоматически.
\end{itemize}
\textbf{Философия:} 
\begin{itemize}
    \item Лебег: Реальность несовершенна, её надо достраивать.
    \item Каратеодори: Реальность — это то, что корректно взаимодействует с наблюдателем (тестовым множеством). В этом подходе полнота возникает сама собой.
\end{itemize}

\begin{remark}[Жордан vs Каратеодори: Сэндвич против Ножа]
Часто измеримость интуитивно понимают как совпадение внешней и внутренней оценок (метод Жордана):
$$ \mu_*(E) = \mu^*(E) $$
Это работает для простых фигур, но требует построения внутренней меры.

Метод Каратеодори (используемый для меры Лебега) работает только с внешней мерой $\mu^*$ и использует «тестовое множество» $X$.
\begin{itemize}
    \item $X$ — это произвольное «тело» в пространстве.
    \item $E$ признается измеримым, если оно разбивает любое тело $X$ аддитивно: «Объем целого равен сумме объемов частей».
\end{itemize}
Это позволяет избежать проблем с внутренней мерой для сложных фрактальных структур.
\end{remark}

\subsubsection{Пополнение с помощью ультрафильтров}

\begin{definition}[Фильтр и ультрафильтр]
\textbf{Фильтром} на множестве $X$ называется непустое семейство подмножеств $F \subseteq \Pcal(X)$ такое, что:
\begin{itemize}
    \item $\emptyset \notin F$;
    \item $A, B \in F \to A \cap B \in F$;
    \item $(A \in F \land A \subseteq B \subseteq X) \to B \in F$.
\end{itemize}
\end{definition}

Фильтр можно представлять себе как семейство больших множеств.

\begin{definition}[Ультрафильтр]
\textbf{Ультрафильтром} на множестве $X$ называется любой максимальный по вложению фильтр $\mathcal{U}$ на $X$.

Замечаение: в общем случае существование ультрафильтра доказывается с помощью леммы Цорна (т.е. аксиомы выбора).
\end{definition}

Фильтр/ультрафильтр называется главным, если он порождается одним элементом, т.е. состоит из всех надмножеств некоторого синглетона. Соответственно, остальные фильтры называются неглавными.

Неглавные ультрафильтры используются для построения ультрастепеней множеств, также с помощью факторизации. Таким способом конструируются, например, гипердействительные числа, пополняющие множество действительных чисел бесконечно малыми и бесконечно большими числами.

Понятие фильтра и ультрафильтра легко обобщается на произвольные множества с частичным порядком.


\subsection{Метод Порождения}

Этот метод реализует принцип «архивации» реальности. Вместо того чтобы описывать бесконечный объект целиком, мы указываем лишь малый набор «семян» (образующих) и алгоритм их развертывания (замыкание относительно операций). Он удобен и как для описания существующего пространства, и для построения новых пространств (или подпространств).

\begin{definition}[Замыкание и Порождение]
Пусть $X$ — множество с набором операций. Пусть $S \subset X$.
\textit{Порождённая подструктура} $\langle S \rangle$ — это наименьшая подструктура в $X$, содержащая $S$ и замкнутая относительно всех операций (вновь вспоминаем о том, что изменение базового множества обязано <<уважать>> операции сигнатуры).
$$ \langle S \rangle = \bigcap \{ Y \subseteq X \mid S \subseteq Y, Y \text{ --- подструктура} \} $$
Конструктивно это означает, что $\langle S \rangle$ состоит из всех элементов, которые можно получить из $S$, применяя операции конечное число раз.
\end{definition}

Мы выделяем три уровня строгости этого метода:

\subsubsection{Алгебраическое порождение (Образующие)}
Применяется в группах, кольцах, алгебрах.
\begin{itemize}
    \item \textbf{Принцип:} Любой элемент сложной структуры — это «слово», составленное из букв алфавита (образующих).
    \item \textbf{Пример:} Группа целых чисел $\Z$ порождается одним элементом $\{1\}$ и операциями $+/-$. Число $5 = 1+1+1+1+1$. Вся бесконечность $\Z$ свернута в единицу.
    \item \textbf{Свобода:} Если между образующими нет связей (соотношений), структура называется \textit{свободной}. Это «чистый лист» генерации.
\end{itemize}

\subsubsection{Линейное порождение (Базис)}
Применяется в векторных пространствах. Самый жесткий и удобный вид порождения.
\begin{itemize}
    \item \textbf{Принцип:} Мы требуем не просто возможности породить всё (системы образующих), но и \textit{минимальности} и \textit{единственности} разложения. 
    \item \textbf{Базис:} Набор векторов $\{e_i\}$, такой что любой вектор $v$ единственным образом представляется как $v = \sum \lambda_i e_i$. 
    
    Как видим, здесь на помощь приходит вспомогательная числовая структура --- поле скаляров, поскольку любой вектор представляется как линейная комбинация базисных векторов с коэффициентами из этого поля. В этом состоит принципиальное отличие базиса от системы образующих.
    \item \textbf{Философия:} Базис вводит систему координат. Он превращает абстрактный вектор в набор чисел и по сути сводит произвольное линейное пространство к пространству $\R^n$ (или $C^n$ для комплексного случая). Это предельная форма редукции.
\end{itemize}

\subsubsection{Топологическое порождение (База и Предбаза)}
Применяется для задания «текстуры» пространства. Описывать все открытые множества $\tau$ невозможно (их слишком много).
\begin{itemize}
    \item \textbf{База топологии:} Набор простых открытых множеств, из которых любые другие получаются объединением.
    \item \textbf{Предбаза:} Набор множеств, из которых база(!) получается конечными пересечениями (например, открытые шары).
    \item \textbf{Пример:} Чтобы задать топологию на $\R$, достаточно задать интервалы $(a, b)$. Вся сложнейшая структура непрерывности выводится из понятия интервала.
    \item \textbf{$\sigma$-алгебра:} Борелевская $\sigma$-алгебра порождается открытыми множествами. Мы не можем явно описать все борелевские множества, но мы знаем, что они «вырастают» из интервалов с помощью счетных операций.
\end{itemize}

\begin{remark}[Информационная сложность]
Метод порождения напрямую связан с понятием \textbf{размерности}.
\begin{itemize}
    \item Размерность векторного пространства = число элементов в базисе.
    \item Вес топологического пространства = минимальная мощность базы.
\end{itemize}
Это мера сложности мира: сколько независимых параметров нужно, чтобы полностью его описать.
\end{remark}






\subsection*{Резюме: Матрица методов конструирования}
\addcontentsline{toc}{subsection}{Резюме: Матрица методов конструирования}

В таблице ниже систематизированы четыре главных способа создания математических объектов, рассмотренных в этом разделе.


\begin{table}[h!]
    \hspace*{-8cm}
    \begin{minipage}{24cm} 
        \caption{Матрица методов конструирования математических структур}
        \label{tab:methods_matrix}
        
        \centering
        \small 
        \renewcommand{\arraystretch}{1.5} 
        
        \begin{tabular}{@{} p{3.5cm} p{3cm} p{5.5cm} p{7cm} p{3.5cm} @{}}
            \toprule
            \textbf{\color{darkblue} Метод} & 
            \textbf{\color{darkblue} Метафора} & 
            \textbf{\color{darkblue} Формальная суть} & 
            \textbf{\color{darkblue} Философский смысл} & 
            \textbf{\color{darkblue} Ключевой пример} \\ 
            \midrule
            
            \textbf{1. Конструирование} \newline (Базовые конструкции) & 
            \textit{Кладка кирпичей} & 
            $X \times Y$ (Пары), $A^B$ (Функции), \newline $X^*$ (Последовательности). \newline Набор «массы» и размерности. & 
            \textbf{Синтез.} Создание сложного из простого через связывание. Объект обретает структуру связей или протяженность. & 
            $\C = \R \times \R$ \newline (Пары чисел) \\ 
            \midrule
            
            \textbf{2. Выделение} \newline (Спецификация) & 
            \textit{Скульптура} & 
            $S = \{x \in U \mid \phi(x)\}$. \newline Понижение энтропии. & 
            \textbf{Отрицание.} \textit{Omnis determinatio est negatio.} Ограничение хаоса универсума законом $\phi$. & 
            Сфера $S^1 \subset \R^2$ \newline ($x^2+y^2=1$) \\ 
            \midrule
            
            \textbf{3. Факторизация} & 
            \textit{Сжатие / Склейка} & 
            $X \to X/{\sim}$. \newline Отождествление неразличимых. & 
            \textbf{Абстракция.} Игнорирование несущественных деталей ради выявления сущности (типа). Переход от индивида к классу. & 
            Время «на часах» \newline $\Z_{12} = \Z / 12\Z$ \\ 
            \midrule
            
            \textbf{4. Пополнение} & 
            \textit{Затыкание дыр} & 
            $X \hookrightarrow \bar{X}$. \newline Добавление пределов. & 
            \textbf{Актуализация.} Превращение бесконечного процесса (последовательности) в завершенный объект (предел). & 
            Вещественные числа \newline $\R = \overline{\Q}$ \\ 
            \midrule
            
            \textbf{5. Порождение и редукция} & 
            \textit{Архивация / Семя} & 
            $X = \langle S \rangle$. \newline Развертывание из базиса. & 
            \textbf{Порождение и редукция.} Развертывание бесконечного из конечного набора образующих и правил; обратно — сведение бесконечного многообразия к конечному базису. & 
            Базис $\vec{i}, \vec{j}, \vec{k}$ в $\R^3$ \\ 
            \bottomrule
        \end{tabular}
        
        \vspace{0.5cm} 
        
        \begin{remark}[Взаимодействие методов в реальных моделях]
            Почти любой сложный объект требует применения всей цепочки методов. 
            Например, построение пространства $L^2$:
            $$ 
            \underbrace{C[a,b]}_{\text{Конструирование}} 
            \xrightarrow[\text{Функционал } (\Phi)]{\text{Норма } \|\cdot\|} 
            \text{Метрическое пр-во} 
            \xrightarrow{\text{Выделение}} 
            \text{Фундамент. посл.} 
            \xrightarrow[\text{{$/_\sim$} (Факторизация)}]{\text{Пополнение}} 
            L^2 
            $$
        \end{remark}
    \end{minipage}
\end{table}



\begin{exercise}[Мастерская: Конструирование миров]
    Ниже предлагаются задачи на применение методов построения пространств. В каждой задаче важен не столько расчет, сколько интерпретация полученной структуры.
    
    \begin{enumerate}
        \item \textbf{Факторизация: Геометрия «Pac-Man».}
        Рассмотрим квадрат $X = [0, 1] \times [0, 1]$. Введем отношение эквивалентности, отождествляющее противоположные края:
        $$(0, y) \sim (1, y) \quad \text{и} \quad (x, 0) \sim (x, 1)$$
        \begin{enumerate}
            \item Что представляет собой пространство $X/{\sim}$ топологически? (Представьте, что вы склеиваете лист бумаги).
            \item Опишите движение частицы, которая летит вправо и достигает границы $x=1$. Где она окажется?
            \item Что изменится, если одну пару краев склеить «с перекруткой»: $(0, y) \sim (1, 1-y)$? (Бутылка Клейна / Лента Мёбиуса).
        \end{enumerate}
    
        \item \textbf{Факторизация + Пополнение: Пространство Звуков.}
        Физически звук — это частота $f \in \R_{>0}$. Однако для человеческого слуха ноты, отличающиеся в 2 раза (октава), звучат «одинаково» (нота Ля 440 Гц и 880 Гц).
        \begin{enumerate}
            \item Введем отношение: $f_1 \sim f_2 \iff f_2 = 2^k f_1$ ($k \in \Z$).
            \item Опишите фактор-пространство $\R_{>0} / {\sim}$. Является ли оно прямой линией или окружностью?
            \item Как на этой модели выглядит музыкальная гамма?
        \end{enumerate}
    
        \item \textbf{Выделение: Конфигурационное пространство.}
        Пусть два робота перемещаются по одной кольцевой рельсе (окружность $S^1$).
        \begin{enumerate}
            \item Универсум состояний системы — это тор $U = S^1 \times S^1$ (положение первого и второго робота).
            \item Роботы не могут занимать одно и то же место (столкновение). Запишите предикат $\phi(x, y)$, который нужно исключить.
            \item Опишите топологию «безопасного» пространства $S = \{ (x,y) \in U \mid x \neq y \}$. (Что останется от тора, если вырезать диагональ?)
        \end{enumerate}
    
        \item \textbf{Пополнение: Проективная прямая.}
        Рассмотрим прямую линию $L$. Наблюдатель находится в точке $O$ вне прямой.
        \begin{enumerate}
            \item Каждой точке $x \in L$ соответствует луч зрения из $O$.
            \item Есть ли луч, который не пересекает прямую $L$? (Луч, параллельный $L$).
            \item Метод пополнения: мы добавляем к $L$ одну точку $\infty$, соответствующую этому параллельному лучу.
            \item Топологически, во что превращается прямая после добавления $\infty$ (если соединить «концы»)?
        \end{enumerate}
        
        \item \textbf{Порождение: Вавилонская библиотека.}
        Пусть $A = \{a, b, \dots, z, \text{пробел}\}$ — алфавит.
        \begin{enumerate}
            \item Пространство $A^*$ (Звезда Клини) содержит все возможные тексты, включая «Гамлета» и бессмысленные наборы букв.
            \item Оцените мощность множества $A^*$ (счетно или несчетно?).
            \item Пусть $L \subset A^*$ — множество осмысленных текстов (выделение). Является ли свойство «быть осмысленным» формализуемым предикатом?
        \end{enumerate}
    \end{enumerate}
    \end{exercise}




\newpage
\section{Числовые системы}

\subsection{Натуральные числа $\N$}

\begin{definition}[Натуральные числа фон Неймана]
Натуральные числа определяются рекурсивно через аксиому бесконечности:
\begin{itemize}
    \item $0 = \emptyset$
    \item $1 = \{0\} = \{\emptyset\}$
    \item $2 = \{0, 1\} = \{\emptyset, \{\emptyset\}\}$
    \item $n + 1 = n \cup \{n\}$
    \item $\N = \bigcap \{X \mid 0 \in X \land \forall x (x \in X \to x \cup \{x\} \in X)\}$ (наименьшее индуктивное множество)
\end{itemize}
\end{definition}

\begin{definition}[Арифметика на $\N$]
\begin{itemize}
    \item \textbf{Сложение:} Определяется рекурсией: $n + 0 = n$, $n + (m + 1) = (n + m) + 1$
    \item \textbf{Умножение:} Определяется рекурсией: $n \cdot 0 = 0$, $n \cdot (m + 1) = n \cdot m + n$
    \item \textbf{Порядок:} $n \le m \iff n \in m \lor n = m$
\end{itemize}
\end{definition}

\begin{exercise}[Арифметика Пеано: «2x2=4»]
    Используя только рекурсивные определения сложения и умножения, докажите формально, что $2 \cdot 2 = 4$.
    \begin{enumerate}
        \item Вспомните определения: $1 = S(0)$, $2 = S(1)$, $3 = S(2)$, $4 = S(3)$.
        \item Используйте правило умножения: $n \cdot S(m) = n \cdot m + n$.
        \item Используйте правило сложения: $n + S(m) = S(n + m)$.
        \item \textbf{Шаги:}
        \begin{itemize}
            \item Запишите $2 \cdot 2$ как $2 \cdot S(1)$.
            \item Раскройте по рекурсии до $2 \cdot 1 + 2$, затем до $2 \cdot 0 + 2 + 2$.
            \item Вычислите сумму $2+2$ через функцию следования $S$, чтобы получить 4.
        \end{itemize}
    \end{enumerate}
    Это упражнение показывает, как сложная таблица умножения сводится к механическому перекладыванию «камушков» (функции $S$).
    \end{exercise}

\subsection{Немного алгебры: группы, кольца, поля}

\marginfigure[1][0.75]{
\begin{tikzpicture}[
    node distance=0.7cm,
    box/.style={
        rectangle,
        draw=darkblue,
        very thick,
        rounded corners=3pt,
        fill=blue!5,
        text width=8.6cm,
        align=left,
        font=\footnotesize
    },
    boxgrp/.style={box, fill=green!12},
    boxring/.style={box, fill=purple!12},
    boxfield/.style={box, fill=orange!12},
    arrow/.style={->, very thick, color=darkgray}
]
    \node[box] (gpd) {\textbf{Группоид (магма)}: одна операция $\circ$\\
    {\scriptsize Аксиомы: нет.}\\
    {\scriptsize Пример: $(\N,2^x)$}};
    \node[box, below=of gpd] (sg) {+\textbf{ассоциативность} $\Rightarrow$ \textbf{полугруппа}\\
    {\scriptsize Аксиомы: $(a\circ b)\circ c=a\circ (b\circ c)$.}\\
    {\scriptsize Пример: $(\{a,b\}^*,\star)$ (конкатенация слов)}};
    \node[box, below=of sg] (mon) {+\textbf{нейтральные элементы} (лев./прав.)  $\Rightarrow$ \textbf{моноид}\\
    {\scriptsize Аксиомы: $a\circ \e=\e\circ a=a$.}\\
    {\scriptsize Пример: $(\N,+,0)$, $(\N,\cdot,1)$.}};
    \node[boxgrp, below=of mon] (grp) {+\textbf{обратимость} по $\circ$ $\Rightarrow$ \textbf{группа} $(G,\e,\circ)$\\
    {\scriptsize Аксиома: $\forall a\, \exists a^{-1}: a\circ a^{-1}=\e=a^{-1}\circ a$.}\\
    {\scriptsize Пример: $(\Z,0,+)$, группа перестановок}};
    \node[box, below=of grp] (ab) {+\textbf{коммутативность} по $+$ $\Rightarrow$ \textbf{абелева группа} $(A,+,0)$\\
    {\scriptsize Аксиома: $a+b=b+a$.}\\
    {\scriptsize Пример: $(\Z,0,+)$.}};
    \node[boxring, below=of ab] (ring) {Две операции: $+$ (абелева группа) и $\cdot$ (группоид).\\
    +\textbf{дистрибутивность} (лев./прав.) $\Rightarrow$ \textbf{кольцо} $(R,0,+,\cdot)$\\
    {\scriptsize Аксиомы: $a(b+c)=ab+ac$; $(a+b)c=ac+bc$.}\\
    {\scriptsize Определение кольца: абелева группа по $+$, группоид по $\cdot$, дистрибутивность.}\\
    {\scriptsize Пример: $(2\Z,0,+,\cdot)$.}};
    \node[box, below=of ring] (uring) {+\textbf{единица} $1$ (лев./прав.) $\Rightarrow$ \textbf{унитальное кольцо} (кольцо с единицей) $(R,0,1,+,\cdot)$\\
    {\scriptsize Аксиомы: $a\cdot1=1\cdot a=a$.}\\
    {\scriptsize Пример: $(\Z,0,1,+,\cdot)$, кольцо многочленов}};
    \node[box, below=of uring] (cring) {+\textbf{коммутативность} по $\cdot$ $\Rightarrow$ \textbf{коммутативное кольцо} $(R,0,+,\cdot)$\\
    {\scriptsize Аксиома: $ab=ba$.}\\
    {\scriptsize Пример: $\Z$ или $\Z_n$.}};
    \node[box, below=of cring] (idom) {+\textbf{отсутствие делителей нуля} $\Rightarrow$ \textbf{область целостности}\\
    {\scriptsize Определение: коммутативное унитальное кольцо без делителей нуля.}\\
    {\scriptsize Аксиома: $ab=0\Rightarrow a=0\lor b=0$.}\\
    {\scriptsize Пример: $\Z$, $\Z[x]$, $\Q[x]$.}};
    \node[boxfield, below=of idom] (field) {+\textbf{обратимость ненулевых} по $\cdot$ $\Rightarrow$ \textbf{поле} $(K,0,1,+,\cdot)$\\
    {\scriptsize Аксиома: $\forall a\neq0\, \exists a^{-1}: aa^{-1}=1=a^{-1}a$.}\\
    {\scriptsize Определение поля: коммутативное унитальное кольцо, в котором каждый ненулевой элемент обратим.}\\
    {\scriptsize Пример: $(\Q,0,1,+,\cdot)$.}};
    \node[box, below=of field] (oring) {+\textbf{порядок, совместимый с $+$ и $\cdot$} $\Rightarrow$ \textbf{упорядоченное кольцо/поле} $(R,0,1,+,\cdot,<)$\\
    {\scriptsize Аксиомы: порядок линейный; $a<b\Rightarrow a+c<b+c$; $0<a,0<b\Rightarrow 0<ab$.}\\
    {\scriptsize Пример: $(\Z,0,1,+,\cdot,<)$, $(\Q,0,1,+,\cdot,<)$.}};
    \node[box, below=of oring] (dord) {+\textbf{дискретность порядка} (лев./прав.) $\Rightarrow$ \textbf{дискретно упорядоченное кольцо/поле} $(R,0,1,+,\cdot,<)$\\
    {\scriptsize Аксиомы: если $a$ не максимум, то $\exists$ правый сосед; если $a$ не минимум, то $\exists$ левый сосед.}\\
    {\scriptsize Пример: $(\Z,0,1,+,\cdot,<)$.}};
    \node[box, below=of dord] (dens) {+\textbf{плотность порядка} $\Rightarrow$ \textbf{плотно упорядоченное кольцо/поле} $(R,0,1,+,\cdot,<)$\\
    {\scriptsize Аксиома: $\forall a<b\, \exists c: a<c<b$.}\\
    {\scriptsize Пример: $(\Q,0,1,+,\cdot,<)$.}};
    \node[box, below=of dens] (cont) {+\textbf{непрерывность порядка} $\Rightarrow$ \textbf{непрерывно упорядоченное кольцо/поле} $(R,0,1,+,\cdot,<)$\\
    {\scriptsize Аксиома: у каждого непустого ограниченного сверху множества есть супремум.}\\
    {\scriptsize Пример: $(\R,+,\cdot,<)$.}};

    \draw[arrow] (gpd) -- (sg);
    \draw[arrow] (sg) -- (mon);
    \draw[arrow] (mon) -- (grp);
    \draw[arrow] (grp) -- (ab);
    \draw[arrow] (ab) -- (ring);
    \draw[arrow] (ring) -- (uring);
    \draw[arrow] (uring) -- (cring);
    \draw[arrow] (cring) -- (idom);
    \draw[arrow] (idom) -- (field);
    \draw[arrow] (field) -- (oring);
    \draw[arrow] (oring) -- (dord);
    \draw[arrow] (dord) -- (dens);
    \draw[arrow] (dens) -- (cont);
\end{tikzpicture}
}

\begin{exercise}[Понимание групп, колец, полей]
Определите, какими структурами являются объекты ниже (группоид, полугруппа, моноид, группа, кольцо, унитальное/коммутативное кольцо, поле, упорядоченное кольцо/поле). Кратко обоснуйте.
\begin{enumerate}
    \item $(\N,+)$ и $(\N,\cdot)$.
    \item $(\Z,+)$ и $(\Z,\cdot)$.
    \item $(2\Z,+,\cdot)$: это кольцо? унитальное?
    \item $(\R\setminus\{0\},\cdot)$: это группа? абелева?
    \item Являются ли полем $\Z/5\Z$ и $\Z/6\Z$? Если нет, то чем они являются?
    \item Проверьте, что $K[x]$ --- унитальное коммутативное кольцо, если $K$ --- унитальное коммутативное кольцо (структура $K[x]$ --- множество многочленов с коэффициентами из $K$ с обычными операциями сложения и умножения).
    \item В кольце: докажите, что $a\cdot 0=0$ и $0\cdot a=0$.
    \item В кольце: докажите, что $-a=(-1)\cdot a$ и $(-a)b=-(ab)$.
    \item В кольце: выведите формулы $a(b-c)=ab-ac$ и $(a-b)c=ac-bc$.
    \item В коммутативном кольце: докажите, что из $ab=ac$ следует $a(b-c)=0$.
    \item В упорядоченном кольце: докажите, что если $a>0$ и $b<c$, то $ab<ac$.
    \item В группе: докажите правила сокращения $ab=ac \Rightarrow b=c$ и $ba=ca \Rightarrow b=c$.
    \item В группе: докажите, что уравнение $ax=b$ имеет единственное решение.
    \item В кольце с $1$: докажите, что из $ax=0$ следует $x=0$, если $a$ — обратим. 
    \item \textbf{Теорема на доказательство:} Пусть $R$ — конечное коммутативное кольцо с $1$. Докажите, что $R$ — область целостности $\Leftrightarrow$ $R$ — поле. (Область целостности: нет делителей нуля.)\\
    \textit{Подсказка:} рассмотрите отображение $f_a(x)=ax$, воспользуйтесь принципом Дирихле.
    \item \textbf{Теорема на доказательство:} множество обратимых элементов коммутативного кольца с $1$ образует группу относительно умножения.
\end{enumerate}
\end{exercise}

\subsection{Целые числа $\Z$ через классы эквивалентности}

\begin{definition}[Отношение эквивалентности на $\N \times \N$]
Определим отношение $\sim$ на $\N \times \N$:
$$\langle a, b \rangle \sim \langle c, d \rangle \iff a + d = b + c$$
Это отношение эквивалентности (рефлексивно, симметрично, транзитивно).
\end{definition}

\begin{definition}[Целые числа]
Множество целых чисел $\Z$ определяется как фактор-множество:
$$\Z = (\N \times \N) / \sim$$
Класс эквивалентности $\langle a, b \rangle$ интерпретируется как число $a - b$.
\end{definition}

\begin{definition}[Арифметика на $\Z$]\;
\begin{enumerate}
    \item \textbf{Сложение:} $[\langle a, b \rangle] + [\langle c, d \rangle] = [\langle a + c, b + d \rangle]$
    \item \textbf{Умножение:} $[\langle a, b \rangle] \cdot [\langle c, d \rangle] = [\langle ac + bd, ad + bc \rangle]$
    \item \textbf{Порядок:} $[\langle a, b \rangle] \le [\langle c, d \rangle] \iff a + d \le b + c$
\end{enumerate}
\end{definition}

\begin{theorem}[Вложение $\N$ в $\Z$]
Отображение $n \mapsto [\langle n, 0 \rangle]$ является изоморфным вложением $\N$ в $\Z$, сохраняющим арифметические операции и порядок.
\end{theorem}


\begin{exercise}[Корректность]
    Мы определили целое число как класс пар: $x = [\langle a, b \rangle]$ (интуитивно $a-b$).
    \begin{enumerate}
        \item \textbf{Тест на минус:} Покажите формально, что $(-1) \cdot (-1) = 1$.
        \begin{itemize}
            \item Представьте $-1$ как класс $[\langle 0, 1 \rangle]$.
            \item Вычислите произведение $[\langle 0, 1 \rangle] \cdot [\langle 0, 1 \rangle]$ по формуле из Определения 3.3.
            \item Убедитесь, что результат эквивалентен $[\langle 1, 0 \rangle]$ (т.е. $1$).
        \end{itemize}
        
        \item \textbf{Корректность:} 
        Операция сложения определена на представителях: $[\langle a, b \rangle] + [\langle c, d \rangle]$.
        Докажите, что результат не зависит от выбора представителя.
        \begin{itemize}
            \item Пусть $\langle a, b \rangle \sim \langle a', b' \rangle$ (то есть $a+b' = b+a'$).
            \item Покажите, что сумма с $\langle c, d \rangle$ даст эквивалентный результат:
            $$ \langle a+c, b+d \rangle \sim \langle a'+c, b'+d \rangle $$
        \end{itemize}
        \textit{Подсказка:} Используйте свойства равенства в $\N$. Без этого доказательства вся конструкция $\Z$ не имеет математического смысла.
    \end{enumerate}
    \end{exercise}

\subsection{Рациональные числа $\Q$ через классы эквивалентности}

\begin{definition}[Отношение эквивалентности на $\Z \times (\Z \setminus \{0\})$]
Определим отношение $\sim$ на $\Z \times (\Z \setminus \{0\})$:
$$\langle a, b \rangle \sim \langle c, d \rangle \iff ad = bc$$
Это отношение эквивалентности.
\end{definition}

\begin{definition}[Рациональные числа]
Множество рациональных чисел $\Q$ определяется как фактор-множество:
$$\Q = (\Z \times (\Z \setminus \{0\})) / \sim$$
Класс эквивалентности $\langle a, b \rangle$ интерпретируется как дробь $\frac{a}{b}$.
\end{definition}

\begin{definition}[Арифметика на $\Q$]\;
\begin{enumerate}
    \item \textbf{Сложение:} $[\langle a, b \rangle] + [\langle c, d \rangle] = [\langle ad + bc, bd \rangle]$
    \item \textbf{Умножение:} $[\langle a, b \rangle] \cdot [\langle c, d \rangle] = [\langle ac, bd \rangle]$
    \item \textbf{Порядок:} Для положительных знаменателей: $[\langle a, b \rangle] \le [\langle c, d \rangle] \iff ad \le bc$
\end{enumerate}
\end{definition}

\begin{theorem}[Вложение $\Z$ в $\Q$]
Отображение $z \mapsto [\langle z, 1 \rangle]$ является изоморфным вложением $\Z$ в $\Q$, сохраняющим арифметические операции и порядок.
\end{theorem}

\begin{exercise}[Рациональная механика]
    \begin{enumerate}
        \item \textbf{Равенство дробей:} 
        Докажите формально, что $\frac{2}{4} = \frac{1}{2}$.
        \begin{itemize}
            \item Запишите это на языке пар: $\langle 2, 4 \rangle \sim \langle 1, 2 \rangle$.
            \item Проверьте условие эквивалентности $ad = bc$.
        \end{itemize}
        
        \item \textbf{Почему нельзя делить на ноль?}
        Попробуйте создать рациональное число $\frac{1}{0}$, то есть пару $\langle 1, 0 \rangle$.
        Попытайтесь проверить его эквивалентность с любой другой дробью $\frac{a}{b}$.
        Что происходит с отношением транзитивности, если допустить $0$ в знаменатель?
        
        \item \textbf{Сложение:}
        Выведите школьную формулу $\frac{a}{b} + \frac{c}{d} = \frac{ad+bc}{bd}$ из определения суммы классов $[\langle a, b \rangle] + [\langle c, d \rangle]$.
    \end{enumerate}
    \end{exercise}

\subsection{Вещественные числа $\R$ через классы эквивалентности}

Существует несколько эквивалентных конструкций $\R$. Рассмотрим построение через фундаментальные последовательности (метод Кантора).

\begin{definition}[Фундаментальная последовательность]
Последовательность рациональных чисел $\langle q_n \rangle$ называется \textit{фундаментальной} (последовательностью Коши), если:
$$\forall \ep > 0 \exists N \in \N: \forall m, n \ge N: |q_m - q_n| < \ep$$
\end{definition}

\begin{definition}[Отношение эквивалентности на фундаментальных последовательностях]
Две фундаментальные последовательности $\langle q_n \rangle$ и $\langle r_n \rangle$ эквивалентны:
$$\langle q_n \rangle \sim \langle r_n \rangle \iff \lim_{n \to \infty} (q_n - r_n) = 0$$
То есть: $\forall \ep > 0 \exists N: \forall n \ge N: |q_n - r_n| < \ep$.
\end{definition}

\begin{definition}[Вещественные числа]
Множество вещественных чисел $\R$ определяется как фактор-множество:
$$\R = \{\langle q_n \rangle \mid \langle q_n \rangle \text{ фундаментальна}\} / \sim$$
\end{definition}

\begin{definition}[Арифметика на $\R$]\;
\begin{enumerate}
    \item \textbf{Сложение:} $[\langle q_n \rangle] + [\langle r_n \rangle] = [\langle q_n + r_n \rangle]$
    \item \textbf{Умножение:} $[\langle q_n \rangle] \cdot [\langle r_n \rangle] = [\langle q_n \cdot r_n \rangle]$
    \item \textbf{Порядок:} $[\langle q_n \rangle] \le [\langle r_n \rangle] \iff \exists N: \forall n \ge N: q_n \le r_n$
\end{enumerate}
\end{definition}

\begin{theorem}[Полнота $\R$]
Всякая фундаментальная последовательность в $\R$ сходится к элементу $\R$ (полнота по Коши).
\end{theorem}

\begin{remark}[Альтернативная конструкция через дедекиндовы сечения]
$\R$ также можно построить как множество дедекиндовых сечений рациональных чисел $\Q$. Оба метода эквивалентны.
\end{remark}


\subsection{Комплексные числа $\C$ через классы эквивалентности}

\begin{definition}[Комплексные числа]
Множество комплексных чисел $\C$ определяется как:
$$\C = \R \times \R$$
с операциями:
\begin{itemize}
    \item \textbf{Сложение:} $\langle a, b \rangle + \langle c, d \rangle = \langle a + c, b + d \rangle$
    \item \textbf{Умножение:} $\langle a, b \rangle \cdot \langle c, d \rangle = \langle ac - bd, ad + bc \rangle$
\end{itemize}
\end{definition}

\begin{theorem}[Вложение $\R$ в $\C$]
Отображение $r \mapsto \langle r, 0 \rangle$ является изоморфным вложением $\R$ в $\C$, сохраняющим арифметические операции.
\end{theorem}

\begin{definition}[Мнимая единица]
Элемент $i = \langle 0, 1 \rangle$ называется мнимой единицей. Имеем $i^2 = -1$.
\end{definition}

\begin{theorem}[Алгебраическая замкнутость $\C$]
Поле $\C$ алгебраически замкнуто: всякий многочлен с коэффициентами в $\C$ имеет корень в $\C$ (основная теорема алгебры).
\end{theorem}


\subsection{Алгебра Линденбаума--Тарского}

В этом пункте мы покажем, как из классического исчисления высказываний получить \emph{алгебру}, в которой формулы становятся элементами булевой структуры. Это один из самых «мостовых» объектов между логикой (синтаксисом) и алгеброй (структурой операций).

\subsubsection{Определение (классическое исчисление высказываний)}

Рассмотрим язык исчисления высказываний с переменными $p,q,r,\dots$ и связками $\neg, \wedge, \vee, \to$ (можно взять любой эквивалентный набор связок). Пусть $\vdash$ --- отношение выводимости в фиксированной корректной и полной дедуктивной системе классической логики (например, гильбертовой или натурального вывода).

Определим отношение эквивалентности на множестве формул:
$$
\varphi \sim \psi \iff \vdash (\varphi \leftrightarrow \psi),
$$
то есть формулы отождествляются, если они \emph{логически эквивалентны} (в смысле доказуемости двусторонней импликации).

\begin{definition}[Алгебра Линденбаума--Тарского (ЛТ)]
\textbf{Алгеброй Линденбаума--Тарского} классического исчисления высказываний называется фактор-множество
$$
\mathsf{LT} \eqdef \{\varphi\}/\sim
$$
классов эквивалентности формул, снабжённое операциями
$$
[\varphi]\wedge[\psi] \eqdef [\varphi\wedge\psi],\quad
[\varphi]\vee[\psi] \eqdef [\varphi\vee\psi],\quad
\neg[\varphi] \eqdef [\neg\varphi],
$$
и константами $0 \eqdef [\bot]$, $1 \eqdef [\top]$ (где $\top$ можно понимать как $\varphi\to\varphi$, а $\bot$ как $\neg\top$).
\end{definition}

\begin{remark}[Корректность определения]
Нужно проверить, что операции на классах эквивалентности определены корректно (ill-defined). То есть из $\varphi\sim\varphi'$ и $\psi\sim\psi'$ следует
$\varphi\wedge\psi \sim \varphi'\wedge\psi'$, аналогично для $\vee$ и $\neg$.
Это верно, потому что в классической логике эквивалентность сохраняется при подстановке в контекст (конгруэнтность относительно связок).
\end{remark}

\subsubsection{Алгебра ЛТ как булева алгебра}

\begin{theorem}[Алгебра ЛТ является булевой алгеброй]
Структура $\langle \mathsf{LT}, \vee, \wedge, \neg, 0, 1 \rangle$ является булевой алгеброй. То есть для всех $x,y,z \in \mathsf{LT}$ выполнены аксиомы:
\begin{enumerate}[label=(\alph*), nosep]
    \item \textbf{Коммутативность:} $x\vee y = y\vee x$, $x\wedge y = y\wedge x$.
    \item \textbf{Ассоциативность:} $(x\vee y)\vee z = x\vee(y\vee z)$, $(x\wedge y)\wedge z = x\wedge(y\wedge z)$.
    \item \textbf{Дистрибутивность:} $x\wedge(y\vee z) = (x\wedge y)\vee(x\wedge z)$ и $x\vee(y\wedge z) = (x\vee y)\wedge(x\vee z)$.
    \item \textbf{Нейтральные элементы:} $x\vee 0 = x$, $x\wedge 1 = x$.
    \item \textbf{Дополнения:} $x\vee \neg x = 1$, $x\wedge \neg x = 0$.
\end{enumerate}
\end{theorem}

\begin{exercise}[базовые тождества в булевой алгебре]
\begin{enumerate}
    \item \textbf{Идемпотентность:} докажите, что для всех $x \in \mathsf{LT}$ выполнено
    $$x \vee x = x, \qquad x \wedge x = x.$$
    
    \item \textbf{Нули и единицы:} докажите, что для всех $x \in \mathsf{LT}$ выполнено
    $$x \vee 1 = 1,\qquad x \wedge 1 = x,\qquad x \vee 0 = x,\qquad x \wedge 0 = 0.$$
    
    \item \textbf{Законы поглощения:} докажите, что для всех $x,y \in \mathsf{LT}$ выполнено
    $$x \vee (x \wedge y) = x,\qquad x \wedge (x \vee y) = x.$$
\end{enumerate}
\end{exercise}

\begin{remark}[Двойственность meet/join (принцип двойственности)]
В булевой алгебре операции $\,\wedge\,$ (meet) и $\,\vee\,$ (join) стоят в строгой симметрии.
Если в любой тождественно истинной формуле булевой алгебры одновременно заменить
$$
\wedge \leftrightarrow \vee,\qquad 0 \leftrightarrow 1,
$$
то получится новая тождественно истинная формула.
Например, из дистрибутивности $x\wedge(y\vee z)=(x\wedge y)\vee(x\wedge z)$ получается
$x\vee(y\wedge z)=(x\vee y)\wedge(x\vee z)$; из $x\vee 0=x$ получаем $x\wedge 1=x$.
\end{remark}

\begin{remark}[Почему $(\mathsf{LT},\wedge,\vee)$ не является кольцом]
Идея двойственности подсказывает, что $\wedge$ и $\vee$ не могут играть роли умножения и сложения в кольце:
в кольце операции $+$ и $\cdot$ не дуальны и не симметричны (к ним не применим принцип «поменять местами и одновременно поменять нуль с единицей»).

\medskip
\textbf{Проблема «вычитания» (обратного элемента):} Чтобы структура была кольцом, она должна быть абелевой группой по «сложению». То есть для каждого $x$ должен существовать элемент $-x$ такой, что $x + (-x) = 0$.
Если в качестве сложения взять операцию $\vee$ (join), то нулём будет $0$ (так как $x \vee 0 = x$).
Но как найти обратный элемент для $1$? Нужно решить уравнение $1 \vee y = 0$.
В булевой алгебре (и в логическом смысле) это невозможно: дизъюнкция с единицей всегда даёт $1$.
Значит, $\vee$ не порождает группу по сложению, а следовательно --- и кольцо.

Чтобы получить кольцевую структуру из булевой алгебры, сложение берут не как $\vee$, а как \textbf{симметрическую разность} $\,\oplus\,$ (см. следующий пункт).
\end{remark}

\subsubsection{Симметрическая разность и булево кольцо}

\begin{definition}[Симметрическая разность]
В булевой алгебре определим операцию \textbf{симметрической разности}:
$$
x \oplus y \eqdef (x \wedge \neg y) \vee (\neg x \wedge y).
$$
\end{definition}

\begin{theorem}[Булево кольцо из булевой алгебры]
Пусть $B$ --- булева алгебра. Тогда структура $\langle B, \oplus, \wedge, 0, 1 \rangle$ является \textbf{булевым кольцом}, то есть кольцом с единицей, в котором выполнено тождество $x^2=x$ для всех $x$ (идемпотентность).
Эквивалентно, для всех $x,y,z$ выполнены:
\begin{enumerate}[label=(\alph*), nosep]
    \item \textbf{Абелева группа по $\oplus$:} $(B,\oplus,0)$ (ассоциативность, коммутативность, нейтральный элемент $0$, каждый элемент сам себе обратный: $x\oplus x=0$).
    \item \textbf{Моноид по $\wedge$:} $(B,\wedge,1)$ (ассоциативность, нейтральный элемент $1$). 
    \item \textbf{Дистрибутивность:} $x\wedge(y\oplus z) = (x\wedge y)\oplus(x\wedge z)$.
    \item \textbf{Идемпотентность:} $x\wedge x = x$.
\end{enumerate}
\end{theorem}

\begin{exercise}[Из аксиом булева кольца]
Рассмотрим структуру, удовлетворяющую свойствам (a)--(d) теоремы о булевом кольце (со сложением $\oplus$ как $+$ и умножением $\wedge$ как $\cdot$). Обозначим $xy \eqdef x\wedge y$, $x+y \eqdef x\oplus y$.
\begin{enumerate}
    \item Докажите, что $x+x=0$ для всех $x$. (\textit{Подсказка:} раскройте $(x+x)^2$.)
    \item Докажите, что $xy=yx$ для всех $x,y$. (\textit{Подсказка:} раскройте $(x+y)^2$.)
    \item Положив $\neg x = 1+x$ и $x\vee y = x+y+xy$, выведите закон исключённого третьего: $x \vee \neg x = 1$.
    \item Выведите законы де Моргана: $\neg(x\vee y) = (\neg x)\wedge(\neg y)$ и $\neg(x\wedge y) = (\neg x)\vee(\neg y)$ (перепишите через $+$ и $xy$).
    \item Выведите закон поглощения: $x(x\vee y)=x$.
\end{enumerate}
\end{exercise}

\begin{remark}[Взаимная определимость булевых алгебр и булевых колец]
Булева алгебра $\to$ булево кольцо: умножение $xy$ берём как $x\wedge y$, сложение $x+y$ как $x\oplus y$, нуль $0$ как $0$, единицу как $1$.

Булево кольцо $\to$ булева алгебра: если $R$ --- булево кольцо (коммутативное, с $1$, и $x^2=x$), то можно определить
$$
x\wedge y \eqdef xy,\quad
\neg x \eqdef 1 + x,\quad
x\vee y \eqdef x + y + xy,
$$
и получить булеву алгебру $\langle R, \vee, \wedge, \neg, 0, 1\rangle$.
Эти конструкции обратимы и не требуют дополнительных условий.
\end{remark}

\subsubsection{Теорема Стоуна (представление)}

\begin{theorem}[Теорема Стоуна]
Любое булево кольцо изоморфно некоторому кольцу подмножеств: существует множество $X$ и изоморфизм булевых колец
$$
R \cong \mathcal{P}(X),
$$
где на $\mathcal{P}(X)$ сложение --- симметрическая разность, а умножение --- пересечение.
\end{theorem}

\begin{remark}[Эпиморфизм на $\mathbb{F}_2$, ядро, ультрафильтр и «истинные» значения]
Поле $\mathbb{F}_2 = \{0,1\}$ можно рассматривать как булево кольцо (сложение и умножение по $\bmod 2$, тождество $x^2=x$).
Для любого сюръективного гомоморфизма булевых колец $\varphi \colon R \twoheadrightarrow \mathbb{F}_2$ множество
$\ker \varphi = \{ x \in R : \varphi(x)=0 \}$ есть \textbf{максимальный идеал} в $R$ (образ --- поле).
Дополнение к ядру, т.е.
$$
R \setminus \ker\varphi = \{ x \in R : \varphi(x)=1 \},
$$
--- это \textbf{ультрафильтр} на соответствующей булевой алгебре (множество «истинных» элементов в двузначной модели).

Таким образом, эпиморфизм на $\mathbb{F}_2$ --- это то же самое, что выбор ультрафильтра; в логической интерпретации алгебры ЛТ ультрафильтры служат \textbf{реализациями истины}: им соответствуют полные согласованные атрибуты истинности (двузначные оценки), различающие формулы только по классам эквивалентности $\sim$.
\end{remark}

\subsubsection{Хейтинговы алгебры и интуиционистская логика}

Если вместо классической логики взять интуиционистскую, то получаем не булеву, а \textbf{хейтингову алгебру}: в ней нет закона исключённого третьего $x\vee \neg x = 1$ в общем виде, а центральной становится операция импликации как \emph{псевдодополнение}. Такие структуры естественно связаны не с кольцами, а с \emph{топологией} (например, алгебра открытых множеств топологического пространства образует хейтингову алгебру).

\subsubsection{Полиномы Жегалкина и факторизация по идеалам}

Классическая логика также кодируется полиномами над $\mathbb{F}_2$ (алгебра Жегалкина).

\begin{definition}[Идеал кольца]
Пусть $R$ --- коммутативное кольцо. Непустое подмножество $I \subseteq R$ называется \textbf{идеалом}, если:
\begin{enumerate}[label=(\roman*), nosep]
    \item $a,b\in I \Rightarrow a-b\in I$;
    \item $a\in I,\ r\in R \Rightarrow ra\in I$.
\end{enumerate}
\end{definition}

\begin{exercise}[Следствия аксиом идеала]
Пусть $I$ --- идеал в коммутативном кольце $R$.
\begin{enumerate}
    \item Докажите, что $0\in I$.
    \item Докажите, что если $a\in I$, то $-a\in I$.
    \item Докажите, что если $a,b\in I$, то $a+b\in I$.
\end{enumerate}
\end{exercise}

Обычно идеал задают \textbf{множеством образующих} $i_1, i_2, \dots$ (их может быть и бесконечно много): тогда $I$ --- это множество всех \emph{конечных} линейных комбинаций $\sum_k r_k i_k$ с коэффициентами $r_k\in R$. Пишут $\langle i_1, i_2, \dots\rangle$ или, в коммутативном случае, нередко $(i_1, i_2, \dots)$. Например, $n\mathbb{Z} = \{nk : k\in\mathbb{Z}\}$ --- идеал в $\mathbb{Z}$ (все целые, кратные $n$).

\begin{exercise}[Идеал, порождённый подмножеством]
Пусть $R$ --- коммутативное кольцо, $\varnothing \neq S \subseteq R$. Рассмотрим множество $J$ всех \emph{конечных} линейных комбинаций элементов $S$ с коэффициентами из $R$, то есть
$$
J = \Bigl\{ \sum_{k=1}^m r_k s_k : m\ge 1,\ r_k\in R,\ s_k\in S \Bigr\}.
$$
Докажите, что $J$ --- идеал в $R$ (он совпадает с $\langle S\rangle$, идеалом с образующими из $S$).
\end{exercise}

\begin{theorem}[Жегалкин: булево кольцо как фактор]
Рассмотрим полиномиальное кольцо $\mathbb{F}_2[x_1,x_2,\dots]$. Факторизуя по идеалу
$$
I = (x_1^2-x_1,\ x_2^2-x_2,\ \dots),
$$
получаем булево кольцо:
$$
\mathbb{F}_2[x_1,x_2,\dots]/I,
$$
в котором для каждого класса выполнено $x_i^2 = x_i$. Элементы этого кольца (классы полиномов) называются \textbf{полиномами Жегалкина} и задают булевы функции: сложение соответствует XOR (симметрической разности), умножение соответствует AND.
\end{theorem}

\begin{exercise}[Формулы и полиномы Жегалкина]
Работайте в $\mathbb{F}_2[x,y,z,\dots]$ с соотношениями $t^2=t$; интерпретируйте сложение как XOR, умножение как $\wedge$, а дополнение как $\neg t \leftrightarrow 1+t$.
Вспомните или выведите, что $x\vee y$ в кольце записывается как $x+y+xy$.

\begin{enumerate}
    \item \textbf{К полиному.} Запишите полином Жегалкина для формул $\neg(x\wedge y)$, $x\vee y$, импликации $x\to y$ и равносильности $x\leftrightarrow y$.
    \item \textbf{К формуле.} Для полинома $1+x+xy$ укажите эквивалентную ему формулу со связками $\vee,\wedge,\neg$ (можно коротко объяснить по таблице истинности).
    \item \textbf{К полиному.} Переведите в полином выражение $(x\vee y)\wedge \neg z$.
    \item \textbf{К формуле.} Полином $x+y+z$ задаёт XOR трёх переменных в $\mathbb{F}_2$. Запишите ту же функцию с помощью $\oplus$ между высказываниями или через стандартные $\vee,\wedge,\neg$.
    \item \textbf{Смешанный пример.} Найдите полином для $(x\to y)\wedge (y\to z)$; затем по полученному полиному проверьте значение на наборе $(x,y,z)=(1,0,1)$.
\end{enumerate}
\end{exercise}


\subsection{Альтернативный взгляд: Числа как Операторы (Матрицы)}

Существует другой мощный способ определения числовых систем, который интерпретирует число не как статическую величину, а как \textbf{действие} (линейный оператор).
В этом подходе числа представляются квадратными матрицами, а операции над ними — сложением и умножением матриц.

\subsubsection{Комплексные числа как матрицы поворота}

Вместо того чтобы вводить мистическую «мнимую единицу» $i = \sqrt{-1}$, мы можем определить комплексные числа как подкольцо матриц $2 \times 2$ с вещественными коэффициентами.

\begin{definition}[Матричная модель $\C$]
Сопоставим $1$ и $i$ следующим матрицам:
$$ \mathbf{1} \mapsto \begin{pmatrix} 1 & 0 \\ 0 & 1 \end{pmatrix} = I, \quad \mathbf{i} \mapsto \begin{pmatrix} 0 & -1 \\ 1 & 0 \end{pmatrix} = J $$
Тогда произвольное комплексное число $z = a + bi$ представляется матрицей:
$$ Z = aI + bJ = \begin{pmatrix} a & -b \\ b & a \end{pmatrix} $$
\end{definition}

\begin{theorem}[Изоморфизм]
Множество таких матриц с обычными матричными операциями изоморфно полю $\C$.
\begin{itemize}
    \item Проверка $i^2 = -1$:
    $$ J^2 = \begin{pmatrix} 0 & -1 \\ 1 & 0 \end{pmatrix} \begin{pmatrix} 0 & -1 \\ 1 & 0 \end{pmatrix} = \begin{pmatrix} -1 & 0 \\ 0 & -1 \end{pmatrix} = -I $$
    \item Модуль числа равен корню из определителя: $|z| = \sqrt{\det(Z)} = \sqrt{a^2 + b^2}$.
\end{itemize}
\textbf{Геометрический смысл:} Умножение на матрицу $Z$ соответствует растяжению вектора в $\sqrt{a^2+b^2}$ раз и повороту на угол $\ph$ ($Z = r \cdot R_\ph$).
\end{theorem}

\subsubsection{Процедура Кэли-Диксона (Удвоение алгебр)}

Как пойти дальше? Гамильтон искал способ умножать тройки чисел, но понял, что нужно умножать четверки. Процедура Кэли-Диксона формализует этот процесс «удвоения» размерности.

\begin{definition}[Удвоение]
Пусть $A$ — алгебра с сопряжением (где для каждого $x$ определено $\bar{x}$).
Определим новую алгебру $A'$ на парах элементов $(a, b)$, где $a, b \in A$:
\begin{itemize}
    \item \textbf{Сложение:} покомпонентное.
    \item \textbf{Умножение:} $(a, b) \cdot (c, d) = (ac - d\bar{b}, \, \bar{a}d + cb)$.
    \item \textbf{Сопряжение:} $\overline{(a, b)} = (\bar{a}, -b)$.
\end{itemize}
\end{definition}

Эта рекурсивная процедура порождает цепочку алгебр:
$$ \R \xrightarrow{\text{удвоение}} \C \xrightarrow{\text{удвоение}} \mathbb{H} \xrightarrow{\text{удвоение}} \mathbb{O} \dots $$

\subsubsection{Кватернионы $\mathbb{H}$ (Quaternions)}

Применив процедуру Кэли-Диксона к полю $\C$, мы получаем алгебру Кватернионов.
Элементы $\mathbb{H}$ — это пары комплексных чисел $(u, v)$.
Покажем, как стандартный базис $\{1, i, j, k\}$ выражается через эти пары.

\begin{enumerate}
    \item \textbf{Определение базиса:}
    Введем базисные элементы следующим образом:
    \begin{itemize}
        \item $\mathbf{1} = (1, 0)$ — вещественная единица.
        \item $\mathbf{i} = (i, 0)$ — мнимая единица из $\C$.
        \item $\mathbf{j} = (0, 1)$ — новая мнимая единица.
        \item $\mathbf{k} = (0, -i)$ — произведение, замыкающее таблицу умножения.
    \end{itemize}
    
    \item \textbf{Проверка таблицы умножения:}
    Используем формулу произведения в алгебре Кэли-Диксона: 
    $$(a, b) \cdot (c, d) = (ac - d\bar{b}, \, \bar{a}d + cb)$$
    
    Вычислим произведение $\mathbf{i}$ и $\mathbf{j}$:
    $$ \mathbf{i} \cdot \mathbf{j} = (i, 0) \cdot (0, 1) = (i\cdot 0 - 1\cdot \bar{0}, \, \bar{i}\cdot 1 + 0\cdot 0) = (0, -i) $$
    Согласно нашему определению, $(0, -i)$ — это в точности элемент $\mathbf{k}$.
    Следовательно, выполняется классическое соотношение $ij = k$.
    
    Проверим квадрат $\mathbf{j}$:
    $$ \mathbf{j}^2 = (0, 1) \cdot (0, 1) = (0 - 1\cdot \bar{1}, \, 0 + 1\cdot 0) = (-1, 0) = -\mathbf{1} $$
    
    Аналогично проверяется, что $k^2 = -1$ и $ji = -k$.
\end{enumerate}

Таким образом, любой кватернион $q$ представляется как пара $(u, v)$, где $u=a+bi, v=c+di$, или в алгебраической форме:
$$ q = (a+bi, 0) + (0, 1)(c-di) = a+bi + jc + kd $$
(Заметим, что из-за некоммутативности порядок множителей при разложении важен, но алгебраическая структура полностью изоморфна классическим кватернионам Гамильтона).

\begin{exercise}[Алгебра вращений]
\begin{enumerate}
    \item Проверьте матричным умножением, что $ij = k$, но $ji = -k$.
    \item Докажите, что для кватернионов выполняется свойство нормы: $\|q_1 q_2\| = \|q_1\| \|q_2\|$ (сумма четырех квадратов).
    \item \textbf{Физический смысл:} Единичные кватернионы образуют группу $SU(2)$, которая описывает вращения в 3D-пространстве (спин электрона). Кватернион $q$ кодирует ось вращения и угол. Это гораздо эффективнее, чем матрицы Эйлера (и используется во всех 3D-движках видеоигр для избежания «шарнирного замка»).
\end{enumerate}
\end{exercise}

\begin{exercise}[Рождение $k$]
    Используя формулу произведения пар $(a, b) \cdot (c, d) = (ac - d\bar{b}, \, \bar{a}d + cb)$:
    \begin{enumerate}
        \item Вычислите произведение $j \cdot i = (0, 1) \cdot (i, 0)$.
        (Ожидаемый ответ: $(0, -i) = -k$).
        \item Вычислите произведение $i \cdot j = (i, 0) \cdot (0, 1)$.
        (Ожидаемый ответ: $(0, i) = k$).
    \end{enumerate}
    Это доказывает некоммутативность: $ij \neq ji$.
    \end{exercise}

\marginfigure[0][0.75]{
\begin{tikzpicture}[
    node distance=0.6cm,
    box/.style={
        rectangle,
        draw=darkblue,
        thick,
        rounded corners=2pt,
        fill=blue!5,
        text width=2.8cm,
        align=center,
        font=\footnotesize
    },
    arrow/.style={->, thick, color=darkgray}
]
    \node[box] (N) {$\N$\\\textbf{натуральные числа}};
    \node[box, right=1.4cm of N] (On) {$\On$\\\textbf{класс ординалов}};
    \node[box, below=of N] (Z) {$\Z$\\\textbf{целые числа}};
    \node[box, right=1.4cm of Z] (Zi) {$\Z[i]$\\\textbf{Гауссовы целые числа}};
    \node[box, below=of Z] (Q) {$\Q$\\\textbf{рациональные числа}};
    \node[box, below=of Q] (R) {$\R$\\\textbf{вещественные числа}};
    \node[box, right=1.4cm of Q] (Qsqrt) {$\Q[\sqrt{2}]$\\\textbf{квадратичное расширение}};
    \node[box, below=of Qsqrt] (A) {$\A$\\\textbf{алгебраические числа}};
    \node[box, below=of A] (Rx) {$\R[x]$\\\textbf{многочлены над $\R$}};
    \node[box, below=of R] (C) {$\C$\\\textbf{комплексные числа}};
    \node[box, below=of C] (H) {$\mathbb{H}$\\\textbf{кватернионы}};
    \node[box, below=of Rx] (Cx) {$\C[x]$\\\textbf{многочлены над $\C$}};

    \draw[arrow] (N) -- (On);
    \draw[arrow] (N) -- (Z);
    \draw[arrow] (Z) -- (Zi);
    \draw[arrow] (Z) -- (Q);
    \draw[arrow] (Q) -- (R);
    \draw[arrow] (Q) -- (Qsqrt);
    \draw[arrow] (Qsqrt) -- (A);
    \draw[arrow] (R) -- (Rx);
    \draw[arrow] (A) -- (C);
    \draw[arrow] (C) -- (H);
    \draw[arrow] (C) -- (Cx);
    \draw[arrow] (R) -- (C);
    \draw[arrow] (Rx) -- (Cx);
\end{tikzpicture}
}

%\newpage


\section{К теоремам Гёделя и Тарского}


\subsection{О стандартной модели арифметики}

Во всём разделе рассматривается стандартная модель $\N$ арифметики Пеано, заданная на ординале $\omega$. 
Первопорядковая арифметика \PA корректна относительно $\N$, а потому непротиворечива. Этими условиями мы будем пользоваться всюду в дальнейшем.


\subsection{Лемма о диагонализации}

Рассмотрим язык первого порядка арифметики Пеано $\Lang(\PA)$. Для всех формул этого языка зададим \textit{геделевскую нумерацию}: функцию $\Numcode{\ph}:\Lang(\PA)\to\N$ (стандартная модель натурального ряда, ординал $\om$). Можем считать, например, что нумерация построена на основе ОТА, главное, чтобы процедура кодирования/декодирования была разрешимой [за конечное число шагов применения простых арифметических действий можно определить, является ли данное число кодом формулы или не является, и какой именно формулы].


На множестве $\Lang(\PA)\times\N$ определим (на мета-уровне) \textit{примитивно рекурсивный} оператор подстановки $Subst_a[\ph,n]$, который делает следующее: подставляет нумерал $\numeral n$ (по определению это арифметический терм $\numeral n\eqdef\underbrace{S\dots S}_{n}0$) в исходный текст формулы $\ph$ всюду вместо переменной $a$ (там, где она имеет свободное вхождение). В результате получается новая формула, отличие которой от $\ph$ состоит в том, что свободная переменная $a$ заменилась конкретным нумералом $\numeral n$. Коротко определение этого оператора можно записать так:
$$
Subst_a[\ph,n]\eqdef\ph(\numeral n),
$$
но тут надо понимать, что код подставляется всегда только вместо свободной переменной $a$, так что если в формуле $\ph$ нет свободной переменной $a$, то оператор вернет исходную формулу $\ph$ без изменений, поэтому сокращение $\ph(\numeral n)$ не вполне корректно без этого уточнения.

На основе оператора $Subst_a$ определим новую \textit{примитивно рекурсивную} функцию:
$$
F(n)\eqdef
\begin{cases}
\Numcode{Subst_a[\ph(a),n]}, & \mbox{если существует }\ph(a)\in\Lang(\PA):\;n=\Numcode{\ph(a)},\\
0, & \mbox{иначе}.
\end{cases}
$$
Пользуясь сокращенным способом записи оператора подстановки, заметим, что для произвольной формулы $\ph(a)$ имеет место равенство
$$
F(\Numcode{\ph(a)})=\Numcode{\ph(\numeral{\Numcode{\ph(a)}})}.
$$
Примитивная рекурсивность функции $F$ очевидна из того, что процедура восстановления текста формулы (последовательности кодов букв) разрешима для любого входного $n$ за конечное число шагов, а подстановка нумерала $\numeral n$ вместо переменной $a$ (в этой последователности) также примитивно рекурсивна (не использует цикл while).

Далее, согласно теореме о представимости примитивно рекурсивных функций в арифметике существует $\Sigma_1$-формула $\de(a,b)$ со свободными переменными $a$ и $b$ такая, что:
$$
\N\models\de(\numeral n,\numeral m)\quad\Longleftrightarrow\quad F(n)=m.
$$
Более того,
имеет место сильная представимость функции $F$ в арифметике:
$$
\PA\vdash \forall y\bigl(\de(\numeral n,y)\leftrightarrow y=\numeral{F(n)}\bigr),
$$
то есть \PA <<знает>>, что формула $\de(a,b)$ задает тотальную функцию, откуда легко показать, что
$$
\PA\vdash \exists y\;\de(\numeral n,y)\leftrightarrow \de(\numeral n,\numeral{F(n)}),
$$
и, учитывая определение функции $F$:
$$
\PA\vdash \exists y\;\de(\numeral n,y)\land(y=\numeral{F(n)})\leftrightarrow \de(\numeral n,\numeral{F(n)}).
$$


\begin{theorem}[лемма о диагонализации/неподвижной точке]\;\\
Для любой формулы $\Phi(a)\in\Lang(\PA)$ с единственной свободной переменной $a$ существует предложение $\theta\in\Lang(\PA)$ такое, что
$$
\PA\vdash \theta\leftrightarrow \Phi(\numeral{\Numcode\theta}).
$$
\end{theorem}
\begin{proof} Пусть $\Phi(a)$ --- произвольная формула \PA с единственной свободной переменной $a$. Рассмотрим формулу с одной свободной переменной $a$:
$$
\xi(a)\equiv \exists y\;\de(a,y)\land\Phi(y).
$$
Наконец, пусть 
$$
\theta\eqdef\xi(\numeral{\Numcode{\xi(a)}}).
$$
Тогда
\begin{align*}
\mbox{\PA}\;\vdash \;
& \theta \leftrightarrow \exists y\;\de(\numeral{\Numcode{\xi(a)}},y)\land\Phi(y) \\ 
& \theta \leftrightarrow \exists y\;(y=\numeral{F(\Numcode{\xi(a)})})\land\Phi(y) \\
& \theta \leftrightarrow \Phi(\numeral{F(\Numcode{\xi(a)})}) \\
& \theta \leftrightarrow \Phi(\numeral{\Numcode{\xi(\numeral{\Numcode{\xi(a)}})}}) \\
& \theta \leftrightarrow \Phi(\numeral{\Numcode\theta}).
\end{align*}
\end{proof}

\begin{exercise}[Лемма о диагонализации: неподвижные точки]
    \begin{enumerate}
        \item \textbf{Простые крайние случаи.}
        Пусть $\phi(x)$:
        \begin{enumerate}[label=(\alph*)]
            \item тождественно ложна;
            \item тождественно истинна.
        \end{enumerate}
        Найдите неподвижную точку в каждом случае (при любом корректном кодировании формул).

        \item \textbf{Свойство последней цифры кода.}
        Пусть $\phi(x)$ выражает тот факт, что число $x$ оканчивается на $6$.
        \begin{enumerate}[label=(\alph*)]
            \item Запишите формулу на языке $\PA$, выражающую это свойство.
            \item Пусть гёделева нумерация устроена так, что закрывающая скобка имеет код $6$. Найдите ложную неподвижную точку формулы $\phi$.
            \item В тех же условиях найдите истинную неподвижную точку формулы $\phi$.
        \end{enumerate}

        \item \textbf{Фиксированное кодирование символов.}
        Пусть заданы коды символов:
        $[0]=1,\ [S]=2,\ [+]=3,\ [=]=4$ (коды остальных символов неважны).
        \begin{enumerate}[label=(\alph*)]
            \item Запишите формулу $\phi(x)$, выражающую «$x$ не делится на $3$».
            \item Найдите неподвижную точку формулы $\phi$ (\textit{Подсказка:} отталкивайтесь от простейших атомарных формул.).
        \end{enumerate}
    \end{enumerate}
\end{exercise}

\subsection{Теоремы Гёделя и Тарского}

Умея кодировать формулы, мы умеем также кодировать и  доказательства примитивно рекурсивным способом, а именно: по числу $n$ мы можем установить за конечное число шагов, является ли оно кодом доказательства какой-то формулы и вернуть геделев номер доказанной формулы, или же оно не является кодом доказательства.

Соответственно, мы можем найти формулу $\Pr(a,b)$ со свободными переменными $a$ и $b$, для которой истинность $\Pr(\numeral n,\numeral m)$ равносильна тому, что $m$ является кодом доказательства формулы с кодом $n$.
Более того, эта формула имеет класс $\Delta_0$ (использует только ограниченные кванторы).

\begin{theorem}[Теорема Гёделя о неполноте]\;\\
Существует предложение $G\in\Lang(\PA)$ такое, что:
$\PA\nvdash G,\quad \PA\nvdash\neg G$.
\end{theorem}
\begin{proof}
Пусть $G$ --- неподвижная точка формулы $\neg\exists x\;\Pr(a,x)$ с единственной свободной переменной $a$. Тогда
$$
\PA\vdash G\leftrightarrow \neg\exists x\;\Pr(\numeral{\Numcode{G}},x).
$$

Если бы $\PA\vdash G$, то существовал бы код $p$ доказательства $G$.
Значит, $\PA\vdash \Pr(\numeral{\Numcode G},\numeral p)$, откуда $\PA\vdash \exists x\,\Pr(\numeral{\Numcode G},x)$, и в силу диагональной эквиваленции мы получаем доказательство противоречия в \PA.

Если бы $\PA\vdash\neg G$, то по диагональной эквиваленции получится, что
$$
\PA\vdash \exists x\;\Pr(\numeral{\Numcode G},x).
$$
Это $\Sigma_1$-формула. По корректности \PA она истинна в $\N$, значит существует стандартное доказательство $G$, то есть $\PA\vdash G$. Вместе с $\PA\vdash\neg G$ получаем противоречие.
\end{proof}

\begin{remark}
    При стандартном кодировании отношение $\Pr(a,b)$ является примитивно рекурсивным и даже представимо $\Delta_0$-формулой: проверка кода конечного доказательства ограничена длиной этого кода. Поэтому формула доказуемости $\exists y\,\Pr(a,y)$ имеет класс $\Sigma_1$.
    Следовательно, гёделево предложение $G$, удовлетворяющее условию теоремы, можно считать $\Pi_1$-предложением: оно утверждает отсутствие кода доказательства.
\end{remark}

\begin{theorem}[Вторая теорема Гёделя о неполноте]\;\\
    $\PA\nvdash \Con(\PA)$,
    где $\Con(\PA)\eqdef \neg\exists x\;\Pr(\numeral{\Numcode{0=1}},x)$.
\end{theorem}
\begin{proof} 
Возьмем то же $G$, что и в предыдущей теореме.
Формализуя само доказательство первой теоремы Гёделя в языке \PA, мы получаем, что
$$
\PA\vdash \Con(\PA)\to\neg\exists x\;\Pr(\numeral{\Numcode G},x),
$$
и по диагональной эквиваленции мы получаем, что
$$
\PA\vdash \Con(\PA)\to G.
$$

Следовательно, если предположить, что $\PA\vdash \Con(\PA)$, то мы получаем, что $\PA\vdash G$.
Но по первой же теореме мы знаем, что $\PA\nvdash G$. Значит, $\PA\nvdash \Con(\PA)$.
\end{proof}

Эти рассуждения можно обобщить на широкий класс теорий.
\begin{theorem}[Теорема Гёделя--Россера о неполноте]\;\\
Пусть $T$ --- теория в языке первого порядка, удовлетворяющая условиям:
\begin{enumerate}
    \item $T$ непротиворечива;
    \item $T$ эффективно аксиоматизируема (для примитивной рекурсивности предиката доказуемости);
    \item $T$ интерпретирует арифметику Робинсона Q (для возможности производить геделевское кодирование формул и доказательств).
\end{enumerate}
Тогда существует предложение $G\in\Lang(T)$ такое, что: 
$T\nvdash G,\quad T\nvdash\neg G$.
\end{theorem}



\begin{theorem}[Тарский, о невыразимости истины]
    Не существует формулы $T(a)\in\Lang(\PA)$ с одной свободной переменной $a$ такой, что для любого предложения $\ph\in\Lang(\PA)$:
    $$
    \N\models T(\numeral{\Numcode\ph})\;\Longleftrightarrow\;\N\models\ph.
    $$
\end{theorem}
\begin{proof}
        Предположим, что такая формула $T(a)$ существует. 
        Применим лемму о диагонализации к формуле $\neg T(a)$:
        получаем предложение $\lambda$ такое, что
        $$
        \PA\vdash \lambda\leftrightarrow\neg T(\numeral{\Numcode\lambda}).
        $$
        По корректности \PA эта эквиваленция истинна в $\N$:
        $$
        \N\models \lambda\;\Longleftrightarrow\;\N\models\neg T(\numeral{\Numcode\lambda}).
        $$
        Но по предположению на $T$: $\N\models T(\numeral{\Numcode\lambda})\Longleftrightarrow\N\models\lambda$. 
        Подставляя, получаем $\N\models\lambda\Longleftrightarrow\N\models\neg\lambda$ --- противоречие.
\end{proof}


\begin{remark}
Положение спасает переход ко второпорядковой арифметике, в которой мы уже можем определить предикат истины. Это будет формула второго порядка (с кванторами по подмножествам натурального ряда $\N$), истинная на всех и только геделевых номерах первопорядковых предложений арифметики. По сути она будет переформулировкой теоретико-множественного определения истины в стандартной модели ($\N\models$).

Однако предикат истины для второпорядковых формул также нельзя определить 
средствами второго порядка --- запрет Тарского применяется на каждом уровне иерархии, что порождает иерархию языков и метаязыков по Тарскому.
\end{remark}



\subsection{Предикат доказуемости и логика GL}

\subsubsection{Предикат доказуемости}

Зафиксируем рекурсивно аксиоматизируемую теорию $T\supseteq\PA$. Арифметизация синтаксиса даёт две ключевые формулы.

\begin{itemize}
    \item \textbf{Отношение доказательства} $\Pr(a,b)$: примитивно рекурсивная ($\Delta_0$) формула, истинная в $\N$ тогда и только тогда, когда $b$ — код вывода формулы с кодом $a$ в теории $T$.
    \item \textbf{Предикат доказуемости}: $\square_T\,\ph \eqdef \exists x\,\Pr(\Numcode{\ph},\,x)$.
    Это $\Sigma_1$-формула; она утверждает, что существует доказательство формулы $\ph$ в~$T$.
\end{itemize}

Ниже мы пишем $\square\ph$ вместо $\square_T\ph$, фиксируя $T = \PA$. Определим также
$\Con(\PA) \eqdef \neg\square\bot$, где $\bot \eqdef (0=1)$.



\subsubsection{Условия Гильберта–Бернайса–Лёба}

\begin{theorem}[Условия HBL]
Для любых формул $\ph,\psi$ языка $\PA$ выполняются:
\begin{enumerate}[label=\upshape(D\arabic*)]
    \item Если $\PA\vdash\ph$, то $\PA\vdash\square\ph$.\hfill\textup{(необходимость)}
    \item $\PA\vdash\square(\ph\to\psi)\to(\square\ph\to\square\psi)$.\hfill\textup{(дистрибутивность/нормальность)}
    \item $\PA\vdash\square\ph\to\square\square\ph$.\hfill\textup{(интроспекция/транзитивность)}
\end{enumerate}
\end{theorem}

\begin{proof}
\textbf{(D1).} Пусть $\PA\vdash\ph$, то есть существует конкретный код вывода $d\in\N$. Проверка того, что $d$ является корректным выводом $\ph$ в \PA, есть конечная алгоритмическая процедура; она кодируется $\Delta_0$-формулой. Следовательно, $\PA\vdash\Pr(\Numcode{\ph},\numeral d)$, и применяя правило введения экзистенциального квантора, получаем $\PA\vdash\square\ph$.

\textbf{(D2).} Работаем внутри \PA. Предположим $\square(\ph\to\psi)$ и $\square\ph$: существуют коды $d_1,d_2$ с $\Pr(\Numcode{\ph\to\psi},d_1)$ и $\Pr(\Numcode{\ph},d_2)$. Применение правила \textit{modus ponens} к выводам --- это примитивно рекурсивная операция на кодах: из $d_1$ и $d_2$ конструктивно строится код вывода $d$ для $\psi$. Арифметика \PA доказывает, что эта конструкция работает, откуда $\square\psi$.

\textbf{(D3).} Заметим, что $\square\ph = \exists x\,\Pr(\Numcode{\ph},x)$ является $\Sigma_1$-формулой. Ключевой факт: \PA \emph{$\Sigma_1$-полна} в следующем смысле --- для любой $\Sigma_1$-формулы $\si$,
\[
    \PA\vdash\si\to\square\si.
\]
Содержательно: если $\Sigma_1$-утверждение истинно, то есть существует конкретный свидетель $k$ (для $\exists$-кванторов), из которого \PA умеет построить явный вывод этого утверждения; существование такого вывода \PA сама доказывает. Применяя этот факт к $\si = \square\ph$, получаем $\PA\vdash\square\ph\to\square\square\ph$.
\end{proof}

\subsubsection{Теорема Лёба}

\begin{theorem}[Лёб, 1955]\label{thm:Loeb}
Если $\PA\vdash\square\ph\to\ph$, то $\PA\vdash\ph$.
\end{theorem}

Другими словами: если \PA доказывает «из доказуемости $\ph$ следует $\ph$», то $\ph$ уже доказуема без всякого предположения о доказуемости.

\begin{proof}
По лемме о диагонализации существует предложение $\la$ такое, что
\begin{equation}\label{eq:lob-fp}
    \PA\vdash\la\;\leftrightarrow\;(\square\la\to\ph).
\end{equation}
В частности, из \eqref{eq:lob-fp} следуют обе импликации:
\begin{align}
    \mbox{\PA}&\vdash\la\to(\square\la\to\ph), \label{eq:lob1}\\
    \mbox{\PA}&\vdash(\square\la\to\ph)\to\la. \label{eq:lob2}
\end{align}

\textbf{Шаг 1.} Применим (D1) к \eqref{eq:lob1}: $\PA\vdash\square(\la\to(\square\la\to\ph))$.

\textbf{Шаг 2.} Дважды применим (D2):
\begin{align*}
    \mbox{\PA}&\vdash\square\la\to\square(\square\la\to\ph),\\
    \mbox{\PA}&\vdash\square(\square\la\to\ph)\to(\square\square\la\to\square\ph).
\end{align*}
Стыкуя эти две импликации: $\PA\vdash\square\la\to(\square\square\la\to\square\ph)$.

\textbf{Шаг 3.} Из (D3): $\PA\vdash\square\la\to\square\square\la$. Подставляем в итог шага 2:
\begin{equation}\label{eq:lob4}
    \PA\vdash\square\la\to\square\ph.
\end{equation}

\textbf{Шаг 4.} По гипотезе $\PA\vdash\square\ph\to\ph$. Из \eqref{eq:lob4} цепочкой:
\begin{equation}\label{eq:lob5}
    \PA\vdash\square\la\to\ph.
\end{equation}

\textbf{Шаг 5.} Формула \eqref{eq:lob5} является теоремой \PA; из \eqref{eq:lob2} по \textit{modus ponens}:
\begin{equation}\label{eq:lob6}
    \PA\vdash\la.
\end{equation}

\textbf{Шаг 6.} Из \eqref{eq:lob6} по (D1): $\PA\vdash\square\la$. Подставляем в \eqref{eq:lob5}: $\PA\vdash\ph$.
\end{proof}

\begin{corollary}[Вторая теорема Гёделя о неполноте]
$\PA\nvdash\Con(\PA)$.
\end{corollary}
\begin{proof}
Предположим $\PA\vdash\Con(\PA) = \neg\square\bot$. Это равносильно $\PA\vdash\square\bot\to\bot$. По теореме Лёба (теорема~\ref{thm:Loeb}) получаем $\PA\vdash\bot$ --- противоречие.
\end{proof}


\begin{remark}
В $\mathsf{GL}$ не выполняется аксиома \textbf{T}: $\square A\to A$ (рефлексивность). Это соответствует тому, что $\PA\nvdash\square\ph\to\ph$ для произвольного $\ph$: звуковость (истинность всего доказуемого) не является внутренней теоремой $\PA$ --- она равносильна $\Con(\PA)$ и его усилениям, недоказуемым в~$\PA$. Именно поэтому теорема Лёба нетривиальна: рефлексивность срабатывает лишь тогда, когда $\ph$ и без того уже доказуема.
\end{remark}


\subsection{Дилемма Евтифрона в арифметике}

\subsubsection{Постановка}

\begin{center}
\textit{Выбирают ли боги что-то, потому что оно есть благое, \\
или что-то есть благое, потому что его выбирают боги?}
\end{center}

Чтобы увязать между собой благость и выбор богов, примем соглашение, что боги выбирают благое (с их точки зрения), т.е. выбор богов --- это их решение о том, что что-то есть благое. Тогда дилемма принимает вид:

\begin{center}
    \textit{Считают ли боги благим что-то, потому что оно есть благое, \\
    или что-то есть благое, потому что боги его считают благим?}
\end{center}

Таким образом, у нас появляется два предиката: <<являться благим>> и <<считаться благим богами>>. Это напоминает модальность или самореференцию по типу предиката $\Pr$ в арифметике. Пусть $\ph(x)$ --- предикат <<$x$ является благим>>. Его истинность будет верифицироваться стандартной моделью $\N$, т.е. мета-предикатом $\N\models\ph(\numeral n)$. Здесь требуется подставить в $\ph$ какой-то нумерал, выражающий стандартное значение $x$, поскольку истинность формулы со свободной переменной означает ее тотальную истинность.

Тот факт, что боги считают $\ph(\numeral n)$ истинным, мы можем записать как модальность $\Mod\ph(\numeral n)$. Тогда мы можем переписать дилемму в виде:

\begin{gather*}
    \mbox{(a)}\;\forall n\in\N:\;(\N\models\ph(\numeral n))\implies \Mod\ph(\numeral n)\\
    \mbox{(b)}\;\forall n\in\N:\;\Mod\ph(\numeral n)\implies (\N\models\ph(\numeral n))
\end{gather*}


\subsubsection{Разрешение дилеммы}

Рассмотрим вариант $\Mod=\Box$ (доказуемость в \PA). Тогда дилемма принимает вид:
\begin{gather*}
    \mbox{(a)}\;\forall n\in\N:\;(\N\models\ph(\numeral n))\implies \Box\ph(\numeral n)\\
    \mbox{(b)}\;\forall n\in\N:\;\Box\ph(\numeral n)\implies (\N\models\ph(\numeral n))
\end{gather*}

Пункт (b) есть следствие теоремы о корректности: все доказуемое в \PA истинно в ее модели $\N$, что в нашей интерпретации означает: что-то является благим, потому что боги так решили. Это не означает, что боги видят все благое, но что бы они не выбрали --- обязательно окажется благим (причем не важно, что мы понимаем под словом <<благо>>).

Пункт (a) означает обратное: все истинное доказуемо, что, конечно, не верно в силу теоремы о неполноте \PA. Есть только небольшой класс формул, для которых это верно: это $\Sigma_1$-формулы, т.е. формулы, которые можно записать в виде $\exists y\,\psi(x, y)$, где $\psi$ содержит только ограниченные кванторы. Для таких свойств $\ph$ имеет место равносильность: $\N\models\ph(\numeral n)\Longleftrightarrow\PA\vdash\ph(\numeral n)$. Если благость можно записать такой формулой (то есть выразить ее алгоритмически), то <<боги выбирают благое тогда и только тогда, когда оно на самом деле благое>>.


Рассмотрим в качестве модальности $\Mod$ мета-предикат, утверждающий, что формула $\ph(\numeral n)$ не противоречит \PA, т.е. $\PA\nvdash\neg\ph(\numeral n)$. Иначе эту модальность можно выразить как $\neg\Box\neg$ или $\Diamond$. Тогда дилемма принимает вид:
\begin{gather*}
    \mbox{(a)}\;\forall n\in\N:\;(\N\models\ph(\numeral n))\implies \Diamond\ph(\numeral n)\\
    \mbox{(b)}\;\forall n\in\N:\;\Diamond\ph(\numeral n)\implies (\N\models\ph(\numeral n))
\end{gather*}

Модальность $\Diamond$ является, с одной стороны, не абсолютной, так как не отражает истинность в одной конкретной модели, с другой стороны, даже если ее арифметизировать (через предикат доказуемости $\Pr$), она не выражает доказуемость в \PA в силу неполноты \PA, т.е. в ней есть нечто трансцендентное.


Рассмотрим свойства дилеммы Евтифрона для случая $\Mod=\Diamond$.
Теперь уже пункт (a) является следствием теоремы о корректности: истинное нельзя опровергнуть в \PA (и потому боги его принимают). Таким образом, дилемма разрешается первым кейсом при любой формуле $\ph$. Это не означает, что боги не могут признать благим что-то не являющееся таковым на самом деле, но все благое они признают (при этом опять же неважно, что мы понимаем под термином <<благо>>).


Пункт (b) описывает обратную позицию: «свойство объективно благое, потому что боги не находят в нем изъяна (совместность)». Это утверждение \textbf{верно не для всех} формул, а для отрицаний $\Sigma_1$-формул, т.е. для формул класса $\Pi_1$. Кроме того, оно, очевидно, будет верным для всех истинных в стандартной модели формул $\ph(\numeral n)$.

Таким образом, равносильность условий $\N\models\ph(\numeral n)$ и $\Diamond\ph(\numeral n)$ достигается для $\Pi_1$-формул. Если благость можно записать такой формулой (то есть выразить ее негативно-алгоритмически), то <<боги выбирают благое тогда и только тогда, когда оно на самом деле благое>>.



\newpage

\section{Основы математического анализа}


\subsection{Теория вещественных чисел}

В математической практике построение анализа начинается после того, как поле вещественных чисел $\R$ уже сконструировано методами теории множеств (\ZFC) как минимальное пополнение рациональных чисел (см. главу о числовых структурах). Однако, чтобы изучать сам анализ, внутреннее теоретико-множественное устройство чисел (как сечений Дедекинда или классов эквивалентности Коши) необходимо «забыть». 

В этом разделе мы формализуем это «забывание», рассматривая $\R$ на трёх уровнях выразительности: от элементарной алгебры первого порядка до возвращения в полную теорию множеств для нужд функционального анализа.

\subsubsection{Уровень I: Элементарная алгебра (\RCF)}

Самый экономичный вариант «забыть» конструкцию $\R$ --- ограничить формальный язык. Мы переходим к односортному языку логики первого порядка.

\smallskip\noindent\textbf{Язык.} 
Включает только индивидные переменные $x, y, z, \ldots$, пробегающие элементы числового носителя $\R$.
В сигнатуру языка мы включаем \emph{только} константы $0, 1$, операции $+$ и $\cdot$, а также отношения $=$ и $\le$. Теоретико-множественный символ $\in$ между числами в языке отсутствует, поэтому вопрос о том, из чего состоят числа, становится синтаксически невыразимым.

\smallskip\noindent\textbf{Аксиоматика.} 
Базируется на теории вещественно замкнутых полей (\RCF), которая включает аксиомы упорядоченного поля и бесконечные схемы, задающие вещественную замкнутость.

\begin{axiom}[Поле]\label{ax:field}
$\langle \R, 0, 1, +, \cdot \rangle$ удовлетворяет аксиомам поля:
\begin{itemize}
    \item $(\R, +, 0)$ --- абелева группа;
    \item $(\R \setminus \{0\}, \cdot, 1)$ --- абелева группа;
    \item дистрибутивность: $a(b+c) = ab + ac$.
    \item нетривиальность: $0 \neq 1$.
\end{itemize}
\end{axiom}

\begin{axiom}[Линейный порядок]\label{ax:linear}
Бинарное отношение $\le$ на $\R$ такое, что:
\begin{itemize}
    \item рефлексивность: $a \le a$;
    \item антисимметрия: $a \le b \wedge b \le a \Rightarrow a = b$;
    \item транзитивность: $a \le b \wedge b \le c \Rightarrow a \le c$;
    \item трихотомия: $\forall a, b \in \R$ выполнено ровно одно из $a < b$, $a = b$, $b < a$ (здесь $a < b \iff a \le b \wedge a \neq b$).
\end{itemize}
\end{axiom}

\begin{axiom}[Согласованность порядка с операциями]\label{ax:ordered_field}
\begin{itemize}
    \item $a \le b \Rightarrow a + c \le b + c$;
    \item $(0 \le a \wedge 0 \le b) \Rightarrow 0 \le ab$.
\end{itemize}
\end{axiom}

\begin{remark}[Выводимость $0 < 1$]
    Интересно отметить, что базовое неравенство $0 < 1$ не требуется постулировать отдельно --- оно выводится как теорема. По аксиоме нетривиальности $0 \neq 1$, значит, по трихотомии либо $0 < 1$, либо $1 < 0$. Если допустить, что $1 < 0$, то прибавив $-1$ к обеим частям, получим $0 < -1$. Умножив неравенство $0 \le -1$ само на себя, по аксиоме согласованности с умножением получим $0 \le (-1)\cdot(-1) = 1$. Из аксиом поля следует, что $(-1)\cdot(-1) = 1$, откуда $0 < 1$, что противоречит допущению $1 < 0$. Следовательно, $0 < 1$. Аналогично доказывается, что любой ненулевой квадрат в упорядоченном поле строго положителен ($x^2 > 0$).
\end{remark}

\begin{axiom}[Схема вещественной замкнутости: корень многочлена нечётной степени]\label{ax:rcf-odd-root}
    Для каждого \emph{нечётного} натурального числа $n$ в аксиоматику первого порядка добавляется отдельная аксиома (параметр $n$ пробегает все нечётные числа, образуя \emph{схему}):
    \[
    \forall a_0,\ldots,\forall a_n\in\R\;
    \Bigl(
    a_n \neq 0
    \;\Rightarrow\;
    \exists x\in\R:\;
    \sum_{k=0}^{n} a_k\, x^k = 0
    \Bigr),
    \]
    понимая $x^k$ как итеративную степень с $x^0\eqdef 1$.
\end{axiom}
    
\begin{axiom}[Существование квадратного корня]\label{ax:square-root}
    Каждое неотрицательное число имеет квадратный корень:
    \[
    \forall a \in \R\; \bigl(a \ge 0 \Rightarrow \exists x \in \R:\; x \cdot x = a\bigr).
    \]
\end{axiom}

\smallskip\noindent\textbf{Свойства теории.} По теореме Тарского теория \RCF{} \emph{дедуктивно полна} (для всякой замкнутой формулы выводится либо она, либо её отрицание) и \emph{разрешима} алгоритмически. В ней доказуема Теорема о промежуточном значении, но \emph{только для многочленов}. 
Следует отметить, что для подавляющего большинства инженерных задач и систем автоматизированного проектирования выразительной силы \RCF{} вполне достаточно, а её вычислимая полнота делает работу на этом уровне абсолютно безопасной. Однако для построения теоретического матанализа этого критически не хватает: нельзя сформулировать квантор по подмножествам, аксиому непрерывности или общее определение предела.


\subsubsection{Уровень II: Предикативный анализ (подход Г.~Вейля)}

Чтобы ввести в анализ пределы, необходимо уметь говорить о множествах чисел. Мы переходим к полноценной рабочей теории второго порядка.

\smallskip\noindent\textbf{Язык.} Используется двусортный язык (строчные буквы --- один сорт, заглавные --- другой). 
Индивидные переменные $x,y,\ldots$ пробегают числа, переменные $X,Y,\ldots$ пробегают множества чисел. Для индивидных переменных сохраняются все константы, операции и отношения Уровня~I ($0, 1, +, \cdot, =,\le$). Равенство для чисел и для множеств --- разное.
Сигнатура дополняется предикатным символом принадлежности $x \in X$ и специальной константой $\N$ (второго сорта) для множества натуральных чисел.

Формула называется \textbf{числовой}, если она не содержит кванторов по множествам ($\exists X,\,\forall X$), при этом допустимо использовать множества в качестве параметров.


\smallskip\noindent\textbf{Базовые аксиомы.} Теория включает аксиомы упорядоченного поля \ref{ax:field}, \ref{ax:linear}, \ref{ax:ordered_field}. Оба равенства (чисел и множеств) удовлетворяют стандартным аксиомам равенства (отношение эквивалентности и конгруэнтность относительно языка), но для равенства множеств добавляется стандартная аксиома экстенсиональности:
\[
\forall X \forall Y\; (\forall z\; (z \in X \leftrightarrow z \in Y) \to X = Y)
\]

\begin{axiom}[Схема свертывания]\label{ax:comprehension}
Для любой числовой формулы $\ph(x, \vec A)$:
$$
\exists X\;\forall x\; (x\in X)\leftrightarrow\ph(x, \vec A),
$$
т.е. существует множество $\{x \in \R \mid \ph(x, \vec A)\}$.
\end{axiom}

\smallskip\noindent\textbf{Введение натуральных чисел ($\N$).} 
Даже имея второпорядковый язык и свертывание по числовым формулам, мы не можем определить натуральные числа как пересечение всех индуктивных множеств, так как формула $\forall X (\text{Ind}(X) \to x \in X)$ содержит запрещённый квантор по множествам. Казалось бы, нам достаточно квантора по вещественным числам, поскольку каждое вещественное число из отрезка $[0;1]$ кодирует подмножество натурального ряда, но проблема в том, что для введения такого кодирования нам требуется функция $2^n$ (для определения $n$-го знака вещественного числа), которую можно определить только имея уже систему натуральных чисел. Поэтому мы признаем (следуя Герману Вейлю) натуральные числа <<несократимым минимумом>> для построения математики и вводим их аксиоматически.

\begin{definition}[Индуктивное множество]
Множество $X$ называется \textbf{индуктивным}, если оно удовлетворяет формуле:
\[
  Ind(X)\lradef (0 \in X) \land \forall x\; (x \in X) \to (x + 1 \in X).
\]
\end{definition}

\begin{axiom}[Аксиома натуральных чисел]
\[Ind(\N)\land\forall X\; (Ind(X) \to \N \subseteq X)\]
\end{axiom}

\begin{theorem}[Минимальность и единственность $\N$]
Аксиоматически введённое множество $\N$ является наименьшим индуктивным множеством и определяется однозначно.
\end{theorem}
\pf
\textbf{Минимальность:} очевидно.\\
\textbf{Единственность:} если $\N'$ --- другое множество, удовлетворяющее аксиоме натуральных чисел, то из минимальности $\N$ следует $\N \subseteq \N'$. С другой стороны, симметрично рассуждая, имеем $\N' \subseteq \N$. Значит, $\N = \N'$.
\epf

\smallskip\noindent\textbf{Определимые функции и последовательности.}
Хотя в нашем двусортном языке отсутствуют переменные для произвольных функций (поскольку функция --- это множество упорядоченных пар, то есть объект более высокого типа, нежели подмножество прямой), мы можем полноценно работать с \emph{определимыми} функциями и последовательностями. Последовательность вещественных чисел $x_n$ задаётся формулой $\ph(n, x)$, для которой в теории доказуемо условие функциональности: $\forall n \in \N\; \exists! x \in \R\; \ph(n, x)$. Аналогичным образом, через формулы $\psi(x, y)$ задаются конкретные числовые функции $f: \R \to \R$. И хотя мы не имеем права использовать кванторы по всем функциям ($\forall f$), мы можем формулировать и доказывать общие теоремы анализа как метаматематические \emph{схемы} (допустимые для любой определимой функции).


\smallskip\noindent\textbf{Аксиоматика непрерывности.} 
Бесконечная схема алгебраических корней заменяется на одну мощную аксиому непрерывности.

\begin{definition}[Верхняя грань и супремум]
Пусть $X \subseteq \R$ --- непустое множество.
\begin{enumerate}
    \item \textbf{Верхняя грань} для $X$ --- число $u \in \R$ такое, что $\forall x \in X: x \le u$ (кратко записывают $X \le u$).
    \item Множество $X$ \textbf{ограничено сверху}, если у него есть хотя бы одна верхняя грань.
    \item \textbf{Супремум} (точная верхняя грань): $\sup X = \min\{u \in \R \mid X \le u\}$, если этот минимум существует.
\end{enumerate}
\end{definition}

\begin{axiom}[Непрерывность / полнота]
Всякое непустое подмножество $\R$, ограниченное сверху, имеет супремум в $\R$. Формульно:
\[
\forall X \subseteq \R\; \bigl( X \neq \emptyset \wedge \exists u \in \R\; (X \le u) \bigr) \Rightarrow \exists \sup X.
\]
\end{axiom}

\begin{theorem}[Выводимость вещественной замкнутости]
    Схема существования корней для многочленов нечётной степени \ref{ax:rcf-odd-root} и аксиома существования квадратного корня \ref{ax:square-root} в предикативном анализе строго выводятся как теоремы из аксиом упорядоченного поля и аксиомы непрерывности.
\end{theorem}

\noindent\textit{Пояснение:} квадратный корень из $a \ge 0$ строится как $\sup \{x \in \R \mid x^2 \le a\}$, а корни нечётных многочленов гарантируются теоремой Больцано\,---\,Коши о промежуточном значении непрерывной функции, которая также доказывается с использованием точных граней.

\begin{remark}[Связь с арифметикой]
Наша теория вещественных чисел на уровне II позволяет интерпретировать арифметику $\mathsf{Z}_2$ --- второпорядковую арифметику Пеано с полной схемой свертывания и полной индукцией. Это следует из того, что аксиома натурального ряда и аксиома непрерывности позволяют определять $n$-ую двоичную цифру любого вещественного числа $n$, что позволяет кодировать подмножества $\N$ вещественными числами из отрезка $[0;1]$ и, следовательно, квантифицировать по ним в схеме свертывания \ref{ax:comprehension}, а это, в свою очередь, позволяет применять индукцию ко всем второпорядковым формулам $\mathsf{Z}_2$.

Интерпретация арифметики $\mathsf{Z}_2$ в нашей теории вещественных чисел делает ее геделевой теорией и, как следствие, обеспечивает ее дедуктивную неполноту и невыразимость истины (теорема Тарского для второго порядка).

Здесь необходимо прояснить связь нашей формализации с \emph{обратной математикой}. В логике Анализ строят иначе: там рассматривают арифметику второго порядка, где индивидами (тип 0) выступает $\N$, а множествами (тип 1) --- элементы булеана $\Pcal(\N)$. Вещественные числа при этом конструируются как множества натуральных (последовательности Коши), поэтому континуум $\R$ синтаксически отождествляется с $\Pcal(\N)$. Внутри этой парадигмы выделяют иерархию подсистем:

\marginfigure[1][0.75]{
\begin{minipage}{10.5cm}
\small        
    \textbf{Теорема Больцано\,---\,Коши (о промежуточном значении):}\\
    Пусть функция $f(x)$ непрерывна на отрезке $[a, b]$ и принимает на его концах значения разных знаков ($f(a)f(b) < 0$). Тогда существует точка $c \in (a, b)$ такая, что $f(c) = 0$.\\
    \textit{Связь с супремумом:} На Уровне II эта точка строится так: если $f(a) < 0$, то корень --- это в точности $c = \sup \{ x \in [a,b] \mid f(x) < 0 \}$.
\end{minipage}
}
\begin{itemize}
    \item Базовая система $\mathsf{RCA}_0$ (рекурсивное свёртывание) достаточна для доказательства архимедовости $\R$.
    \item Система $\mathsf{WKL}_0$ (слабая лемма Кёнига) доказывает теоремы о компактности: лемму Гейне\,---\,Бореля, теорему Вейерштрасса о достижении максимума, теорему Кантора о равномерной непрерывности и теорему Брауэра о неподвижной точке.
    \item Предикативная система $\mathsf{ACA}_0$ (схема свёртывания без кванторов по $\Pcal(\N)$) строго сильнее предыдущих. Она позволяет доказать критерий Коши, теорему Больцано\,---\,Вейерштрасса о сходящейся подпоследовательности.
    \item Арифметика $\mathsf{Z}_2$ (полная схема свёртывания по $\Pcal(\N)$) исчерпывает всю логику работы с самими вещественными числами и непрерывными функциями.
\end{itemize}
Арифметическая теория $Z_2$ уже является достаточно мощной, чтобы доказать 99\% теорем классического анализа (пределы, производные, интеграл Римана, теоремы Вейерштрасса и т.д.). 

Наша теория Анализа на уровне II по сути представляет собой синтаксический сахар над теорией $\mathsf{Z}_2$. Она позволяет записать многие утверждения короче, но не доказывает никаких новых теорем, которые нельзя было бы доказать в теории $\mathsf{Z}_2$.
\end{remark}

Мы можем усилить нашу теорию, взяв стандартную модель $\R$ (построенную из $\Q$, например, с помощью дедекиндовых сечений) и ее элементарную второпорядковую теорию $Th^{2}(\R)$. Такая теория является полной по построению и содержит новые факты о подмножествах $\R$, недоступные в $Z_2$. При этом она является неперечислимой, как и полная арифметика $Th(\N)$.

Но даже теория $Th^{2}(\R)$ неспособна описывать такие сложные вещи, как произвольное семейство подмножеств пространства $C[a;b]$, топологию этого пространства как объект, произвольный функционал на этом пространстве, а также пространства, получаемые через ультрафильтры и т.д. В современной обратной математике используется специальная шкала арифметик высших порядков $Z_n$, где каждая следующая делает доступными все более сложные конструкции, но все они оказываются слабее теории $\ZFC$. Поэтому окончательный и самый мощный вариант Анализа строится именно на базе $\ZFC$.


\subsubsection{Уровень III: Матанализ без ограничений (\ZFA/\ZFC)}


Для работы с произвольными функционалами, несепарабельными пространствами (типа $L^\infty$) и абстрактными топологиями двусортного языка недостаточно. Здесь математики вынуждены переходить к полному языку теории множеств.

\smallskip\noindent\textbf{Язык и Архитектура.} 
В классических учебниках (Рудин, Зорич) анализ строится в \ZFC{}, где математики негласно договариваются не применять операцию $\in$ к самим числам из соображений «хорошего тона». 
Мы работаем внутри $\mathsf{ZF}$, но весь язык редуцирован к кумулятивной иерархии $V(\R)$, стартующей с $\R$ (по той же схеме, что $\emptyset$ порождает стандартный универсум фон Неймана).

На базовом уровне иерархии элементы $\R$ рассматриваются как урэлементы (атомы). Их свойства полностью исчерпываются элементарной теорией модели $\R$ в сигнатуре упорядоченного поля. Это аппаратно исключает из рассмотрения внутреннюю $\in$-структуру самих чисел. При этом над базой $\R$ становятся доступны все мощные конструкции $\mathsf{ZF}$: ординалы, кардиналы, трансфинитная индукция и пространства произвольных функций и топологий.

\begin{remark}[Роль аксиомы выбора]
Виновником появления неконструктивных «патологий» (существование неизмеримого множества Витали, парадокса Банаха\,---\,Тарского), равно как и необходимым инструментом для доказательства большинства теорем функционального анализа, является Аксиома Выбора (\AC). Если в \ZFC{} стандартом является полная \AC{}, то для нужд анализа часто достаточно её ослабленных вариантов:
\begin{itemize}
    \item \textbf{Счётный выбор ($\AC_\omega$):} Выбор по счётным семействам. Достаточен для ключевых шагов построения меры Лебега на прямой без полной аксиомы выбора. Чаще даже рассматривают $\AC_\om(\R)$, т.е. счетный выбор из подмножеств $\R$.
    \item \textbf{Зависимый выбор ($\mathsf{DC}$):} Позволяет строить последовательности и достаточен для большинства теорем функционального анализа (например, теоремы Бэра о категориях и теоремы Хана\,---\,Банаха для сепарабельных пространств).
    \item \textbf{Полная Аксиома Выбора ($\AC$):} Необходима для теоремы Тихонова о произведении компактов, теоремы Хана\,---\,Банаха в несепарабельном случае, а также является виновником неконструктивных «патологий» (существование неизмеримого множества Витали, парадокса Банаха\,---\,Тарского).
\end{itemize}
\end{remark}


\begin{remark}[Инкапсуляция и независимость от реализации]
В классической обратной математике логики часто идут другим путем: берут за основу индивидов натуральные числа $\N$, а $\R$ конструируют как множества первого сорта. Но в такой системе (и в чистой \ZFC) запись $n \in \pi$ (натуральное число принадлежит числу пи) становится синтаксически корректной формулой, а объединение $\pi \cup e$ даёт некий новый «мусорный» математический объект. 
Использование $\R$ в качестве урэлементов (подход \ZFA) работает как архитектурная инкапсуляция в программировании: оно аппаратно запрещает заглядывать внутрь чисел, защищая наш анализ от зависимости от конкретной реализации (через сечения или через последовательности Коши) и гарантируя, что теоремы анализа абсолютно универсальны.
\end{remark}


\subsubsection{Сводная таблица сравнения уровней}

\begin{table}[htbp]
% расширение влево в широкое поле (geometry: left=12cm)
\noindent\hspace*{-10cm}%
\begin{minipage}[t]{\dimexpr\textwidth+10cm\relax}
\centering
\small
\caption{Сравнение выразительной силы и метатеоретических свойств уровней формализации}
\label{tab:formalization_levels}
\begin{tabular}{@{}p{0.38\textwidth} p{0.18\textwidth} p{0.21\textwidth} p{0.15\textwidth}@{}}
\toprule
\textbf{Характеристика / Знаковые теоремы} & \textbf{Уровень I: $\RCF$} \newline (1-й порядок над $\R$) & \textbf{Уровень II: Вейль} \newline (2-й порядок над $\R$) & \textbf{Уровень III: $\ZFA$} \newline (атомы $\R$ + выбор) \\
\midrule
\multicolumn{4}{@{}l}{\textbf{Основания}} \\
\midrule
Типизация & Только числа & Числа и множества & Иерархия $V(\R)$ \\
Схема выделения подмножеств & Нет (нет объектов) & По числовым форм. & Полная (ZF) \\
Разрешимость и дедуктивная полнота& Да (Тарский) & Нет (Гёдель) & Нет (Гёдель) \\
Несчётность $\R$ (Теорема Кантора) & Нет (нет II порядка) & Да (вложен. отрезки) & Да \\
Определение множества $\N$ & Нет (о-минимальность)& Да (аксиоматически) & Да \\
\midrule
\multicolumn{4}{@{}l}{\textbf{Базовые теоремы анализа (XIX век)}} \\
\midrule
Пределы, производные, интегралы & Нет (нет II порядка) & Да & Да \\
Архимедовость $\R$ & Нет (нет $\N \subset \R$) & Да (в $\mathsf{RCA}_0$) & Да \\
Теорема Больцано\,---\,Коши & Да (многочлены) & Да ($\equiv \mathsf{WKL}_0$) & Да \\
Лемма Гейне\,---\,Бореля & Нет (нет II порядка) & Да ($\equiv \mathsf{WKL}_0$) & Да \\
Теорема Вейерштрасса об экстремуме & Нет (нет II порядка) & Да ($\equiv \mathsf{WKL}_0$) & Да \\
Теорема Брауэра (неподвижная точка) & Нет (нет II порядка) & Да ($\equiv \mathsf{WKL}_0$) & Да \\
Теор. Кантора (равном. непрерывность) & Нет (нет II порядка) & Да ($\equiv \mathsf{WKL}_0$) & Да \\
Теорема Больцано\,---\,Вейерштрасса & Нет (нет II порядка) & Да ($\equiv \mathsf{ACA}_0$) & Да \\
Критерий Коши (сходимость в $\R$) & Нет (нет II порядка) & Да ($\equiv \mathsf{ACA}_0$) & Да \\
\midrule
\multicolumn{4}{@{}l}{\textbf{Мера, интеграл, функциональный анализ (XX век)}} \\
\midrule
Мера и Интеграл Лебега & Нет (нет $\sigma$-алгебр) & Ограниченно (коды) & Да \\
Пространства $C[a,b]$ и $L^p$ ($p<\infty$) & Нет (нет $\sigma$-алгебр) & Да (через коды) & Да \\
Пространство $L^\infty$ и функционалы & Нет (нет $\sigma$-алгебр) & Нет (нужен 3-й сорт) & Да ($\Pcal(\R\times\R)$) \\
Теорема Хана\,---\,Банаха & Нет (нет $\sigma$-алгебр) & Нет (нужна $\AC$) & Да ($\mathsf{DC}$ / $\AC$) \\
\midrule
\multicolumn{4}{@{}l}{\textbf{Следствия аксиомы выбора}} \\
\midrule
Теорема Тихонова (произвед. компактов)& Нет (нет II порядка) & Нет (полная $\AC$) & Да (полная $\AC$) \\
Множество Витали (неизмеримость) & Нет (нет II порядка) & Нет (полная $\AC$) & Да (полная $\AC$) \\
Парадокс Банаха\,---\,Тарского & Нет (нет II порядка) & Нет (полная $\AC$) & Да (полная $\AC$) \\
\bottomrule
\end{tabular}
\end{minipage}
\end{table}


\subsubsection{Расширенная числовая прямая}

Для удобства формулировок многих теорем анализа (особенно касающихся пределов и точных граней) к множеству действительных чисел добавляют два формальных символа: $-\infty$ и $+\infty$. Полученное множество называется \textbf{расширенной числовой прямой}:
\[ \overline{\R} = \R \cup \{-\infty, +\infty\}. \]

\begin{remark}[Связь с уровнями формализации]
Важно понимать, что $\overline{\R}$ \emph{не является полем} (в нём нарушаются свойства групповых операций). Поэтому на алгебраическом Уровне~I (\RCF) добавление этих символов бессмысленно. Конструкция $\overline{\R}$ становится актуальной только начиная с Уровня~II, где в языке появляются пределы и точные грани (в частности, в $\overline{\R}$ абсолютно \emph{любое} подмножество имеет супремум).
\end{remark}

\smallskip\noindent\textbf{Порядок и Операции.}
Линейный порядок с $\R$ продолжается на $\overline{\R}$ естественным образом: $-\infty < x < +\infty$.
Там, где это не ведёт к логическим противоречиям, арифметические операции доопределяются (в духе теории пределов). Например: $x + (+\infty) = +\infty$, а при $x < 0$: $x \cdot (+\infty) = -\infty$. 
Операции, которые невозможно доопределить однозначно, оставляют \textbf{не определёнными} (неопределённости):
\[ (+\infty) - (+\infty), \quad (-\infty) - (-\infty), \quad 0 \cdot (\pm\infty), \quad \frac{\pm\infty}{\pm\infty}, \quad \frac{0}{0}. \]

\subsubsection{Фундаментальные следствия непрерывности (Уровень II)}

В этом разделе приведены ключевые факты, которые опираются на предикативную базу Уровня II (наличие $\N$ и аксиому супремума). 

\begin{theorem}[Свойство Архимеда]
В любом непрерывно упорядоченном поле множество $\N$ не ограничено сверху: для любого $x$ существует $n \in \N$ такое, что $n > x$.
\end{theorem}
\noindent\textit{Доказательство:} если бы $\N$ было ограничено сверху, по аксиоме непрерывности существовало бы $s = \sup \N$; тогда $s - 1$ не была бы верхней гранью, то есть для некоторого $n \in \N$ выполнено $n > s - 1$, откуда $n + 1 > s$. Но $n+1 \in \N$ (по аксиоме индуктивности), что противоречит тому, что $s$ --- верхняя грань.

\begin{theorem}[Эквивалентность: супремум $\Leftrightarrow$ инфимум]
Всякое непустое подмножество $\R$, ограниченное сверху, имеет супремум тогда и только тогда, когда всякое непустое подмножество $\R$, ограниченное снизу, имеет инфимум.
\end{theorem}
\noindent\textit{Доказательство:} $\inf X = -\sup(-X)$ и наоборот.

\marginfigure[1][0.75]{
\begin{minipage}{10.5cm}
\small
\textbf{Базовые определения (Уровень II/III):}
\begin{itemize}[leftmargin=*, nosep]
    \item \textbf{Предел последовательности:} $\lim a_n = a$ означает $\forall \ep > 0\; \exists N \in \N\; \forall n \ge N: |a_n - a| < \ep$.
    \item \textbf{Фундаментальная последовательность:} $\langle a_n \rangle$ фундаментальна (по Коши), если $\forall \ep > 0\; \exists N \in \N\; \forall m,n \ge N: |a_m - a_n| < \ep$.
    \item \textbf{Предел функции:} $\lim_{x \to a} f(x) = L$ означает: $\forall \ep > 0\; \exists \de > 0\; \forall x\ (0 < |x - a| < \de \Rightarrow |f(x) - L| < \ep)$. Условие $0 < |x - a|$ исключает саму точку $a$, чтобы предел описывал поведение $f$ \emph{вблизи} $a$.
\end{itemize}
\end{minipage}
}

\begin{theorem}[Эквивалентные формулировки непрерывности]
В рамках аксиом упорядоченного поля принцип непрерывности (аксиома супремума) равносилен каждому из следующих утверждений:
\begin{enumerate}
    \item \textbf{Свойство Архимеда} + \textbf{принцип вложенных отрезков}: для любой последовательности замкнутых отрезков $[a_n, b_n]$ с $a_n \le b_n$ и $[a_{n+1}, b_{n+1}] \subseteq [a_n, b_n]$ пересечение $\bigcap_{n=1}^\infty [a_n, b_n]$ непусто.
    \item \textbf{Свойство Архимеда} + \textbf{полнота по Коши}: всякая фундаментальная последовательность имеет предел в $\R$.
    \item \textbf{Свойство Архимеда} + \textbf{сходимость монотонных ограниченных последовательностей}: всякая неубывающая последовательность, ограниченная сверху, сходится к своему супремуму.
\end{enumerate}
\end{theorem}
\noindent\textit{Пояснение:} Эквивалентности с пунктами 2 и 3 жестко опираются на свойство Архимеда; без него существуют упорядоченные поля (с бесконечно малыми), полные по Коши, но не удовлетворяющие аксиоме супремума.


\begin{theorem}[Категоричность непрерывно упорядоченного поля]
С точностью до изоморфизма существует ровно одно упорядоченное поле, удовлетворяющее аксиоме непрерывности (супремума) --- это поле вещественных чисел $\R$.
\end{theorem}
    

\begin{remark}[Метатеоретический статус и уникальность $\R$]
Следует оговориться, что сама теорема о категоричности выходит за рамки выразительных возможностей Уровня II. Доказательство существования изоморфизма между любыми двумя полями требует конструирования биекции (функции, связывающей две абстрактные структуры) --- объекта, недоступного в двусортной арифметике. Поэтому этот факт строго формулируется и доказывается лишь на Уровне III (в $\ZFC$) или в рамках семантической теории моделей.
    
Тем не менее, эта теорема делает континуум $\R$ абсолютно уникальным, безальтернативным фундаментом анализа. Любое строгое расширение $\R$ (например, поле гипервещественных чисел в нестандартном анализе или поля сюрреальных чисел Конвея) неминуемо содержит бесконечно большие элементы, то есть теряет свойство Архимеда. Но, как мы доказали выше, любое непрерывное (полное по Дедекинду) поле обязано быть архимедовым, иначе натуральный ряд $\N$ оказался бы ограничен сверху, не имея при этом точной верхней грани. Следовательно, любое расширение $\R$ неизбежно оказывается «дырявым» в смысле сечений. Аксиома супремума жёстко фиксирует размер континуума, одновременно запрещая как лакуны, так и бесконечно малые величины.

Казалось бы, если мы продолжим трансфинитный процесс пополнения поля (как это делается, например, при построении иерархий в теории множеств или сюрреальных чисел по кардиналам), то рано или поздно должны получить гигантское поле, где снова исчезнут «дыры». Однако этот алгебраический барьер непреодолим: добавление новых элементов неизбежно нарушает свойство Архимеда, порождая бесконечно большие и бесконечно малые величины. А значит, стандартный континуум $\R$ навсегда останется ограниченным в этом расширении, но не будет иметь точной верхней грани. Любое упорядоченное поле (включая поле-класс \textbf{No} всех сюрреальных чисел), не изоморфное $\R$, обречено на наличие дедекиндовых сечений, которые нельзя заполнить. Любое строгое расширение $\R$ неизбежно оказывается «дырявым» в смысле сечений.
\end{remark}

\begin{example}[Физико-философский комментарий: Космологические горизонты и неархимедовость]
    Математическая неизбежность появления разрывов в расширениях $\R$ имеет глубокий смысл для теоретической физики. Когда физики пытаются описать космологический горизонт (границы наблюдаемой Вселенной), сингулярности чёрных дыр или природу тёмной энергии, они иногда обращаются к моделям пространства-времени, основанным на неархимедовых полях. 
    
    В стандартном континууме $\R$ масштаб Вселенной описывается как потенциальная бесконечность. Введение неархимедовых структур (например, гипервещественных или сюрреальных чисел) позволяет математически строго оперировать актуально бесконечно большими расстояниями (моделируя «то, что находится за космологическим горизонтом») и бесконечно малыми (для описания структуры «квантовой пены» на планковских масштабах). 
    
    Однако теорема о категоричности ставит жёсткий метафизический ультиматум: расплатой за возможность работать с актуальными бесконечностями становится потеря непрерывности. Пространство-время в неархимедовых теориях гравитации неизбежно оказывается «дырявым» в смысле Дедекинда --- в нём пропадают точные границы (супремумы), отделяющие макроскопический (архимедов) мир от бесконечно большого или инфинитезимального. Это математическое свойство изящно согласуется с физическими проблемами при попытках «склеить» общую теорию относительности (макро-мир непрерывного $\R$) и квантовую механику (микро-мир).
\end{example}




\subsubsection{Простые следствия аксиом}


Эти факты удобно иметь «под рукой», чтобы не опираться на интуицию и не испытывать неловкости при формальных рассуждениях.

\begin{enumerate}
    \item \textbf{Теорема о делении с остатком:} Для любых $\alpha \in \R$ и $d > 0$ существуют единственное целое $k$ и единственное $r \in [0; d)$ такие, что $\alpha = kd + r$.
    \item \textbf{Плотность $\Q$ в $\R$:} Между любыми двумя вещественными $a < b$ существует рациональное: $\exists q \in \Q: a < q < b$. \textit{(Вывод: взять $n \in \N$ с $n(b-a)>1$, применить деление с остатком к $na$ при $d=1$, положить $q=(k+1)/n$, где $na = k + r$.)}
    \item \textbf{Модуль и неравенство треугольника:} $|x| = \max\{x, -x\}$; для всех $x, y \in \R$ выполнено $|x + y| \le |x| + |y|$ и $|xy| = |x| \cdot |y|$. \textbf{Обратное неравенство треугольника:} $\bigl| |x| - |y| \bigr| \le |x - y|$ (следствие: $|x| - |y| \le |x - y|$).
    \item \textbf{Неравенство Бернулли:} Для $x \ge -1$ и $n \in \N$: $(1+x)^n \ge 1 + nx$.
    \item \textbf{Оценки для дробей:} Если $0 < a \le b$, то $\frac{a}{b} \le \frac{a+c}{b+c}$ при $c > 0$; полезные оценки при работе с пределами.
    \item \textbf{Единственность предела:} Последовательность не может иметь двух различных пределов.
    \item \textbf{Ограниченность сходящейся последовательности:} Всякая сходящаяся последовательность в $\R$ ограничена.
    \item \textbf{Теорема Больцано--Вейерштрасса (для последовательностей):} Если последовательность в $\R$ ограничена, то у неё есть сходящаяся подпоследовательность. Предел такой подпоследовательности называется \textbf{частичным пределом} исходной последовательности.
    \item \textbf{Арифметика пределов последовательностей:} Если $a_n \to a$ и $b_n \to b$, то $a_n + b_n \to a + b$, $a_n - b_n \to a - b$, $a_n b_n \to ab$; если кроме того $b \neq 0$, то $a_n/b_n \to a/b$ (для частного считаем $b_n \neq 0$ при всех $n$). Аналогично для пределов функций при $x \to a$.
    \item \textbf{Эквивалентность определений предела функции (Коши и Гейне):} $\lim_{x \to a} f(x) = L$ по Коши (через $\ep$-$\de$) равносильно тому, что для любой последовательности $x_n \to a$, $x_n \neq a$, выполнено $f(x_n) \to L$. То же для односторонних пределов и для $x \to \infty$.
    \item \textbf{$O$- и $o$-символика (Ландау):} Пусть $f,g$ заданы в проколотой окрестности точки $a$ (или при $n \to \infty$). Пишут $f = O(g)$ при $x \to a$ (соответственно при $n \to \infty$), если $|f| \le C|g|$ в некоторой окрестности $a$ для некоторой константы $C$. Пишут $f = o(g)$, если $\frac{f}{g} \to 0$ при $x \to a$ (или при $n \to \infty$). Асимптотическое равенство: $f \sim g$ означает $f/g \to 1$.
\end{enumerate}

\marginfigure[-1][0.75]{
\begin{minipage}{10.5cm}
\small
\textbf{Верхний и нижний предел последовательности.}
\begin{itemize}[leftmargin=*, nosep]
    \item \textbf{Верхний предел:} $\varlimsup_{n \to \infty} a_n = \overline{\lim}_{n \to \infty} a_n = \lim_{n \to \infty} \sup_{k \ge n} a_k$ (предел монотонно неубывающей последовательности $\sup_{k \ge n} a_k$; в $\overline{\R}$ всегда существует). Равноправно пишут и $\limsup$.
    \item \textbf{Нижний предел:} $\varliminf_{n \to \infty} a_n = \underline{\lim}_{n \to \infty} a_n = \lim_{n \to \infty} \inf_{k \ge n} a_k$. Равноправно и $\liminf$.
    \item У \textbf{ограниченной} последовательности множество частичных пределов непусто (Больцано--Вейерштрасса), и его \textbf{максимум} совпадает с $\varlimsup a_n$, \textbf{минимум} --- с $\varliminf a_n$.
\end{itemize}
\end{minipage}
}

\subsection{Основные теоремы анализа с примерами из физики}

\marginfigure[1][0.75]{
\begin{minipage}{10.5cm}
\small
\textbf{Окрестности и формулировки на их языке.}
\begin{itemize}[leftmargin=*, nosep]
    \item \textbf{$\ep$-окрестность} точки $a \in \R$: $U_\ep(a) = \{x : |x - a| < \ep\}$ (интервал $(a-\ep, a+\ep)$).
    \item \textbf{Проколотая $\ep$-окрестность}: $\dot{U}_\ep(a) = U_\ep(a) \setminus \{a\} = \{x : 0 < |x - a| < \ep\}$.
    \item \textbf{Непрерывность в точке (на языке окрестностей):} $f$ непрерывна в $a$ тогда и только тогда, когда для любой окрестности $V$ точки $f(a)$ существует окрестность $U$ точки $a$ такая, что $f(U) \subseteq V$ (т.\,е. для любого $\ep > 0$ найдётся $\de > 0$ с $f(U_\de(a)) \subseteq U_\ep(f(a))$).
    \item \textbf{Открытое множество:} $A \subseteq \R$ открыто, если для любой точки $a \in A$ существует $\ep > 0$ такое, что $U_\ep(a) \subseteq A$ (каждая точка входит в множество вместе с какой-то своей окрестностью).
    \item \textbf{Замкнутое множество:} $A \subseteq \R$ замкнуто, если его дополнение $\R \setminus A$ открыто; эквивалентно: $A$ содержит все свои предельные точки (т.\,е. $\overline{A} = A$).
    \item \textbf{Компакт:} в $\R$ (и в $\R^n$) подмножество \textbf{компактно} тогда и только тогда, когда оно замкнуто и ограничено (теорема Гейне--Бореля).
\end{itemize}
\medskip
\textbf{Плотность и изолированная точка.}
\begin{itemize}[leftmargin=*, nosep]
    \item Подмножество $A \subseteq \R$ \textbf{плотно в $\R$} (или плотно на прямой), если замыкание $\overline{A}$ совпадает с $\R$, т.е. любая точка $x \in \R$ является пределом некоторой последовательности элементов из $A$. Эквивалентно: любой непустой открытый интервал $(a,b)$ пересекается с $A$: $(a,b) \cap A \neq \emptyset$.
    \item Точка $a \in A \subseteq \R$ называется \textbf{изолированной точкой} множества $A$, если существует $\ep > 0$ такое, что $U_\ep(a) \cap A = \{a\}$ (в некоторой окрестности $a$ нет других точек из $A$). Иначе говоря, $a$ не является предельной точкой $A$.
\end{itemize}
\end{minipage}
}

\begin{definition}[Непрерывность функции]
    Функция $f: \R \to \R$ непрерывна в точке $a \in \R$, если существует предел $\lim_{x \to a} f(x)$ и он равен $f(a)$:
    $$\lim_{x \to a} f(x) = f(a).$$
    Функция непрерывна на множестве, если она непрерывна в каждой точке этого множества.
    \end{definition}
    
    
\begin{example}[Физический пример: траектория частицы]
Пусть $x(t)$ --- координата частицы в момент времени $t$. Непрерывность означает, что малые изменения времени приводят к малым изменениям положения. Разрыв соответствовал бы «телепортации» частицы.
\end{example}


\begin{theorem}[Теорема о промежуточном значении (Больцано--Коши)]
Если функция $f: [a, b] \to \R$ непрерывна на $[a, b]$ и $f(a) < c < f(b)$ (или $f(b) < c < f(a)$), то существует $x \in (a, b)$ такое, что $f(x) = c$.
\end{theorem}

\begin{example}[Физический пример: поиск равновесия]
Если температура тела изменяется непрерывно от $T_1$ до $T_2$, то в некоторый момент она принимает любое промежуточное значение. Это используется в термодинамике для поиска состояний равновесия.
\end{example}

\begin{theorem}[Теорема Вейерштрасса об экстремумах]
Непрерывная функция $f: [a, b] \to \R$ на замкнутом отрезке достигает своего максимума и минимума.
\end{theorem}
(Обычно доказывается напрямую: например, $s = \sup f([a,b])$ существует по полноте $\R$; берём последовательность $x_n \in [a,b]$ с $f(x_n) \to s$, по теореме Больцано--Вейерштрасса о последовательностях у $x_n$ есть сходящаяся подпоследовательность $x_{n_k} \to c \in [a,b]$, по непрерывности $f(c) = s$.)

\begin{theorem}[Теорема Вейерштрасса: непрерывный образ компакта]
Непрерывная функция отображает компакт в компакт: если $K \subseteq \R^n$ компактно и $f: K \to \R^m$ непрерывна, то $f(K)$ компактно в $\R^m$. (Обобщение теоремы об экстремумах: в $\R$ компакт --- это замкнутое ограниченное множество, так что из этой теоремы снова следует достижимость $\max$ и $\min$ на отрезке.)
\end{theorem}


\marginfigure[1][0.75]{
\begin{minipage}{10.5cm}
\small
\textbf{Аксиомы топологии для открытых множеств на $\R$.}

Семейство $\tau$ подмножеств множества $X$ называется \textbf{топологией}, если выполнены:
\begin{enumerate}[label=(\roman*), nosep]
    \item $\emptyset \in \tau$ и $X \in \tau$;
    \item объединение любого набора множеств из $\tau$ лежит в $\tau$;
    \item пересечение любых \emph{двух} (и значит, любого конечного числа) множеств из $\tau$ лежит в $\tau$.
\end{enumerate}
Открытые подмножества прямой по определению: $G \subseteq \R$ открыто, если для каждой точки $x \in G$ найдётся интервал $(x-\ep, x+\ep) \subseteq G$.

\textit{Проверка:} (i) Пустое множество открыто (условие выполняется «викуумно»); $\R$ открыто (любой интервал содержится в $\R$). (ii) Если $G_i \in \tau$ и $x \in \bigcup_i G_i$, то $x \in G_j$ для некоторого $j$, значит есть окрестность $x$ в $G_j$, она же окрестность в объединении. (iii) Если $G_1, G_2$ открыты и $x \in G_1 \cap G_2$, то есть $\ep_1, \ep_2 > 0$ с $U_{\ep_1}(x) \subseteq G_1$ и $U_{\ep_2}(x) \subseteq G_2$; тогда при $\ep = \min(\ep_1, \ep_2)$ получим $U_\ep(x) \subseteq G_1 \cap G_2$.

\medskip
\textbf{Евклидова топология в $\R^n$.} Стандартная топология задаётся евклидовой метрикой и порождёнными ею понятиями:
\begin{itemize}[leftmargin=*, nosep]
    \item \textbf{Метрика:} $d(x,y) = |x - y| = \sqrt{(x_1-y_1)^2 + \ldots + (x_n-y_n)^2}$ для $x, y \in \R^n$.
    \item \textbf{$\ep$-окрестность} точки $a \in \R^n$: $U_\ep(a) = \{ x \in \R^n : d(x,a) < \ep \}$ (открытый шар радиуса $\ep$).
    \item \textbf{Открытое множество:} $A \subseteq \R^n$ открыто, если для любой точки $a \in A$ существует $\ep > 0$ с $U_\ep(a) \subseteq A$.
    \item \textbf{Замкнутое множество:} $A$ замкнуто, если $\R^n \setminus A$ открыто; эквивалентно: $A$ содержит все свои предельные точки.
    \item \textbf{Компакт:} в $\R^n$ по теореме Гейне--Бореля подмножество компактно $\Leftrightarrow$ оно замкнуто и ограничено (в метрике $d$).
\end{itemize}
\end{minipage}
}


\begin{example}[Физический пример: поиск оптимальной траектории]
В механике Лагранжа действие $S = \int L \, dt$ является функционалом на пространстве траекторий. Теорема Вейерштрасса гарантирует существование траектории с минимальным действием (принцип наименьшего действия).
\end{example}

\begin{definition}[Равномерная непрерывность]
Функция $f: E \subseteq \R^n \to \R^m$ \textbf{равномерно непрерывна} на $E$, если
$$\forall \ep > 0\; \exists \de > 0\; \forall x, y \in E\ \bigl( |x - y| < \de \Rightarrow |f(x) - f(y)| < \ep \bigr).$$
В отличие от обычной непрерывности, одно и то же $\de$ работает для всех пар точек из $E$.
\end{definition}

\begin{remark}[Преемственность терминов: равномерная непрерывность $\Rightarrow$ непрерывность]
Равномерно непрерывная функция непрерывна в каждой точке. Достаточно применить общезначимую импликацию $\exists x\; \forall y\; \ph(x,y) \to \forall y\; \exists x\; \ph(x,y)$: при фиксированном $\ep > 0$ равномерная непрерывность даёт «$\exists \de\; \forall x,y \in E\ \bigl( |x-y| < \de \Rightarrow |f(x)-f(y)| < \ep \bigr)$»; по этой импликации для каждой пары $(a,y)$ (в частности при фиксированном $a$) существует $\de$ с $|y - a| < \de \Rightarrow |f(y) - f(a)| < \ep$, т.\,е. выполнено определение непрерывности $f$ в точке $a$.
\end{remark}

\begin{theorem}[Кантор: непрерывность на компакте $\Rightarrow$ равномерная непрерывность]
Если множество $K \subseteq \R^n$ компактно и функция $f: K \to \R^m$ непрерывна на $K$, то $f$ равномерно непрерывна на $K$.
\end{theorem}
(Идея доказательства: от противного; тогда для некоторого $\ep > 0$ при каждом $n$ найдутся $x_n, y_n \in K$ с $|x_n - y_n| < 1/n$ и $|f(x_n) - f(y_n)| \ge \ep$; по компактности у $\langle x_n \rangle$ есть сходящаяся подпоследовательность $x_{n_k} \to c \in K$; тогда и $y_{n_k} \to c$; по непрерывности $f(x_{n_k}) \to f(c)$ и $f(y_{n_k}) \to f(c)$, противоречие с $|f(x_{n_k}) - f(y_{n_k})| \ge \ep$.)

\begin{theorem}[Теорема Брауэра о неподвижной точке]
Любое непрерывное отображение выпуклого компакта в себя имеет по крайней мере одну неподвижную точку. В частности, для любого непрерывного $f: [a, b] \to [a, b]$ существует точка $x \in [a, b]$, для которой $f(x) = x$.
\end{theorem}
    
\begin{example}[Физический пример: перемешивание жидкости]
Если взять стакан воды и аккуратно перемешать его содержимое (непрерывной деформацией без разрывов сплошности), то после остановки жидкости хотя бы одна молекула воды окажется ровно в той же точке пространства, где она находилась до начала перемешивания.
\end{example}


\begin{definition}[Производная]
Производная функции $f$ в точке $a$ определяется как предел:
$$f'(a) = \lim_{h \to 0} \frac{f(a + h) - f(a)}{h}$$
если этот предел существует.
\end{definition}

\begin{example}[Физический пример: скорость и ускорение]
Если $x(t)$ --- координата частицы, то $v(t) = x'(t)$ --- скорость, а $a(t) = v'(t) = x''(t)$ --- ускорение. Это основа кинематики.
\end{example}

\begin{theorem}[Теорема Ролля]
Если функция $f: [a, b] \to \R$ непрерывна на $[a, b]$, дифференцируема на $(a, b)$ и $f(a) = f(b)$, то существует $c \in (a, b)$ такое, что $f'(c) = 0$.
\end{theorem}

\begin{theorem}[Теорема Лагранжа (о среднем значении)]
Если функция $f: [a, b] \to \R$ непрерывна на $[a, b]$ и дифференцируема на $(a, b)$, то существует $c \in (a, b)$ такое, что:
$$f'(c) = \frac{f(b) - f(a)}{b - a}$$
\end{theorem}

\begin{example}[Физический пример: средняя скорость]
Если автомобиль проехал расстояние $s$ за время $t$, то его средняя скорость равна $\frac{s}{t}$. Теорема Лагранжа гарантирует, что в некоторый момент мгновенная скорость совпадает со средней.
\end{example}

\begin{theorem}[Правило Лопиталя]
Если $\lim_{x \to a} f(x) = \lim_{x \to a} g(x) = 0$ (или оба равны $\infty$) и существует $\lim_{x \to a} \frac{f'(x)}{g'(x)}$, то:
$$\lim_{x \to a} \frac{f(x)}{g(x)} = \lim_{x \to a} \frac{f'(x)}{g'(x)}$$
\end{theorem}

\begin{example}[Физический пример: анализ сингулярностей]
В физике часто возникают неопределенности вида $\frac{0}{0}$ или $\frac{\infty}{\infty}$ при анализе предельных случаев. Правило Лопиталя позволяет их разрешить.
\end{example}

\begin{theorem}[Формула Тейлора]
Если функция $f$ имеет $n+1$ производную в окрестности точки $a$, то:
$$f(x) = \sum_{k=0}^n \frac{f^{(k)}(a)}{k!}(x - a)^k + R_n(x)$$
где остаточный член $R_n(x) = \frac{f^{(n+1)}(\xi)}{(n+1)!}(x - a)^{n+1}$ для некоторого $\xi$ между $a$ и $x$.
\end{theorem}

\begin{example}[Физический пример: разложение потенциала]
В физике потенциал часто разлагают в ряд Тейлора вокруг точки равновесия. Например, для малых колебаний маятника: $V(\theta) \approx V(0) + \frac{1}{2}V''(0)\theta^2$, что дает гармонический осциллятор.
\end{example}

\begin{definition}[Интеграл Римана]
Интеграл функции $f: [a, b] \to \R$ определяется как предел интегральных сумм:
$$\int_a^b f(x) \, dx = \lim_{n \to \infty} \sum_{i=1}^n f(\xi_i) \Delta x_i$$
где разбиение отрезка уточняется, а $\xi_i \in [x_{i-1}, x_i]$.
\end{definition}

\begin{theorem}[Фундаментальная теорема анализа]
Если $F'(x) = f(x)$ для всех $x \in [a, b]$, то:
$$\int_a^b f(x) \, dx = F(b) - F(a)$$
\end{theorem}

\begin{example}[Физический пример: работа силы]
Работа силы $F(x)$ при перемещении от $a$ до $b$ равна $W = \int_a^b F(x) \, dx$. Если $F = -\frac{dV}{dx}$ (консервативная сила), то $W = V(a) - V(b)$ (изменение потенциальной энергии).
\end{example}

\begin{theorem}[Теорема о замене переменной]
Если $g: [\al, \be] \to [a, b]$ дифференцируема и $f$ интегрируема, то:
$$\int_a^b f(x) \, dx = \int_{\al}^{\be} f(g(t)) g'(t) \, dt$$
\end{theorem}

\begin{example}[Физический пример: переход к полярным координатам]
При вычислении интегралов в физике часто удобно переходить к другим координатам. Например, в полярных координатах $x = r\cos\theta$, $y = r\sin\theta$ элемент площади $dx \, dy = r \, dr \, d\theta$.
\end{example}

\begin{theorem}[Интегрирование по частям]
$$\int_a^b u(x) v'(x) \, dx = u(b)v(b) - u(a)v(a) - \int_a^b u'(x) v(x) \, dx$$
\end{theorem}

\subsection{Несколько примеров классического анализа}

\begin{enumerate}[leftmargin=*, label=\textbf{Пример \arabic*.}]

\item \textbf{Неподвижная точка непрерывного отображения отрезка в себя.}

\textit{Пространство состояний:} $M = [0,1] \subseteq \R$. \textit{Процесс:} задано непрерывное отображение $f: [0,1] \to [0,1]$; в момент $n \in \N$ состояние $x_n \in [0,1]$, закон эволюции $x_{n+1} = f(x_n)$.

\textit{Задача:} Доказать, что у $f$ существует неподвижная точка: $\exists c \in [0,1]: f(c) = c$. (Достаточно теорема о промежуточном значении: рассмотреть $g(x) = f(x) - x$.)

\textit{Усиление (принцип сжатия):} Пусть $f: [0,1] \to [0,1]$ удовлетворяет условию Липшица с константой $k \in (0,1)$: $|f(x) - f(y)| \le k|x-y|$ для всех $x,y \in [0,1]$. Доказать, что неподвижная точка единственна и что для любого $x_0 \in [0,1]$ последовательность $x_{n+1} = f(x_n)$ сходится к этой точке. (Использовать: $|x_{n+1} - c| \le k|x_n - c|$ и предел.)

\textit{Пример: ветвящийся процесс Гальтона--Ватсона.} Популяция: каждая особь независимо даёт потомков, число потомков имеет распределение $(p_k)_{k\ge 0}$. Производящая функция числа потомков одной особи: $f(z) = \sum_{k=0}^\infty p_k z^k$, $z \in [0,1]$. Производящая функция числа особей через $n$ поколений (от одной предка) --- это $n$-я итерация: $f_n(z) = f(f(\ldots f(z)\ldots))$ ($n$ раз). Неподвижная точка $z_0 = f(z_0)$ --- это \textbf{вероятность вырождения} (extinction probability) популяции. Итерация $f_{n+1}(z) = f(f_n(z))$ при $z \in [0,1)$ сходится к этой неподвижной точке по \textbf{принципу сжатия}: на $[0,1]$ производная $f'(z) = \sum_{k\ge 1} k p_k z^{k-1}$ не превосходит $f'(1) = \mu$ (среднее число потомков), поэтому $|f'(z)| \le \mu$ и $f$ удовлетворяет условию Липшица с константой $\mu$. Если $\mu = f'(1) \le 1$, то $f$ --- сжатие, единственная неподвижная точка в $[0,1]$ есть $z_0 = 1$: популяция вырождается почти наверное. Если $\mu > 1$, то кроме $z=1$ есть вторая неподвижная точка $z_0 \in (0,1)$; к ней и сходится итерация при $z \in [0,1)$, причём $z_0 < 1$ --- популяция выживает с положительной вероятностью $1 - z_0$.

\item \textbf{Иррациональная ротация окружности: плотность орбиты.}

\textit{Пространство состояний:} Окружность $S^1 = \R/\Z$ (факторизация по $t \sim t+1$; элемент --- класс $[t]$). Топологически можно отождествить с $[0,1)$ с отождествлением $0$ и $1$. \textit{Процесс:} фиксирован $\al \in \R \setminus \Q$; в момент $n$ состояние $\theta_n \in S^1$, закон $\theta_{n+1} = \theta_n + \al \pmod 1$ (сложение в $\R$, затем фактор).

\textit{Задача:} Доказать, что орбита $\{\theta_n\}_{n \ge 0}$ плотна в $S^1$: для любого открытого интервала $J \subseteq S^1$ найдётся $n$ такое, что $\theta_n \in J$. (Идея: по принципу Дирихле среди $\theta_0, \ldots, \theta_N$ два попадают в один интервал длины $< 1/N$; их разность даёт «шаг» $k\al \bmod 1$, кратные которого образуют равномерную сетку.)

\item \textbf{Канторово множество: континуум меры нуль.}

Кантор строил множество $K \subseteq [0,1]$ так: из отрезка $[0,1]$ удаляют открытый средний третий интервал $(1/3, 2/3)$; из двух оставшихся отрезков снова удаляют средние трети; и так далее. Множество $K$ --- пересечение всех оставшихся замкнутых множеств (на шаге $n$ остаётся объединение $2^n$ замкнутых отрезков длины $3^{-n}$ каждый).

\textit{Свойства, которые изучал Кантор:}
\begin{itemize}[nosep]
    \item \textbf{Мощность континуума:} $K$ равномощно $[0,1]$. Достаточно сопоставить точке $x \in K$ её троичное разложение, в котором встречаются только цифры $0$ и $2$; такое разложение единственно. Множество таких последовательностей $\{0,2\}^{\N}$ равномощно $\{0,1\}^{\N}$, а последнее уже биективно $[0,1]$ (двоичное разложение).
    \item \textbf{Мера нуль:} В данном случае «меру» можно посчитать без общей теории Лебега. На шаге $1$ удалили один интервал длины $1/3$; на шаге $2$ --- два интервала по $1/9$, итого $2/9$; на шаге $n$ удалено $2^{n-1}$ интервалов длины $3^{-n}$ каждый. Сумма длин \emph{всех} удалённых (непересекающихся) интервалов:
    \[
    \frac{1}{3} + \frac{2}{9} + \frac{4}{27} + \dots = \sum_{n \ge 1} \frac{2^{n-1}}{3^n} = \frac{1}{3} \cdot \frac{1}{1 - 2/3} = 1.
    \]
    Исходный отрезок имел длину $1$; мы удалили непересекающиеся интервалы суммарной длины $1$. Оставшееся множество $K$ не содержит ни одного интервала (в любой окрестности точки из $K$ есть удалённая точка). Естественно приписать $K$ \textbf{длину (меру) нуль}: $\mu(K) = 1 - 1 = 0$.
\end{itemize}
Итого: $K$ «столь же велико», что и отрезок (по мощности), но «нулевой длины» (по мере). Такую конструкцию Кантор использовал в исследованиях континуума и размерности.

\item \textbf{Иррациональность $\sqrt{2}$ и расширение поля $\Q$.}

\textit{Иррациональность:} Предположим $\sqrt{2} \in \Q$; тогда $\sqrt{2} = p/q$ для некоторых $p, q \in \N$ (дробь можно считать несократимой). Тогда $p^2 = 2q^2$, значит $p^2$ чётно, поэтому $p$ чётно: $p = 2m$. Подставляя: $4m^2 = 2q^2$, т.е. $q^2 = 2m^2$, значит $q$ тоже чётно. Числа $p$ и $q$ имеют общий делитель $2$ --- противоречие с несократимостью. Итак, $\sqrt{2} \notin \Q$.

\textit{Расширение поля $\Q$:} В поле $\R$ уже есть элемент $\sqrt{2}$ (положительный корень уравнения $x^2 = 2$; его существование даёт, например, аксиома полноты: множество $\{ x \in \Q : x > 0,\ x^2 < 2 \}$ ограничено сверху в $\R$, его супремум $s$ удовлетворяет $s^2 = 2$). Рассмотрим подмножество $\R$:
\[
\Q(\sqrt{2}) = \{ a + b\sqrt{2} \mid a, b \in \Q \}.
\]
Оно содержит $\Q$ (при $b=0$) и $\sqrt{2}$ (при $a=0$, $b=1$). Сложение и вычитание очевидным образом дают элементы того же вида. Произведение: $(a + b\sqrt{2})(c + d\sqrt{2}) = (ac + 2bd) + (ad + bc)\sqrt{2} \in \Q(\sqrt{2})$. Обратный элемент: для $a + b\sqrt{2} \neq 0$ (т.е. $(a,b) \neq (0,0)$) число $a - b\sqrt{2}$ ненулевое (иначе $\sqrt{2} = a/b \in \Q$), и
\[
\frac{1}{a + b\sqrt{2}} = \frac{a - b\sqrt{2}}{(a+b\sqrt{2})(a-b\sqrt{2})} = \frac{a - b\sqrt{2}}{a^2 - 2b^2} \in \Q(\sqrt{2}),
\]
поскольку знаменатель $a^2 - 2b^2 \in \Q$ и отличен от нуля (при $b \neq 0$ иначе было бы $\sqrt{2} \in \Q$; при $b=0$ имеем $a \neq 0$). Таким образом, $\Q(\sqrt{2})$ --- поле, содержащее $\Q$ и получающееся из $\Q$ добавлением корня уравнения $x^2 - 2 = 0$; это расширение поля $\Q$ степени $2$ (каждый элемент однозначно записывается как $a + b\sqrt{2}$ с $a,b \in \Q$).

\item \textbf{Построение функции $a^x$: изоморфизм $(\R,+) \to (\R^+, \cdot)$, сохраняющий порядок.}

Рассмотрим группу вещественных чисел по сложению $(\R, +)$ и группу положительных вещественных по умножению $(\R^+, \cdot)$. Требуется построить функцию $f\colon \R \to \R^+$ такую, что:
\begin{enumerate}[label=(\roman*), nosep]
    \item $f(x + y) = f(x) f(y)$ для всех $x, y \in \R$ (гомоморфизм групп),
    \item $f(1) = a$, где $a > 0$, $a \neq 1$ фиксировано,
    \item $f$ \textbf{сохраняет порядок}: из $x < y$ следует $f(x) < f(y)$.
\end{enumerate}

\textit{Вывод $a > 1$:} Из (i) при $x=y=0$ получаем $f(0)=1$. Из (iii) при $0 < 1$ имеем $f(0) < f(1)$, т.е. $1 < a$. Таким образом, \textbf{условие сохранения порядка возможно только при $a > 1$}. (При $0 < a < 1$ гомоморфизм с $f(1)=a$ был бы строго убывающим.)

\textit{Пошаговое построение:}
\begin{itemize}[nosep]
    \item $f(0)=1$; по индукции $f(n)=a^n$ при $n \in \N$; из $f(-n)f(n)=f(0)=1$ получаем $f(-n)=a^{-n}$.
    \item Для $r = p/q \in \Q$ имеем $f(r)^q = f(p) = a^p$, значит $f(r)$ --- единственное положительное число, $q$-я степень которого равна $a^p$; полагаем $a^{p/q} = (a^p)^{1/q}$. При $a > 1$ функция $\Q \ni r \mapsto a^r$ строго возрастает.
    \item Для произвольного $x \in \R$ определим $f(x) = \sup\,\{ a^r : r \in \Q,\ r < x \}$ (супремум по непустому ограниченному сверху множеству существует в $\R$). Эквивалентно: $f(x) = \lim_{n\to\infty} a^{r_n}$ для любой последовательности рациональных $r_n \to x$ (достаточно монотонно возрастающей).
\end{itemize}

\textit{Вывод непрерывности:} Функция $f$ на $\R$ по построению возрастает и совпадает с $a^r$ на $\Q$. Для любой точки $x \in \R$ и последовательности $x_n \to x$ выберем рациональные $r_n', r_n''$ так, что $r_n' \le x_n \le r_n''$ и $r_n', r_n'' \to x$. Тогда $a^{r_n'} \le f(x_n) \le a^{r_n''}$; при $n \to \infty$ оба края стремятся к $f(x)$ (так как $f(x)$ по определению совпадает с пределом $a^r$ по рациональным $r \to x$). По «лемме о двух милиционерах» $f(x_n) \to f(x)$. Итак, \textbf{из сохранения порядка и гомоморфизма следует непрерывность} $f$.

Итог: условие «гомоморфизм $(\R,+) \to (\R^+, \cdot)$ с $f(1)=a$, сохраняющий порядок» возможно только при $a > 1$ и однозначно задаёт непрерывную функцию $a^x$ на $\R$. Для $a=1$ полагаем $1^x = 1$.
\end{enumerate}

\begin{remark}[Связь с методами Главы 2]
В процессе 2 пространство состояний получено \textbf{факторизацией} $\R$ по $\Z$; канторово множество --- \textbf{пересечение} вложенных замкнутых множеств (метод выделения на каждом шаге и переход к пределу); мера нуль получается как дополнение к объединению удалённых интервалов заданной длины. Неподвижная точка (процесс 1) опирается на полноту $\R$; расширение $\Q(\sqrt{2})$ (процесс 4) --- пример расширения поля $\Q$ присоединением корня многочлена $x^2 - 2$. Построение $a^x$ использует \textbf{продолжение по монотонности} (супремум по $a^r$, $r<x$) с $\Q$ на $\R$; непрерывность выводится из сохранения порядка.
\end{remark}


\subsection{Ряды и бесконечные произведения}

\begin{definition}[Числовой ряд]
Пусть $\langle a_n \rangle_{n=1}^\infty$ --- последовательность вещественных чисел. \textbf{Рядом} называется формальная сумма
$$\sum_{n=1}^\infty a_n = a_1 + a_2 + a_3 + \ldots$$
\textbf{Частичная сумма}: $S_n = a_1 + a_2 + \ldots + a_n$. Ряд \textbf{сходится}, если последовательность $\langle S_n \rangle$ имеет конечный предел; тогда пишут $\sum_{n=1}^\infty a_n = \lim_{n\to\infty} S_n$. Ряд \textbf{расходится}, если предела нет (включая случай $\lim S_n = \pm\infty$).
\end{definition}

\begin{definition}[Абсолютная сходимость]
Ряд $\sum a_n$ \textbf{сходится абсолютно}, если сходится ряд $\sum |a_n|$.
\end{definition}

\begin{theorem}[Критерий Коши для рядов]
Ряд $\sum a_n$ сходится тогда и только тогда, когда $\forall \ep > 0\; \exists N \in \N\; \forall m \ge n \ge N: |a_n + a_{n+1} + \ldots + a_m| < \ep$.
\end{theorem}

\begin{theorem}[Необходимое условие сходимости]
Если ряд $\sum a_n$ сходится, то $a_n \to 0$ при $n \to \infty$. Обратное неверно (например, $\sum 1/n$ расходится).
\end{theorem}

\begin{theorem}[Геометрический ряд]
Для $|q| < 1$ ряд $\sum_{n=0}^\infty q^n$ сходится и $\sum_{n=0}^\infty q^n = \frac{1}{1-q}$. При $|q| \ge 1$ ряд расходится.
\end{theorem}

\begin{theorem}[Абсолютная сходимость $\Rightarrow$ сходимость]
Если ряд $\sum a_n$ сходится абсолютно, то он сходится. Более того, при абсолютной сходимости сумма ряда не зависит от порядка слагаемых (теорема Дирихле о перестановке).
\end{theorem}

\begin{theorem}[Римана о перестановке членов условно сходящегося ряда]
Пусть ряд $\sum a_n$ сходится условно (сходится, но не абсолютно). Тогда для любого $L \in \R \cup \{-\infty, +\infty\}$ существует перестановка $\pi$ натурального ряда такая, что $\sum_{n=1}^\infty a_{\pi(n)} = L$.

\textit{Идея:} берём положительные и отрицательные члены по отдельности; оба подряда расходятся. «Выписываем» положительные, пока частичная сумма не превысит $L$, затем добавляем отрицательные, пока не опустимся ниже $L$, и т.д.
\end{theorem}

\begin{definition}[Бесконечное произведение]
Бесконечным произведением называется выражение $\prod_{n=1}^\infty (1 + u_n)$, где $u_n \in \R$. Частичное произведение: $P_n = (1+u_1)(1+u_2)\ldots(1+u_n)$. Произведение \textbf{сходится} к $P \neq 0$, если $\lim_{n\to\infty} P_n = P$. (Нулевое значение считают расхождением, если только какой-то множитель не равен нулю.)
\end{definition}

\begin{theorem}[Критерий сходимости бесконечного произведения]
Пусть $u_n > -1$ для всех $n$. Произведение $\prod_{n=1}^\infty (1+u_n)$ сходится к ненулевому числу тогда и только тогда, когда сходится ряд $\sum_{n=1}^\infty \ln(1+u_n)$, что при $u_n \to 0$ эквивалентно сходимости $\sum u_n$. В частности, при $0 \le u_n < 1$: $\prod(1-u_n) > 0 \iff \sum u_n < \infty$.
\end{theorem}

\subsection{Пространство \texorpdfstring{$\ell^2$}{l2}}

\begin{definition}[Пространство $\ell^2$]
Пространство $\ell^2$ --- множество всех последовательностей $x = \langle x_n \rangle_{n=1}^\infty$ вещественных (или комплексных) чисел, суммируемых с квадратом:
$$\ell^2 = \left\{ \langle x_n \rangle : \sum_{n=1}^\infty |x_n|^2 < \infty \right\}$$
Норма: $\|x\|_2 = \sqrt{\sum_{n=1}^\infty |x_n|^2}$. Скалярное произведение: $\langle x \mid y \rangle = \sum_{n=1}^\infty \overline{x_n} y_n$ (для вещественного случая черта не нужна).
\end{definition}

\begin{theorem}[Неравенство Коши--Шварца для $\ell^2$]
Для $x, y \in \ell^2$ выполнено $|\langle x \mid y \rangle| \le \|x\|_2 \|y\|_2$.
\end{theorem}

\begin{theorem}[Полнота $\ell^2$]
Пространство $\ell^2$ с метрикой $d(x,y) = \|x-y\|_2$ является полным метрическим пространством (гильбертово пространство последовательностей). Всякая последовательность Коши в $\ell^2$ сходится к элементу $\ell^2$.
\end{theorem}

\begin{remark}[Связь с интегралом]
Пространство $\ell^2$ --- дискретный аналог пространства $L^2$ функций, интегрируемых с квадратом. Интеграл заменяется суммой. Аналогия будет уточнена в разделе о мере и интеграле Лебега.
\end{remark}

\section{Метрика, топология, пространство $\R^n$, норма}
\textbf{Цель:} Знакомство с понятиями метрики, топологии и борелевской структуры.

\subsection{Метрика и сходимость}

\begin{definition}[Метрика]
Метрика на множестве $X$ --- это функция $d: X \times X \to \R_{\ge 0}$ такая, что для всех $x, y, z \in X$:
\begin{enumerate}
    \item \textbf{Тождество неразличимых:} $d(x, y) = 0 \iff x = y$
    \item \textbf{Симметрия:} $d(x, y) = d(y, x)$
    \item \textbf{Неравенство треугольника:} $d(x, z) \le d(x, y) + d(y, z)$
\end{enumerate}
Пара $\langle X, d \rangle$ называется метрическим пространством.
\end{definition}

\begin{definition}[Норма]
Норма на векторном пространстве $V$ над $\R$ (или $\C$) --- функция $\|\cdot\|: V \to \R_{\ge 0}$: $\|x\| = 0 \iff x = 0$; $\|\la x\| = |\la| \|x\|$; $\|x + y\| \le \|x\| + \|y\|$. Норма порождает метрику: $d(x,y) = \|x - y\|$.
\end{definition}

\begin{definition}[Евклидова метрика и норма на $\R^n$]
Для $x = (x_1,\ldots,x_n)$, $y = (y_1,\ldots,y_n) \in \R^n$:
$$|x| = \|x\|_2 = \sqrt{x_1^2 + \ldots + x_n^2}, \quad d(x,y) = |x - y|$$
\end{definition}

\begin{example}[Другие метрики на $\R^n$]
\begin{itemize}
    \item \textbf{Манхэттен:} $d_1(x,y) = \sum_{i=1}^n |x_i - y_i|$
    \item \textbf{Чебышева:} $d_\infty(x,y) = \max_i |x_i - y_i|$
    \item На $\R^n$ все эти метрики \textbf{эквивалентны} (порождают одну и ту же топологию): существуют константы $C_1, C_2 > 0$ такие, что $C_1 d_2 \le d_1 \le C_2 d_2$ и т.п.
\end{itemize}
\end{example}

\begin{definition}[Сходимость по метрике]
Последовательность $\langle x_n \rangle$ в метрическом пространстве $\langle X, d \rangle$ \textbf{сходится} к $x \in X$, если $d(x_n, x) \to 0$ при $n \to \infty$. Обозначение: $x_n \to x$.
\end{definition}

\subsection{Топология: открытые и замкнутые множества}

\begin{definition}[Открытый шар и база топологии]
Пусть $\langle X, d \rangle$ --- метрическое пространство. \textbf{Открытый шар} радиуса $\ep > 0$ с центром в $x$: $B(x, \ep) = \{y \in X : d(x,y) < \ep\}$. Множество $U \subseteq X$ \textbf{открыто}, если для каждой точки $x \in U$ существует $\ep > 0$ с $B(x,\ep) \subseteq U$. Семейство всех открытых множеств образует \textbf{топологию} на $X$.
\end{definition}

\begin{definition}[Замкнутое множество, замыкание, внутренность]
Множество $F$ \textbf{замкнуто}, если его дополнение $X \setminus F$ открыто. Эквивалентно: $F$ содержит все свои предельные точки. \textbf{Замыкание} $\overline{A}$ множества $A$ --- наименьшее замкнутое множество, содержащее $A$. \textbf{Внутренность} $\mathrm{int}(A)$ --- наибольшее открытое множество, содержащееся в $A$.
\end{definition}

\begin{theorem}[Куратовский: максимум 14 множеств]
Рассмотрим операции \textbf{внутренности} $A \mapsto \mathrm{int}(A)$ и \textbf{дополнения} $A \mapsto A^c = X \setminus A$ на подмножествах топологического пространства. Начиная с произвольного множества $A$, сколько различных множеств можно получить, последовательно применяя эти две операции?

\textbf{Ответ:} Не более 14. Более того, на вещественной прямой существует множество $A$, для которого достигается ровно 14 различных множеств (теорема Куратовского, 1922). Аналогичный результат верен для пары замыкание--дополнение.
\end{theorem}

\begin{definition}[Компактность]
Множество $K \subseteq X$ \textbf{компактно}, если из любого открытого покрытия $K$ можно выбрать конечное подпокрытие. В метрических пространствах: $K$ компактно $\iff$ из любой последовательности в $K$ можно выбрать сходящуюся подпоследовательность.
\end{definition}

\begin{theorem}[Гейне--Бореля]
В $\R^n$ с евклидовой метрикой множество компактно $\iff$ оно замкнуто и ограничено.
\end{theorem}

\subsection{Борелевская \texorpdfstring{$\sigma$}{сигма}-алгебра}

\begin{definition}[$\sigma$-алгебра и борелевские множества]
\textbf{$\sigma$-алгебра} на множестве $X$ --- семейство $\Bcal \subseteq \Pcal(X)$, замкнутое относительно дополнения и счётных объединений (значит, и пересечений), содержащее $X$ и $\emptyset$.

\textbf{Борелевская $\sigma$-алгебра} $\Bcal(\R^n)$ --- наименьшая $\sigma$-алгебра на $\R^n$, содержащая все открытые множества (относительно евклидовой топологии). Её элементы называются \textbf{борелевскими множествами}.
\end{definition}

\begin{remark}[Порождение борелевской структуры]
Борелевскую $\sigma$-алгебру порождают: открытые множества, замкнутые множества, интервалы (на прямой), прямоугольники (на плоскости). Она замкнута относительно счётных объединений и пересечений.
\end{remark}

\subsection{Категории Бэра и континуум-гипотеза}

\begin{definition}[Множества первой и второй категории Бэра]
Подмножество $A$ полного метрического пространства $X$ имеет \textbf{первую категорию} (является \textbf{тощим}), если $A$ содержится в объединении счётного семейства нигде не плотных множеств. Остальные множества имеют \textbf{вторую категорию} (являются \textbf{жирными}).
\end{definition}

\begin{theorem}[Теорема Бэра]
Полное метрическое пространство не может быть представлено как счётное объединение нигде не плотных множеств. Иначе говоря, $X$ имеет вторую категорию в себе.
\end{theorem}

\begin{remark}[Континуум-гипотеза и теорема Бэра]
\textbf{Континуум-гипотеза} (Кантор, 1878): не существует множества мощности строго между $\aleph_0$ и $2^{\aleph_0}$. Гёдель (1940) и Коэн (1963) показали, что гипотеза независима от аксиом \ZF: её нельзя ни доказать, ни опровергнуть.

\textbf{Связь с категориями Бэра:} Теорема Бэра утверждает, что $\R$ нельзя покрыть счётным объединением тощих множеств. Естественный вопрос: каково минимальное число тощих множеств, объединением которых можно покрыть $\R$? При справедливости CH это число равно $\aleph_1$; при отрицании CH --- равно $2^{\aleph_0}$. Тем самым CH уточняет количественную сторону результата Бэра.
\end{remark}

\section{Мера, Вероятность и Хаос}
\textbf{Цель:} Формализация неопределенности и структуры событий.

\subsection{Теория меры и Интеграл Лебега}

\begin{definition}[$\si$-алгебра]
$\si$-алгебра на множестве $X$ --- это семейство $\Si \subseteq \Pcal(X)$ такое, что:
\begin{enumerate}
    \item $X \in \Si$
    \item Если $A \in \Si$, то $X \setminus A \in \Si$ (замкнутость относительно дополнения)
    \item Если $\langle A_n \rangle_{n=1}^\infty \subseteq \Si$, то $\bigcup_{n=1}^\infty A_n \in \Si$ (замкнутость относительно счетных объединений)
\end{enumerate}
Пара $\langle X, \Si \rangle$ называется измеримым пространством.
\end{definition}

\begin{definition}[Мера]
Мера на измеримом пространстве $\langle X, \Si \rangle$ --- это функция $\mu: \Si \to [0, +\infty]$ такая, что:
\begin{enumerate}
    \item $\mu(\emptyset) = 0$
    \item \textbf{$\si$-аддитивность:} Если $\langle A_n \rangle_{n=1}^\infty$ --- дизъюнктные множества из $\Si$, то
    $$\mu\left(\bigcup_{n=1}^\infty A_n\right) = \sum_{n=1}^\infty \mu(A_n)$$
\end{enumerate}
Тройка $\langle X, \Si, \mu \rangle$ называется пространством с мерой.
\end{definition}

\begin{definition}[Мера Лебега]
Мера Лебега на $\R$ --- это единственная мера $\la$ на борелевской $\si$-алгебре $\Bcal(\R)$ такая, что $\la([a, b]) = b - a$ для любого интервала $[a, b]$.

\textbf{Борелевская $\si$-алгебра} $\Bcal(\R)$ --- наименьшая $\si$-алгебра, содержащая все открытые множества в $\R$.
\end{definition}

\begin{definition}[Множество меры нуль]
Множество $A \in \Si$ имеет меру нуль, если $\mu(A) = 0$.

Свойство выполняется \textbf{почти всюду} (almost everywhere, a.e.), если множество точек, где оно не выполняется, имеет меру нуль.
\end{definition}

\begin{theorem}[Канторово множество]
Канторово множество $C$ строится так: начинаем с $C_0 = [0, 1]$. На каждом шаге удаляем среднюю треть каждого интервала. $C = \bigcap_{n=0}^\infty C_n$.

\textbf{Свойства:}
\begin{enumerate}
    \item $C$ несчетно (мощность континуума)
    \item $\la(C) = 0$ (мера Лебега равна нулю)
\end{enumerate}
\end{theorem}

\begin{definition}[Измеримая функция]
Функция $f: X \to \R$ измерима относительно $\si$-алгебры $\Si$, если для любого $a \in \R$ множество $\{x \in X \mid f(x) > a\} \in \Si$.
\end{definition}

\begin{definition}[Интеграл Лебега]\;
\begin{enumerate}
    \item Для простой функции $f = \sum_{i=1}^n a_i \chi_{A_i}$ (где $\chi_{A_i}$ --- характеристическая функция): 
    $$\int_X f \, d\mu = \sum_{i=1}^n a_i \mu(A_i)$$
    \item Для неотрицательной измеримой функции $f$:
    $$\int_X f \, d\mu = \sup\left\{\int_X g \, d\mu \mid g \text{ простая}, 0 \le g \le f\right\}$$
    \item Для произвольной измеримой функции $f$ (через разложение $f = f^+ - f^-$, где $f^+ = \max(f, 0)$, $f^- = \max(-f, 0)$):
    $$\int_X f \, d\mu = \int_X f^+ \, d\mu - \int_X f^- \, d\mu$$
    (если хотя бы один интеграл конечен)
\end{enumerate}
\end{definition}

\begin{theorem}[Теорема о монотонной сходимости]
Если $\langle f_n \rangle$ --- последовательность неотрицательных измеримых функций, $f_n \le f_{n+1}$ и $f_n \to f$ почти всюду, то
$$\lim_{n \to \infty} \int_X f_n \, d\mu = \int_X f \, d\mu$$
\end{theorem}

\begin{theorem}[Теорема о доминированной сходимости]
Если $\langle f_n \rangle$ --- последовательность измеримых функций, $f_n \to f$ почти всюду, и существует интегрируемая функция $g$ такая, что $|f_n| \le g$ почти всюду, то
$$\lim_{n \to \infty} \int_X f_n \, d\mu = \int_X f \, d\mu$$
\end{theorem}

\begin{remark}
Интеграл Лебега обобщает интеграл Римана: если функция интегрируема по Риману, то она интегрируема по Лебегу, и значения интегралов совпадают. Однако класс интегрируемых по Лебегу функций шире.
\end{remark}

\begin{remark}
Мера позволяет различать «невозможное» (пустое множество) и «почти невозможное» (множество меры нуль). В моделировании это критически важно для описания редких событий («Черных лебедей»).
\end{remark}

\begin{exercise}[Парадокс Влияния]
Рассмотрим отрезок времени $[0, 1]$. Пусть моменты «озарения» распределены по структуре Канторова множества $K$.
\begin{enumerate}
    \item Какова суммарная длительность озарений (мера $K$)? (Ответ: 0).
    \item Каково количество моментов озарения (мощность $K$)? (Ответ: Континуум, как и весь отрезок).
    \item \textbf{Интерпретация:} Постройте модель, где событие происходит «постоянно» (в несчетном числе моментов), но «никогда» (занимает 0 времени). Применимо ли это к процессу мышления?
\end{enumerate}
\end{exercise}

\subsection{Случайные блуждания и процессы}

\begin{definition}[Вероятностное пространство]
Вероятностное пространство --- это тройка $\langle \Om, \Fcal, \P \rangle$, где:
\begin{enumerate}
    \item $\Om$ --- пространство элементарных исходов
    \item $\Fcal$ --- $\si$-алгебра событий на $\Om$
    \item $\P: \Fcal \to [0, 1]$ --- вероятностная мера: $\P(\Om) = 1$ и $\P$ $\si$-аддитивна
\end{enumerate}
\end{definition}

\begin{definition}[Случайная величина]
Случайная величина --- это измеримая функция $X: \Om \to \R$ (относительно $\si$-алгебры $\Fcal$ и борелевской $\si$-алгебры на $\R$).
\end{definition}

\begin{definition}[Случайный процесс]
Случайный процесс --- это семейство случайных величин $\{X_t\}_{t \in T}$, где $T$ --- множество индексов (обычно время: $T = \N$, $\Z$, $\R_{\geq 0}$ и т.д.).

Для каждого $\om \in \Om$ функция $t \mapsto X_t(\om)$ называется траекторией процесса.
\end{definition}

\begin{definition}[Марковское свойство]
Случайный процесс $\{X_t\}_{t \in T}$ обладает марковским свойством, если для любого $t \in T$ и любого $s > t$ условное распределение $X_s$ при заданном $\{X_u\}_{u \le t}$ зависит только от $X_t$:

$$\P(X_s \in A \mid \{X_u\}_{u \le t}) = \P(X_s \in A \mid X_t)$$

\textbf{Интерпретация:} «Будущее зависит только от настоящего, а не от прошлого» (процесс без памяти).
\end{definition}

\begin{definition}[Цепь Маркова]
Дискретная цепь Маркова --- это последовательность случайных величин $\langle X_n \rangle_{n=0}^\infty$, где:
$$\P(X_{n+1} = j \mid X_n = i, X_{n-1} = i_{n-1}, \ldots, X_0 = i_0) = \P(X_{n+1} = j \mid X_n = i)$$
для всех $n$ и всех состояний $i, j, i_{n-1}, \ldots, i_0$.
\end{definition}

\begin{remark}[Применение в физике: броуновское движение]
Броуновское движение (винеровский процесс) --- это непрерывный предел случайного блуждания. Частица «забывает», откуда прилетела, на каждом шаге. Это марковский процесс: будущее положение зависит только от текущего, а не от всей траектории.
\end{remark}

\begin{definition}[Стационарное распределение]
Для цепи Маркова стационарное распределение $\pi$ --- это распределение вероятностей на множестве состояний такое, что:
$$\pi_j = \sum_i \pi_i P_{ij}$$
где $P_{ij} = \P(X_{n+1} = j \mid X_n = i)$ --- вероятности переходов.

Если цепь начинается со стационарного распределения, то распределение остается неизменным во времени.
\end{definition}

\begin{remark}[Физическая интерпретация: тепловое равновесие]
В статистической механике стационарное распределение соответствует тепловому равновесию. Система приходит к состоянию, которое не меняется со временем (распределение Больцмана).
\end{remark}

\begin{definition}[Винеровский процесс]
Винеровский процесс (броуновское движение) $W_t$ --- это случайный процесс на $[0, \infty)$ такой, что:
\begin{enumerate}
    \item $W_0 = 0$ почти наверное
    \item Для $0 \le s < t$ приращение $W_t - W_s$ имеет нормальное распределение $\mathcal{N}(0, t-s)$
    \item Для непересекающихся интервалов $[s_1, t_1], [s_2, t_2], \ldots$ приращения $W_{t_1} - W_{s_1}, W_{t_2} - W_{s_2}, \ldots$ независимы
    \item Траектории непрерывны почти наверное
\end{enumerate}
\end{definition}

\begin{theorem}[Недифференцируемость траекторий винеровского процесса]
Траектории винеровского процесса почти наверное недифференцируемы в любой точке $t > 0$.
\end{theorem}

\begin{remark}
Винеровский процесс является пределом случайного блуждания при масштабировании времени и пространства (центральная предельная теорема).
\end{remark}

\begin{exercise}[Эволюция Идеи]
Смоделируем развитие философской идеи как случайное блуждание на графе $\mathbb{Z}$ (целые числа), где $0$ — нейтральность, $+\infty$ — абсолютное принятие, $-\infty$ — забвение.
\begin{enumerate}
    \item Пусть на каждом шаге идея сдвигается на $+1$ с вероятностью $p$ и на $-1$ с вероятностью $q=1-p$.
    \item Найдите вероятность того, что идея когда-либо вернется в исходную точку (задача о разорении игрока).
    \item Введите «поглощающее состояние» (догма): если идея достигает значения $N$, она перестает меняться. Как зависит время догматизации от «радикальности» шагов?
\end{enumerate}
\end{exercise}


\newpage

\section{Графы, Деревья и Динамика}
\textbf{Цель:} Изучить структуру связей и эволюцию во времени. Показать, как абстрактные конструкции (случайные графы, марковские цепи) --- это язык, на котором записаны законы физики.

\subsection{Случайные графы и фазовые переходы}

\begin{definition}[Модель Эрдеша-Реньи $G(n, p)$]
Модель случайного графа $G(n, p)$ строится так:
\begin{enumerate}
    \item Берем $n$ вершин
    \item Каждую пару вершин соединяем ребром с вероятностью $p$ независимо от других пар
\end{enumerate}
Результат --- случайный граф на $n$ вершинах.
\end{definition}

\begin{theorem}[Гигантская компонента]
Существует критический порог $p = \frac{1}{n}$. 
\begin{enumerate}
    \item Если $p < \frac{1}{n}$ (подкритический режим): граф состоит из множества маленьких компонент, размер самой большой компоненты имеет порядок $O(\log n)$
    \item Если $p = \frac{c}{n}$ с $c > 1$ (надкритический режим): граф внезапно превращается из кучки разрозненных островов в единый материк --- появляется гигантская компонента, содержащая положительную долю всех вершин
\end{enumerate}
\end{theorem}

\begin{remark}[Физическая интерпретация: теория перколяции]
Теория перколяции --- это математическая модель фазового перехода. Например:
\begin{itemize}
    \item Превращение воды в лед: при критической температуре структура внезапно меняется
    \item Возникновение проводимости в материале: при критической концентрации проводящих связей материал внезапно становится проводящим
\end{itemize}
Фазовый переход соответствует появлению гигантской компоненты в случайном графе.
\end{remark}

\subsection{Графы и деревья}

\begin{definition}[Граф]
Граф --- это пара $G = \langle V, E \rangle$, где:
\begin{enumerate}
    \item $V$ --- множество вершин (узлов)
    \item $E \subseteq V \times V$ --- множество рёбер (связей между вершинами)
\end{enumerate}
Граф называется \textbf{ориентированным}, если рёбра упорядочены; \textbf{неориентированным} --- если порядок не важен.
\end{definition}

\begin{definition}[Путь в графе]
Путь в графе $G = \langle V, E \rangle$ --- это последовательность вершин $(v_0, v_1, \ldots, v_n)$ такая, что $\langle v_i, v_{i+1} \rangle \in E$ для всех $i = 0, \ldots, n-1$.

Путь называется \textbf{простым}, если все вершины различны. Путь называется \textbf{циклом}, если $v_0 = v_n$ и $n > 0$.
\end{definition}

\begin{definition}[Дерево]
Дерево --- это связный граф без циклов. Эквивалентные определения:
\begin{enumerate}
    \item Связный граф с $n$ вершинами и $n-1$ рёбрами
    \item Граф, в котором любые две вершины соединены единственным простым путём
    \item Связный граф без циклов
\end{enumerate}
\end{definition}

\begin{definition}[Корневое дерево]
Корневое дерево --- это дерево $\Tcal = \langle V, E, r \rangle$ с выделенной вершиной $r \in V$ (корнем).

Для вершин $x, y \in V$:
\begin{enumerate}
    \item $x$ --- \textbf{предок} $y$ (обозначение: $x \le y$), если $x$ лежит на единственном пути от $r$ к $y$
    \item $x$ --- \textbf{родитель} $y$, если $x$ --- непосредственный предок $y$ (т.е. $\langle x, y \rangle \in E$)
    \item $y$ --- \textbf{потомок} $x$, если $x \le y$
    \item $\Succ(x) = \{y \in V \mid \langle x, y \rangle \in E\}$ --- множество \textbf{непосредственных потомков} $x$
    \item $\Pred(x)$ --- \textbf{родитель} $x$ (если $x \neq r$)
\end{enumerate}
\end{definition}

\begin{definition}[Уровень и высота]
В корневом дереве:
\begin{enumerate}
    \item \textbf{Уровень} вершины $x$ --- это длина пути от $r$ до $x$: $|x| = d(r, x)$
    \item \textbf{Высота} дерева --- это максимальный уровень вершины (может быть бесконечной)
\end{enumerate}
\end{definition}

\begin{definition}[Лист и внутренняя вершина]
\begin{enumerate}
    \item \textbf{Лист} --- вершина без потомков: $\Succ(x) = \emptyset$
    \item \textbf{Внутренняя вершина} --- вершина с потомками: $\Succ(x) \neq \emptyset$
\end{enumerate}
\end{definition}

\begin{definition}[Бесконечный путь]
Бесконечный путь (максимальная цепь) в дереве --- это бесконечная последовательность вершин $(v_0, v_1, v_2, \ldots)$ такая, что $v_0 = r$ и $\langle v_i, v_{i+1} \rangle \in E$ для всех $i \in \N$.

Обозначение: $\Om$ --- множество всех бесконечных путей в дереве.
\end{definition}

\begin{definition}[Цилиндр]
Для вершины $x \in V$ цилиндр $C_x$ --- это множество всех бесконечных путей, проходящих через $x$:
$$C_x = \{\om \in \Om \mid x \in \om\}$$
\end{definition}

\begin{theorem}[Компактность пространства путей]
Если дерево имеет конечное ветвление (каждая вершина имеет конечное число потомков), то пространство $\Om$ всех бесконечных путей компактно в топологии, порождённой цилиндрами.
\end{theorem}

\begin{exercise}[Модель развития концепции]
Постройте дерево, где корень --- исходная философская концепция, а каждая вершина --- её интерпретация. Рёбра --- логические следствия или уточнения.
\begin{enumerate}
    \item Определите метрику на пространстве путей (например, через расстояние между первыми точками расхождения).
    \item Покажите, что цилиндры образуют базу топологии.
    \item Интерпретируйте компактность: что означает, что «все пути развития концепции компактны»?
\end{enumerate}
\end{exercise}

\subsection{Лемма Кёнига}

\begin{theorem}[Лемма Кёнига]
Бесконечное дерево с конечным ветвлением имеет бесконечный путь.

\textbf{Формулировка:} Если $\Tcal = \langle V, E, r \rangle$ --- бесконечное корневое дерево, и каждая вершина имеет конечное число потомков, то существует бесконечный путь, начинающийся в корне.
\end{theorem}

\begin{proof}[Набросок доказательства]
Построим путь индуктивно:
\begin{enumerate}
    \item Начинаем с корня $v_0 = r$
    \item Если выбрана вершина $v_n$ на уровне $n$, то среди её потомков найдётся хотя бы один, у которого бесконечно много потомков (иначе дерево было бы конечным)
    \item Выбираем такой потомок как $v_{n+1}$
    \item Продолжаем процесс бесконечно
\end{enumerate}
Полученная последовательность $\langle v_0, v_1, v_2, \ldots \rangle$ --- искомый бесконечный путь.
\end{proof}

\begin{remark}
Лемма Кёнига эквивалентна Аксиоме Выбора для счётных семейств (\AC$_\om$). Она критически важна для доказательства существования предельных объектов в деревьях.
\end{remark}

\begin{exercise}[Применение леммы Кёнига]
Рассмотрим дерево решений для бесконечной игры между двумя игроками.
\begin{enumerate}
    \item Покажите, что если дерево бесконечно и имеет конечное ветвление, то игра может продолжаться бесконечно (лемма Кёнига).
    \item Постройте пример дерева с бесконечным ветвлением, в котором нет бесконечного пути.
    \item Интерпретируйте результат: что означает для моделирования процессов принятия решений?
\end{enumerate}
\end{exercise}

\subsection{Кодирование деревьев ординалами}

\begin{definition}[Ординальная высота дерева]
Для корневого дерева $\Tcal = \langle V, E, r \rangle$ определим ранг вершины $v$ ($rank(v)$) трансфинитной рекурсией:
\begin{itemize}
    \item Если $v$ --- лист, то $rank(v) = 0$
    \item Если $v$ --- не лист, то $rank(v) = \sup \{rank(u) + 1 \mid u \in \Succ(v)\}$
\end{itemize}
Ординальная высота дерева определяется как $height(\Tcal) = rank(r)$.
\end{definition}

\begin{theorem}[Кодирование дерева ординалом]
Всякое фундированное дерево (дерево без бесконечных путей) имеет ординальную высоту. Более того, можно сопоставить каждому дереву ординал, который кодирует его структуру.
\end{theorem}

\begin{definition}[Кодирование через порядковые номера]
Пусть $\Tcal$ --- дерево с корнем $r$. Определим функцию $code: V \to \On$ рекурсивно:
\begin{itemize}
    \item $code(r) = 0$
    \item Если $v$ имеет потомков $u_1, u_2, \ldots, u_n$ (в некотором порядке), то:
    $$code(u_i) = code(v) + \sum_{j=1}^{i-1} \omega^{rank(u_j)}$$
    где $rank(u_j)$ --- ранг поддерева с корнем $u_j$
\end{itemize}
\end{definition}

\begin{example}[Кодирование конечного дерева]
Для дерева с корнем и двумя потомками:
\begin{itemize}
    \item $code(r) = 0$
    \item Если левый потомок --- лист: $code(\text{левый}) = 0 + \omega^0 = 1$
    \item Если правый потомок имеет поддерево высоты 1: $code(\text{правый}) = 0 + \omega^0 + \omega^1 = 1 + \omega$
\end{itemize}
\end{example}

\begin{theorem}[Изоморфизм деревьев и ординалов]
Два фундированных дерева изоморфны (как частично упорядоченные множества) тогда и только тогда, когда их ординальные коды совпадают.
\end{theorem}

\begin{remark}[Применение к ветвящимся процессам]
Кодирование деревьев ординалами позволяет классифицировать генеалогические деревья в процессах Гальтона-Ватсона по их структуре, независимо от конкретных значений случайных величин.
\end{remark}

\begin{remark}[Философский смысл кодирования деревьев ординалами]
Проблема: Как отличить одно бесконечное дерево от другого?

Решение: Мы приписываем листьям ранг 0. Вершинам выше --- ранг, равный супремуму рангов потомков + 1. Для бесконечных деревьев это требует трансфинитных чисел (ординалов).

\textbf{Философский смысл:} Это способ «посчитать» сложность иерархии, даже если она бесконечна. Дерево высоты $\omega$ (потенциальная бесконечность) проще, чем дерево высоты $\omega + 1$ (актуальная бесконечность). Ординалы позволяют различать структуры бесконечных процессов по их «глубине» и «сложности».
\end{remark}

\section{Мартингалы}
\textbf{Цель:} Освоить теорию мартингалов как инструмент для работы с условными ожиданиями и справедливыми играми.

\subsection{Условное математическое ожидание}

\begin{definition}[$\si$-алгебра, порождённая случайной величиной]
Для случайной величины $X$ на $\langle \Om, \Fcal, \P \rangle$:
$$\si(X) = \{X^{-1}(B) \mid B \in \Bcal(\R)\}$$
--- наименьшая $\si$-алгебра, относительно которой $X$ измерима.
\end{definition}

\begin{definition}[Условное математическое ожидание]
Пусть $X$ --- интегрируемая случайная величина, $\Gcal \subseteq \Fcal$ --- под-$\si$-алгебра. Условное математическое ожидание $\E[X \mid \Gcal]$ --- это случайная величина такая, что:
\begin{enumerate}
    \item $\E[X \mid \Gcal]$ измерима относительно $\Gcal$
    \item Для любого $A \in \Gcal$: $\int_A \E[X \mid \Gcal] \, d\P = \int_A X \, d\P$
\end{enumerate}
\end{definition}

\begin{theorem}[Существование и единственность]
Условное математическое ожидание существует и единственно с точностью до множеств меры нуль (теорема Радона-Никодима).
\end{theorem}

\begin{theorem}[Свойства условного ожидания]\;
\begin{enumerate}
    \item \textbf{Линейность:} $\E[aX + bY \mid \Gcal] = a\E[X \mid \Gcal] + b\E[Y \mid \Gcal]$
    \item \textbf{Монотонность:} Если $X \le Y$, то $\E[X \mid \Gcal] \le \E[Y \mid \Gcal]$
    \item \textbf{Башня:} Если $\Hcal \subseteq \Gcal$, то $\E[\E[X \mid \Gcal] \mid \Hcal] = \E[X \mid \Hcal]$
    \item \textbf{Взятие измеримого:} Если $Y$ измерима относительно $\Gcal$, то $\E[XY \mid \Gcal] = Y \E[X \mid \Gcal]$
\end{enumerate}
\end{theorem}

\subsection{Мартингалы}

\begin{definition}[Фильтрация]
Фильтрация --- это возрастающая последовательность $\si$-алгебр $\langle \Fcal_n \rangle_{n=0}^\infty$ такая, что $\Fcal_n \subseteq \Fcal_{n+1} \subseteq \Fcal$ для всех $n$.
\end{definition}

\begin{definition}[Мартингал]
Последовательность случайных величин $\langle X_n \rangle_{n=0}^\infty$ называется мартингалом относительно фильтрации $\{\Fcal_n\}$, если:
\begin{enumerate}
    \item $X_n$ измерима относительно $\Fcal_n$ для всех $n$
    \item $\E[|X_n|] < \infty$ для всех $n$
    \item \textbf{Условие мартингала:} $\E[X_{n+1} \mid \Fcal_n] = X_n$ почти наверное для всех $n$
\end{enumerate}
\end{definition}

\begin{definition}[Субмартингал и супермартингал]
\begin{enumerate}
    \item \textbf{Субмартингал:} $\E[X_{n+1} \mid \Fcal_n] \geq X_n$ (тенденция к росту)
    \item \textbf{Супермартингал:} $\E[X_{n+1} \mid \Fcal_n] \le X_n$ (тенденция к убыванию)
\end{enumerate}
\end{definition}

\begin{theorem}[Справедливая игра]
Мартингал моделирует справедливую игру: в каждый момент времени ожидаемый выигрыш на следующем шаге равен текущему капиталу.
\end{theorem}

\begin{theorem}[Теорема о сходимости мартингала]
Если $\langle X_n \rangle$ --- мартингал, ограниченный в $L^1$ (т.е. $\sup_n \E[|X_n|] < \infty$), то существует случайная величина $X_\infty$ такая, что $X_n \to X_\infty$ почти наверное и в $L^1$.
\end{theorem}

\begin{theorem}[Остановка мартингала (Optional Stopping)]
Если $\langle X_n \rangle$ --- мартингал, $\tau$ --- момент остановки (случайное время, зависящее только от прошлого), и выполнены условия интегрируемости, то $\E[X_\tau] = \E[X_0]$.
\end{theorem}

\begin{exercise}[Модель справедливого рынка]
Смоделируйте рынок, где цена актива $S_n$ в момент $n$ является мартингалом относительно фильтрации, порождённой историей цен.
\begin{enumerate}
    \item Покажите, что если рынок «справедлив», то $\E[S_{n+1} \mid S_0, \ldots, S_n] = S_n$.
    \item Интерпретируйте: что означает для философской модели «справедливого» процесса развития идеи?
    \item Постройте пример субмартингала (тенденция к росту) и супермартингала (тенденция к убыванию).
\end{enumerate}
\end{exercise}

\section{Мартингалы на деревьях и граница Мартина}
\textbf{Цель:} Применить теорию мартингалов к деревьям и изучить границу Мартина как пространство предельных точек.

\subsection{Мартингалы на деревьях}

\begin{definition}[Поток на дереве]
Пусть $\Tcal = \langle V, E, r \rangle$ --- корневое дерево. Поток --- это функция $q: V \to [0, \infty)$ такая, что:
$$q(x) = \sum_{y \in \Succ(x)} q(y) \quad \text{для всех } x \in V$$
с условием $q(r) = 1$ (нормировка).
\end{definition}

\begin{definition}[Вероятностная мера на пространстве путей]
Поток $q$ индуцирует вероятностную меру $\P$ на пространстве $\Om$ бесконечных путей:
$$\P(C_x) = q(x)$$
для цилиндра $C_x$, где $x \in V$.
\end{definition}

\begin{definition}[Мартингал на дереве]
Последовательность случайных величин $\langle M_n \rangle_{n=0}^\infty$ на $\langle \Om, \P \rangle$ называется мартингалом на дереве, если:
\begin{enumerate}
    \item $M_n$ измерима относительно $\si$-алгебры, порождённой цилиндрами вершин уровня $n$
    \item $\E[M_{n+1} \mid \Fcal_n] = M_n$, где $\Fcal_n$ --- $\si$-алгебра, порождённая путями до уровня $n$
\end{enumerate}
\end{definition}

\begin{theorem}[Поток как мартингал]
Если $q$ --- поток на дереве, то последовательность $M_n(\om) = q(x_n)$, где $x_n$ --- вершина уровня $n$ на пути $\om$, является мартингалом.
\end{theorem}

\begin{theorem}[Сходимость мартингала на дереве]
Если $\langle M_n \rangle$ --- мартингал на дереве, ограниченный в $L^1$, то $M_n(\om) \to M_\infty(\om)$ почти наверное для некоторой функции $M_\infty: \Om \to \R$.
\end{theorem}

\subsection{Граница Мартина}

\begin{definition}[Граница Мартина]
Граница Мартина $\partial \Tcal$ дерева $\Tcal$ --- это множество всех бесконечных путей $\Om$, снабжённое топологией, порождённой цилиндрами.

\textbf{Альтернативное определение:} Граница Мартина --- это компактификация дерева, получаемая добавлением «бесконечно удалённых» точек (концов путей).
\end{definition}

\begin{definition}[Гармоническая функция]
Функция $h: V \to \R$ называется гармонической, если для любой внутренней вершины $x$:
$$h(x) = \frac{1}{|\Succ(x)|} \sum_{y \in \Succ(x)} h(y)$$
(среднее значение на потомках равно значению в вершине).
\end{definition}

\begin{theorem}[Представление гармонических функций]
Любая ограниченная гармоническая функция $h$ на дереве имеет единственное представление:
$$h(x) = \int_{\partial \Tcal} K(x, \xi) \, d\mu(\xi)$$
где $K(x, \xi)$ --- ядро Мартина, а $\mu$ --- вероятностная мера на границе.
\end{theorem}

\begin{definition}[Ядро Мартина]
Ядро Мартина $K(x, \xi)$ для вершины $x$ и граничной точки $\xi \in \partial \Tcal$ определяется как предел отношения вероятностей достижения $\xi$ из различных вершин.
\end{definition}

\begin{theorem}[Сходимость к границе]
Для почти всех путей $\om \in \Om$ мартингал $M_n(\om)$ сходится к пределу $M_\infty(\om)$, который можно интерпретировать как «значение на границе Мартина».
\end{theorem}

\begin{exercise}[Модель развития с границей]
Постройте модель развития философской системы как дерево, где:
\begin{enumerate}
    \item Вершины --- состояния системы
    \item Поток $q$ --- «вероятность развития» по каждому пути
    \item Покажите, что граница Мартина соответствует «предельным состояниям» системы
    \item Интерпретируйте сходимость мартингала: что означает «предельное значение» на границе?
\end{enumerate}
\end{exercise}

\section{Леса Гальтона-Ватсона и ветвящиеся процессы}
\textbf{Цель:} Изучить ветвящиеся процессы как модель размножения и применить к философским моделям распространения идей.

\subsection{Процесс Гальтона-Ватсона}

\begin{definition}[Процесс Гальтона-Ватсона]
Процесс Гальтона-Ватсона $\langle Z_n \rangle_{n=0}^\infty$ --- это марковская цепь, моделирующая размер популяции в поколении $n$:
\begin{enumerate}
    \item $Z_0 = 1$ (одна особь в начальном поколении)
    \item Каждая особь в поколении $n$ независимо порождает случайное число потомков с распределением $\{p_k\}_{k=0}^\infty$
    \item $Z_{n+1} = \sum_{i=1}^{Z_n} \xi_i$, где $\xi_i$ --- число потомков $i$-й особи
\end{enumerate}
\end{definition}

\begin{definition}[Производящая функция]
Производящая функция распределения числа потомков:
$$f(s) = \sum_{k=0}^\infty p_k s^k, \quad s \in [0, 1]$$
\end{definition}

\begin{theorem}[Математическое ожидание]
Если $m = \E[\xi] = f'(1) < \infty$, то $\E[Z_n] = m^n$.
\end{theorem}

\begin{theorem}[Вероятность вымирания]
Вероятность вымирания $q = \P(Z_n = 0 \text{ для некоторого } n)$ является наименьшим неотрицательным корнем уравнения $f(s) = s$.

\textbf{Критический случай:}
\begin{enumerate}
    \item Если $m < 1$ (докритический): $q = 1$ (вымирание почти наверное)
    \item Если $m = 1$ (критический): $q = 1$ (вымирание почти наверное)
    \item Если $m > 1$ (надкритический): $q < 1$ (положительная вероятность выживания)
\end{enumerate}
\end{theorem}

\begin{remark}[Физическая интерпретация: цепная ядерная реакция]
Процесс Гальтона-Ватсона моделирует цепную ядерную реакцию:
\begin{itemize}
    \item Каждый нейтрон может породить случайное число новых нейтронов
    \item Если коэффициент размножения $m > 1$ (надкритический случай): происходит взрыв (бесконечное дерево нейтронов)
    \item Если $m < 1$ (докритический случай): реакция затухает (дерево вымирает)
    \item Если $m = 1$ (критический случай): реакция самоподдерживается, но не растет экспоненциально
\end{itemize}
Это показывает, как математическая модель ветвящегося процесса описывает физический феномен ядерного взрыва.
\end{remark}

\subsection{Лес Гальтона-Ватсона}

\begin{definition}[Лес Гальтона-Ватсона]
Лес Гальтона-Ватсона --- это случайный лес (объединение деревьев), где каждое дерево --- это генеалогическое дерево процесса Гальтона-Ватсона, начинающегося с одной особи.

Формально: счётное объединение независимых деревьев, каждое из которых построено по процессу Гальтона-Ватсона.
\end{definition}

\begin{theorem}[Структура леса]
В надкритическом случае ($m > 1$) лес почти наверное содержит бесконечное дерево (дерево с ненулевой вероятностью выживания).
\end{theorem}

\begin{theorem}[Применение леммы Кёнига]
Если процесс Гальтона-Ватсона не вымирает, то генеалогическое дерево бесконечно. По лемме Кёнига, такое дерево содержит бесконечный путь (линию потомков, которая никогда не прерывается).
\end{theorem}

\begin{exercise}[Модель распространения идеи]
Смоделируйте распространение философской идеи как процесс Гальтона-Ватсона:
\begin{enumerate}
    \item Каждый «носитель» идеи может «заразить» случайное число новых носителей
    \item Найдите условие, при котором идея распространяется (надкритический случай)
    \item Интерпретируйте вымирание: что означает для идеи «умереть»?
    \item Постройте лес идей: несколько независимых «зародышей» идеи развиваются параллельно
\end{enumerate}
\end{exercise}

\begin{remark}
Леса Гальтона-Ватсона естественно связаны с мартингалами на деревьях: поток на генеалогическом дереве соответствует вероятностям выживания различных линий потомков. Граница Мартина такого дерева описывает «предельные» линии развития популяции.
\end{remark}

\end{document}